
% =======================================================================
% APPENDIX C: Detailergebnisse der Monte-Carlo Simulation (Orthogonalität)
% =======================================================================
\section{Stabilitäts- und Konvergenzanalysen der Design-Modifikationen}
\label{app:c}

In diesem Anhang sind ergänzende statistische Visualisierungen zur Bewertung der geometrischen Design-Verzerrungen (vergleiche Abschnitt 5.3) aufgeführt. Sämtliche Analysen wurden für ein $\sym{snr}=1{,}33$ und über $n_{MC}=10.000$ Monte-Carlo-Iterationen durchgeführt.

\begin{figure}[htbp]
    \centering
    \setlength\figurewidth{0.5\textwidth}
    \setlength\figureheight{16cm}
    % This file was created by matlab2tikz.
%
\definecolor{mycolor1}{rgb}{0.12941,0.12941,0.12941}%
%
\begin{tikzpicture}

\begin{axis}[%
width=0.391\figurewidth,
height=0.265\figureheight,
at={(0\figurewidth,0.735\figureheight)},
scale only axis,
xmin=1.00,
xmax=10000.00,
xtick={2000.00,4000.00,6000.00,8000.00,10000.00},
xticklabels={\empty},
ymin=0.00,
ymax=105.00,
ylabel style={font=\color{mycolor1}},
ylabel={\sym{power} [\%]},
axis background/.style={fill=white},
title style={font=\bfseries\color{mycolor1}},
title={\textbf{$\sym{beta}_0$}},
xmajorgrids,
ymajorgrids,
grid style={dotted, opacity=0.25},
legend style={at={(0.03,0.03)}, anchor=south west, legend cell align=left, align=left}
]
\addplot [color=white!10!black, line width=1.2pt]
  table[row sep=crcr]{%
1.00	100.00\\
10000.00	100.00\\
};
\addlegendentry{Basis \ac{CCD}}

\addplot [color=black!25!darkgray, dashed, line width=1.2pt]
  table[row sep=crcr]{%
1.00	100.00\\
10000.00	100.00\\
};
\addlegendentry{Fall I ($-20\,\%$)}

\addplot [color=black!45!gray, dotted, line width=1.2pt]
  table[row sep=crcr]{%
1.00	100.00\\
10000.00	99.99\\
};
\addlegendentry{Fall I ($+20\,\%$)}

\addplot [color=white!15!darkgray, dashdotted, line width=1.2pt]
  table[row sep=crcr]{%
1.00	100.00\\
10000.00	100.00\\
};
\addlegendentry{Fall I (\ac{FC})}

\addplot [color=white!45!black, line width=1.2pt]
  table[row sep=crcr]{%
1.00	100.00\\
10000.00	99.99\\
};
\addlegendentry{Fall II (leicht)}

\addplot [color=gray!85!lightgray, dashed, line width=1.2pt]
  table[row sep=crcr]{%
1.00	100.00\\
10000.00	99.99\\
};
\addlegendentry{Fall II (stark)}

\addplot [color=white!50!darkgray, dotted, line width=1.2pt]
  table[row sep=crcr]{%
1.00	100.00\\
10000.00	99.99\\
};
\addlegendentry{Fall III}

\addplot [color=black!5!lightgray, dashdotted, line width=1.2pt]
  table[row sep=crcr]{%
1.00	100.00\\
10000.00	99.99\\
};
\addlegendentry{Fall IV: \ac{SP}}

\addplot [color=white!80!black, line width=1.2pt]
  table[row sep=crcr]{%
1.00	100.00\\
10000.00	99.97\\
};
\addlegendentry{Fall IV: \ac{ZP}}

\end{axis}

\begin{axis}[%
width=0.391\figurewidth,
height=0.265\figureheight,
at={(0.54\figurewidth,0.735\figureheight)},
scale only axis,
xmin=1.00,
xmax=10000.00,
xtick={2000.00,4000.00,6000.00,8000.00,10000.00},
xticklabels={\empty},
ymin=0.00,
ymax=105.00,
ytick={0.00,20.00,40.00,60.00,80.00,100.00},
yticklabels={\empty},
axis background/.style={fill=white},
title style={font=\bfseries\color{mycolor1}},
title={\textbf{$\sym{beta}_1$}},
xmajorgrids,
ymajorgrids,
grid style={dotted, opacity=0.25}
]
\addplot [color=white!10!black, line width=1.2pt, forget plot]
  table[row sep=crcr]{%
1.00	100.00\\
19.00	100.00\\
21.00	95.24\\
23.00	95.65\\
26.00	96.15\\
30.00	96.67\\
32.00	96.88\\
34.00	94.12\\
36.00	91.67\\
41.00	92.68\\
46.00	93.48\\
52.00	94.23\\
56.00	92.86\\
63.00	93.65\\
67.00	92.54\\
76.00	93.42\\
86.00	94.19\\
97.00	94.85\\
110.00	95.45\\
117.00	94.87\\
124.00	95.16\\
132.00	94.70\\
149.00	95.30\\
159.00	95.60\\
169.00	95.27\\
180.00	95.00\\
191.00	95.29\\
217.00	94.93\\
230.00	94.35\\
261.00	94.25\\
295.00	94.92\\
314.00	94.59\\
334.00	94.61\\
355.00	94.37\\
378.00	93.92\\
402.00	94.28\\
427.00	93.91\\
455.00	93.63\\
484.00	93.60\\
515.00	93.40\\
547.00	93.60\\
659.00	93.78\\
701.00	94.01\\
746.00	93.57\\
793.00	93.57\\
844.00	93.72\\
897.00	93.42\\
955.00	93.61\\
1016.00	93.90\\
1149.00	93.91\\
1222.00	93.62\\
1383.00	93.85\\
1472.00	93.68\\
1565.00	93.87\\
1665.00	93.75\\
1771.00	93.90\\
2268.00	93.47\\
2413.00	93.54\\
2730.00	93.77\\
2905.00	93.67\\
3090.00	93.43\\
3287.00	93.43\\
3496.00	93.62\\
3719.00	93.68\\
4209.00	93.47\\
5066.00	93.51\\
6099.00	93.54\\
6901.00	93.44\\
7809.00	93.47\\
8307.00	93.46\\
8837.00	93.57\\
9401.00	93.50\\
10000.00	93.58\\
};
\addplot [color=black!25!darkgray, dashed, line width=1.2pt, forget plot]
  table[row sep=crcr]{%
1.00	100.00\\
11.00	100.00\\
12.00	91.67\\
14.00	92.86\\
16.00	93.75\\
18.00	94.44\\
19.00	94.74\\
21.00	90.48\\
23.00	91.30\\
25.00	92.00\\
26.00	88.46\\
28.00	85.71\\
30.00	83.33\\
32.00	84.38\\
34.00	85.29\\
36.00	83.33\\
38.00	84.21\\
41.00	82.93\\
43.00	83.72\\
46.00	84.78\\
49.00	85.71\\
52.00	86.54\\
56.00	85.71\\
59.00	84.75\\
63.00	85.71\\
67.00	86.57\\
71.00	87.32\\
76.00	86.84\\
81.00	87.65\\
91.00	89.01\\
103.00	90.29\\
117.00	91.45\\
132.00	92.42\\
140.00	91.43\\
159.00	92.45\\
169.00	92.31\\
180.00	92.78\\
191.00	92.67\\
217.00	92.63\\
245.00	93.47\\
261.00	93.10\\
277.00	93.50\\
295.00	93.22\\
314.00	93.63\\
334.00	93.41\\
355.00	92.96\\
378.00	93.12\\
402.00	93.53\\
455.00	93.41\\
484.00	93.39\\
515.00	93.20\\
547.00	93.24\\
582.00	93.47\\
619.00	93.21\\
659.00	93.47\\
701.00	93.44\\
746.00	93.70\\
844.00	93.72\\
955.00	94.03\\
1016.00	93.70\\
1080.00	93.70\\
1149.00	93.39\\
1565.00	93.35\\
1665.00	93.45\\
1771.00	93.34\\
2005.00	93.82\\
2132.00	93.71\\
2413.00	93.66\\
2567.00	93.42\\
2905.00	93.49\\
3090.00	93.40\\
3496.00	93.51\\
3719.00	93.49\\
3957.00	93.33\\
4763.00	93.51\\
5066.00	93.33\\
6099.00	93.39\\
6901.00	93.31\\
7341.00	93.27\\
7809.00	93.37\\
8307.00	93.34\\
8837.00	93.19\\
9401.00	93.15\\
10000.00	93.23\\
};
\addplot [color=black!45!gray, dotted, line width=1.2pt, forget plot]
  table[row sep=crcr]{%
1.00	100.00\\
15.00	100.00\\
16.00	93.75\\
18.00	94.44\\
21.00	95.24\\
23.00	95.65\\
26.00	96.15\\
30.00	96.67\\
36.00	97.22\\
43.00	97.67\\
52.00	98.08\\
63.00	98.41\\
71.00	98.59\\
76.00	97.37\\
81.00	97.53\\
86.00	96.51\\
91.00	96.70\\
97.00	95.88\\
103.00	96.12\\
110.00	95.45\\
124.00	95.97\\
140.00	95.00\\
159.00	95.60\\
169.00	95.27\\
180.00	95.00\\
191.00	93.72\\
217.00	94.47\\
245.00	95.10\\
261.00	95.02\\
295.00	95.59\\
334.00	96.11\\
355.00	96.06\\
378.00	95.77\\
402.00	95.77\\
427.00	95.55\\
455.00	95.38\\
484.00	94.83\\
547.00	94.70\\
582.00	94.67\\
619.00	94.51\\
659.00	94.54\\
746.00	94.10\\
793.00	94.45\\
844.00	94.19\\
897.00	94.31\\
955.00	93.93\\
1016.00	94.00\\
1080.00	94.17\\
1222.00	94.35\\
1300.00	94.54\\
1383.00	94.50\\
1472.00	94.29\\
1565.00	94.19\\
1665.00	94.29\\
1884.00	94.06\\
2005.00	94.11\\
2132.00	93.90\\
2268.00	94.09\\
2413.00	94.20\\
2567.00	94.20\\
2905.00	93.98\\
3496.00	94.05\\
4209.00	93.80\\
5389.00	93.91\\
5733.00	93.86\\
6488.00	94.02\\
8837.00	93.96\\
9401.00	93.90\\
10000.00	93.97\\
};
\addplot [color=white!15!darkgray, dashdotted, line width=1.2pt, forget plot]
  table[row sep=crcr]{%
1.00	100.00\\
5.00	100.00\\
6.00	83.33\\
7.00	85.71\\
8.00	87.50\\
9.00	88.89\\
10.00	90.00\\
11.00	90.91\\
12.00	91.67\\
14.00	92.86\\
16.00	93.75\\
18.00	94.44\\
21.00	95.24\\
23.00	95.65\\
26.00	96.15\\
28.00	96.43\\
30.00	93.33\\
32.00	93.75\\
34.00	91.18\\
38.00	92.11\\
43.00	93.02\\
49.00	93.88\\
56.00	94.64\\
63.00	95.24\\
71.00	95.77\\
81.00	96.30\\
86.00	96.51\\
91.00	95.60\\
97.00	95.88\\
103.00	95.15\\
110.00	94.55\\
117.00	94.87\\
124.00	94.35\\
132.00	94.70\\
140.00	94.29\\
149.00	94.63\\
159.00	94.34\\
169.00	94.67\\
180.00	93.89\\
191.00	94.24\\
204.00	94.12\\
217.00	93.55\\
230.00	93.48\\
261.00	94.25\\
277.00	94.58\\
334.00	94.61\\
378.00	95.24\\
427.00	95.32\\
455.00	95.16\\
484.00	95.25\\
515.00	95.15\\
547.00	95.25\\
582.00	94.85\\
619.00	94.67\\
659.00	94.99\\
701.00	94.86\\
746.00	94.91\\
793.00	94.58\\
844.00	94.67\\
897.00	94.65\\
955.00	94.76\\
1016.00	94.39\\
1080.00	94.26\\
1149.00	94.34\\
1222.00	94.11\\
1300.00	94.31\\
1383.00	94.36\\
1472.00	94.29\\
1565.00	94.44\\
1665.00	94.41\\
1771.00	94.18\\
2268.00	93.87\\
2413.00	93.87\\
2730.00	94.21\\
3090.00	93.85\\
3287.00	93.76\\
3496.00	93.88\\
3957.00	93.48\\
4477.00	93.50\\
4763.00	93.43\\
5066.00	93.09\\
5389.00	93.10\\
6099.00	92.75\\
6488.00	92.66\\
6901.00	92.83\\
8307.00	92.58\\
9401.00	92.54\\
10000.00	92.40\\
};
\addplot [color=white!45!black, line width=1.2pt, forget plot]
  table[row sep=crcr]{%
1.00	0.00\\
2.00	50.00\\
3.00	66.67\\
4.00	75.00\\
5.00	80.00\\
6.00	83.33\\
7.00	71.43\\
8.00	75.00\\
9.00	77.78\\
10.00	80.00\\
11.00	81.82\\
12.00	83.33\\
13.00	84.62\\
14.00	85.71\\
15.00	86.67\\
16.00	81.25\\
17.00	82.35\\
19.00	84.21\\
21.00	85.71\\
23.00	86.96\\
26.00	88.46\\
28.00	89.29\\
30.00	86.67\\
32.00	87.50\\
34.00	85.29\\
38.00	86.84\\
43.00	88.37\\
46.00	86.96\\
49.00	87.76\\
52.00	86.54\\
56.00	87.50\\
59.00	88.14\\
63.00	88.89\\
67.00	89.55\\
71.00	88.73\\
76.00	88.16\\
81.00	88.89\\
86.00	89.53\\
91.00	89.01\\
97.00	88.66\\
103.00	87.38\\
110.00	88.18\\
124.00	89.52\\
132.00	90.15\\
140.00	90.00\\
149.00	89.93\\
169.00	91.12\\
191.00	92.15\\
204.00	91.18\\
217.00	91.71\\
230.00	91.30\\
245.00	91.43\\
277.00	92.42\\
295.00	92.88\\
314.00	92.68\\
334.00	92.81\\
355.00	92.68\\
378.00	92.86\\
427.00	92.74\\
455.00	93.19\\
484.00	92.77\\
547.00	92.87\\
582.00	92.61\\
619.00	92.57\\
701.00	92.87\\
844.00	92.89\\
897.00	93.20\\
955.00	93.30\\
1016.00	93.31\\
1080.00	93.06\\
1149.00	93.21\\
1222.00	93.29\\
1300.00	93.23\\
1383.00	93.35\\
1472.00	93.21\\
1565.00	93.48\\
1665.00	93.27\\
1771.00	93.28\\
1884.00	93.58\\
2268.00	93.87\\
2567.00	93.88\\
2905.00	94.04\\
3287.00	93.89\\
3496.00	93.62\\
3719.00	93.55\\
3957.00	93.63\\
4209.00	93.51\\
5066.00	93.60\\
5389.00	93.45\\
6099.00	93.47\\
7809.00	93.64\\
8307.00	93.72\\
9401.00	93.67\\
10000.00	93.67\\
};
\addplot [color=gray!85!lightgray, dashed, line width=1.2pt, forget plot]
  table[row sep=crcr]{%
1.00	100.00\\
13.00	100.00\\
14.00	92.86\\
16.00	93.75\\
18.00	94.44\\
21.00	95.24\\
23.00	95.65\\
26.00	96.15\\
28.00	96.43\\
30.00	93.33\\
34.00	94.12\\
36.00	91.67\\
41.00	92.68\\
46.00	93.48\\
52.00	94.23\\
56.00	94.64\\
59.00	93.22\\
63.00	92.06\\
67.00	91.04\\
76.00	92.11\\
81.00	91.36\\
91.00	92.31\\
103.00	93.20\\
117.00	94.02\\
124.00	94.35\\
132.00	93.18\\
140.00	92.14\\
149.00	91.95\\
159.00	91.19\\
180.00	92.22\\
191.00	92.67\\
204.00	92.16\\
230.00	93.04\\
245.00	93.47\\
261.00	93.49\\
277.00	93.86\\
295.00	93.56\\
334.00	94.31\\
355.00	93.52\\
402.00	93.78\\
427.00	93.44\\
455.00	93.19\\
484.00	92.77\\
515.00	92.43\\
547.00	92.50\\
582.00	92.27\\
659.00	92.56\\
701.00	92.87\\
746.00	92.90\\
793.00	93.06\\
844.00	93.01\\
955.00	93.19\\
1016.00	93.11\\
1080.00	92.78\\
1149.00	92.95\\
1222.00	93.04\\
1300.00	93.00\\
1383.00	93.28\\
1472.00	93.21\\
1565.00	93.42\\
1665.00	93.45\\
1771.00	93.22\\
1884.00	93.26\\
2005.00	93.17\\
2132.00	93.34\\
2268.00	93.39\\
2567.00	93.07\\
2730.00	93.26\\
2905.00	93.18\\
3090.00	93.33\\
3287.00	93.37\\
3496.00	93.22\\
3719.00	93.17\\
3957.00	93.28\\
4209.00	93.18\\
4477.00	92.99\\
5389.00	93.08\\
5733.00	93.09\\
7341.00	92.81\\
7809.00	92.87\\
10000.00	92.77\\
};
\addplot [color=white!50!darkgray, dotted, line width=1.2pt, forget plot]
  table[row sep=crcr]{%
1.00	100.00\\
18.00	100.00\\
19.00	94.74\\
22.00	95.45\\
25.00	96.00\\
26.00	96.15\\
28.00	92.86\\
30.00	90.00\\
34.00	91.18\\
38.00	92.11\\
41.00	90.24\\
46.00	91.30\\
52.00	92.31\\
59.00	93.22\\
67.00	94.03\\
71.00	92.96\\
81.00	93.83\\
91.00	94.51\\
103.00	93.20\\
110.00	92.73\\
124.00	93.55\\
132.00	93.18\\
149.00	93.96\\
169.00	93.49\\
180.00	93.33\\
191.00	92.67\\
217.00	92.63\\
245.00	93.47\\
261.00	93.10\\
277.00	93.50\\
295.00	93.56\\
334.00	94.31\\
355.00	94.37\\
402.00	95.02\\
427.00	94.61\\
455.00	94.95\\
484.00	94.63\\
515.00	94.56\\
547.00	94.70\\
582.00	94.50\\
619.00	94.67\\
701.00	94.72\\
793.00	94.33\\
844.00	94.31\\
897.00	93.53\\
955.00	93.40\\
1149.00	93.39\\
1222.00	93.62\\
1300.00	93.38\\
1472.00	93.75\\
1565.00	93.74\\
1665.00	93.57\\
1884.00	93.63\\
2132.00	93.86\\
2268.00	93.56\\
2413.00	93.49\\
2567.00	93.53\\
2905.00	93.73\\
3090.00	93.56\\
3287.00	93.55\\
3496.00	93.74\\
3719.00	93.65\\
3957.00	93.73\\
4209.00	93.61\\
4763.00	93.66\\
5733.00	93.42\\
6099.00	93.29\\
8837.00	93.43\\
9401.00	93.54\\
10000.00	93.54\\
};
\addplot [color=black!5!lightgray, dashdotted, line width=1.2pt, forget plot]
  table[row sep=crcr]{%
1.00	100.00\\
2.00	50.00\\
3.00	66.67\\
4.00	75.00\\
5.00	80.00\\
6.00	83.33\\
7.00	85.71\\
8.00	87.50\\
9.00	88.89\\
10.00	90.00\\
11.00	81.82\\
12.00	83.33\\
13.00	76.92\\
14.00	78.57\\
15.00	73.33\\
16.00	75.00\\
17.00	76.47\\
18.00	77.78\\
19.00	78.95\\
21.00	80.95\\
23.00	82.61\\
26.00	84.62\\
28.00	85.71\\
30.00	83.33\\
32.00	84.38\\
34.00	82.35\\
36.00	83.33\\
38.00	84.21\\
41.00	85.37\\
43.00	86.05\\
46.00	86.96\\
49.00	85.71\\
52.00	84.62\\
56.00	85.71\\
59.00	84.75\\
63.00	85.71\\
71.00	84.51\\
76.00	85.53\\
81.00	83.95\\
86.00	84.88\\
91.00	84.62\\
97.00	85.57\\
103.00	86.41\\
117.00	86.32\\
124.00	87.10\\
149.00	87.25\\
159.00	87.42\\
169.00	86.39\\
191.00	85.86\\
204.00	86.76\\
217.00	86.64\\
230.00	85.65\\
245.00	85.31\\
261.00	85.44\\
295.00	86.44\\
314.00	86.62\\
334.00	87.13\\
355.00	87.04\\
378.00	87.83\\
402.00	87.31\\
455.00	87.47\\
484.00	87.60\\
515.00	87.96\\
547.00	87.93\\
582.00	87.63\\
701.00	88.45\\
746.00	88.34\\
793.00	87.77\\
844.00	87.44\\
897.00	87.96\\
1080.00	88.33\\
1149.00	88.60\\
1222.00	88.79\\
1300.00	89.08\\
1383.00	89.23\\
1472.00	89.20\\
1665.00	89.37\\
1884.00	89.28\\
2005.00	89.13\\
2413.00	89.47\\
2567.00	89.37\\
2730.00	89.49\\
3719.00	89.35\\
3957.00	89.49\\
4763.00	89.65\\
6488.00	89.06\\
6901.00	89.25\\
7341.00	89.16\\
7809.00	88.96\\
8307.00	89.07\\
9401.00	89.14\\
10000.00	89.25\\
};
\addplot [color=white!80!black, line width=1.2pt, forget plot]
  table[row sep=crcr]{%
1.00	100.00\\
7.00	100.00\\
8.00	87.50\\
9.00	88.89\\
10.00	90.00\\
11.00	81.82\\
12.00	83.33\\
13.00	84.62\\
14.00	78.57\\
15.00	80.00\\
16.00	81.25\\
17.00	82.35\\
19.00	84.21\\
21.00	85.71\\
23.00	86.96\\
26.00	88.46\\
28.00	89.29\\
30.00	90.00\\
34.00	91.18\\
38.00	92.11\\
41.00	92.68\\
43.00	90.70\\
46.00	89.13\\
49.00	87.76\\
52.00	88.46\\
59.00	89.83\\
63.00	88.89\\
67.00	89.55\\
71.00	88.73\\
76.00	89.47\\
86.00	90.70\\
97.00	91.75\\
103.00	92.23\\
117.00	91.45\\
124.00	91.13\\
132.00	91.67\\
140.00	91.43\\
149.00	91.28\\
169.00	92.31\\
191.00	93.19\\
204.00	93.14\\
217.00	93.55\\
230.00	93.48\\
245.00	93.88\\
261.00	93.49\\
277.00	92.78\\
334.00	93.11\\
355.00	91.83\\
378.00	92.06\\
402.00	91.79\\
455.00	91.87\\
515.00	92.82\\
547.00	92.87\\
582.00	93.30\\
619.00	93.38\\
701.00	93.30\\
746.00	92.90\\
793.00	93.32\\
844.00	93.25\\
897.00	93.42\\
955.00	93.72\\
1016.00	93.31\\
1149.00	93.39\\
1222.00	93.78\\
1383.00	93.93\\
1472.00	93.89\\
1565.00	93.61\\
1665.00	93.75\\
1884.00	93.63\\
2005.00	93.52\\
2132.00	93.76\\
2567.00	93.77\\
3090.00	93.85\\
3287.00	93.82\\
3719.00	94.06\\
4209.00	93.99\\
5066.00	93.66\\
5733.00	93.63\\
6099.00	93.36\\
7341.00	93.22\\
8307.00	93.27\\
10000.00	93.25\\
};
\end{axis}

\begin{axis}[%
width=0.391\figurewidth,
height=0.265\figureheight,
at={(0\figurewidth,0.368\figureheight)},
scale only axis,
xmin=1.00,
xmax=10000.00,
xtick={2000.00,4000.00,6000.00,8000.00,10000.00},
xticklabels={\empty},
ymin=0.00,
ymax=105.00,
ylabel style={font=\color{mycolor1}},
ylabel={\sym{power} [\%]},
axis background/.style={fill=white},
title style={font=\bfseries\color{mycolor1}},
title={\textbf{$\sym{beta}_2$}},
xmajorgrids,
ymajorgrids,
grid style={dotted, opacity=0.25}
]
\addplot [color=white!10!black, line width=1.2pt, forget plot]
  table[row sep=crcr]{%
1.00	100.00\\
38.00	100.00\\
41.00	97.56\\
49.00	97.96\\
56.00	98.21\\
59.00	94.92\\
67.00	95.52\\
71.00	95.77\\
76.00	94.74\\
81.00	93.83\\
91.00	94.51\\
97.00	94.85\\
103.00	94.17\\
110.00	93.64\\
124.00	94.35\\
140.00	95.00\\
149.00	95.30\\
159.00	94.97\\
169.00	94.08\\
180.00	93.89\\
191.00	93.19\\
204.00	93.14\\
217.00	92.17\\
245.00	92.24\\
261.00	92.72\\
277.00	92.78\\
295.00	93.22\\
334.00	92.81\\
355.00	92.68\\
378.00	93.12\\
402.00	93.28\\
427.00	93.68\\
455.00	93.85\\
484.00	93.80\\
515.00	94.17\\
582.00	93.81\\
619.00	94.02\\
659.00	93.93\\
701.00	93.72\\
746.00	93.30\\
844.00	92.77\\
897.00	92.98\\
955.00	92.98\\
1016.00	93.31\\
1080.00	93.43\\
1222.00	93.45\\
1300.00	93.69\\
1383.00	93.78\\
1472.00	94.02\\
1565.00	93.74\\
1665.00	93.75\\
1884.00	93.58\\
2005.00	93.67\\
2132.00	93.67\\
2268.00	93.78\\
2413.00	93.70\\
2567.00	93.84\\
2730.00	93.59\\
2905.00	93.49\\
3090.00	93.59\\
3287.00	93.61\\
3496.00	93.74\\
3719.00	93.60\\
3957.00	93.66\\
4209.00	93.87\\
4763.00	93.93\\
5066.00	93.82\\
6099.00	93.80\\
6488.00	93.76\\
6901.00	93.81\\
7809.00	93.71\\
8837.00	93.80\\
9401.00	93.70\\
10000.00	93.69\\
};
\addplot [color=black!25!darkgray, dashed, line width=1.2pt, forget plot]
  table[row sep=crcr]{%
1.00	100.00\\
23.00	100.00\\
25.00	96.00\\
28.00	96.43\\
30.00	96.67\\
32.00	93.75\\
36.00	94.44\\
41.00	95.12\\
46.00	95.65\\
52.00	96.15\\
56.00	94.64\\
63.00	95.24\\
71.00	95.77\\
81.00	96.30\\
97.00	96.91\\
103.00	96.12\\
117.00	96.58\\
124.00	95.97\\
140.00	95.00\\
159.00	95.60\\
169.00	95.27\\
191.00	95.81\\
204.00	95.59\\
217.00	94.93\\
230.00	94.78\\
245.00	93.47\\
261.00	93.49\\
277.00	93.14\\
314.00	92.68\\
378.00	93.12\\
402.00	93.53\\
427.00	93.21\\
455.00	93.63\\
484.00	93.80\\
515.00	93.40\\
547.00	93.24\\
619.00	93.70\\
659.00	93.63\\
701.00	93.30\\
746.00	93.16\\
844.00	93.48\\
897.00	93.87\\
955.00	93.51\\
1016.00	93.50\\
1080.00	93.33\\
1149.00	93.30\\
1222.00	93.37\\
1300.00	93.54\\
1472.00	93.48\\
1565.00	93.55\\
1665.00	93.75\\
1771.00	93.45\\
1884.00	93.68\\
2132.00	93.90\\
2268.00	94.05\\
2567.00	94.04\\
2905.00	94.04\\
3287.00	93.79\\
3496.00	93.74\\
3719.00	93.76\\
3957.00	93.66\\
4477.00	93.70\\
4763.00	93.70\\
5066.00	93.82\\
5389.00	93.84\\
5733.00	93.74\\
6488.00	93.76\\
7809.00	93.87\\
10000.00	93.78\\
};
\addplot [color=black!45!gray, dotted, line width=1.2pt, forget plot]
  table[row sep=crcr]{%
1.00	100.00\\
23.00	100.00\\
25.00	96.00\\
26.00	96.15\\
28.00	92.86\\
32.00	93.75\\
36.00	94.44\\
41.00	95.12\\
46.00	95.65\\
49.00	95.92\\
52.00	94.23\\
59.00	94.92\\
63.00	93.65\\
71.00	94.37\\
81.00	95.06\\
86.00	95.35\\
91.00	93.41\\
103.00	94.17\\
110.00	94.55\\
117.00	94.02\\
132.00	94.70\\
140.00	95.00\\
149.00	94.63\\
169.00	95.27\\
191.00	94.76\\
230.00	94.35\\
245.00	94.69\\
261.00	94.64\\
295.00	95.25\\
334.00	94.61\\
355.00	94.65\\
378.00	94.44\\
402.00	94.03\\
427.00	94.15\\
455.00	94.07\\
484.00	93.80\\
515.00	93.98\\
547.00	93.97\\
619.00	93.70\\
659.00	93.78\\
701.00	93.72\\
746.00	93.57\\
793.00	93.57\\
897.00	93.42\\
955.00	93.40\\
1016.00	93.60\\
1080.00	93.70\\
1149.00	94.08\\
1222.00	94.11\\
1300.00	94.31\\
1383.00	94.14\\
1472.00	94.16\\
1665.00	93.99\\
1771.00	94.01\\
2005.00	94.31\\
2132.00	94.37\\
2268.00	94.22\\
2413.00	94.24\\
2567.00	94.00\\
2730.00	94.10\\
2905.00	94.08\\
3090.00	94.21\\
3287.00	94.07\\
3496.00	94.19\\
3957.00	94.21\\
4209.00	94.25\\
4763.00	94.06\\
5066.00	94.02\\
5389.00	94.10\\
6099.00	94.15\\
6901.00	94.22\\
7809.00	94.30\\
8307.00	94.09\\
10000.00	94.14\\
};
\addplot [color=white!15!darkgray, dashdotted, line width=1.2pt, forget plot]
  table[row sep=crcr]{%
1.00	100.00\\
15.00	100.00\\
16.00	93.75\\
18.00	94.44\\
21.00	95.24\\
23.00	95.65\\
25.00	96.00\\
26.00	92.31\\
30.00	93.33\\
34.00	94.12\\
38.00	94.74\\
43.00	95.35\\
49.00	95.92\\
56.00	96.43\\
59.00	94.92\\
63.00	95.24\\
67.00	94.03\\
76.00	94.74\\
86.00	95.35\\
97.00	95.88\\
110.00	96.36\\
124.00	96.77\\
132.00	96.21\\
149.00	96.64\\
159.00	96.86\\
169.00	96.45\\
180.00	96.67\\
191.00	96.34\\
204.00	96.57\\
217.00	96.31\\
245.00	96.73\\
277.00	95.67\\
295.00	95.93\\
314.00	95.22\\
334.00	95.21\\
355.00	94.93\\
378.00	94.71\\
402.00	95.02\\
427.00	95.08\\
455.00	95.38\\
484.00	95.45\\
515.00	95.34\\
547.00	95.43\\
582.00	95.19\\
619.00	94.83\\
659.00	94.99\\
701.00	95.01\\
746.00	95.17\\
793.00	94.83\\
844.00	94.67\\
1016.00	94.69\\
1080.00	94.72\\
1149.00	94.43\\
1222.00	94.52\\
1300.00	94.46\\
1383.00	94.29\\
1472.00	93.95\\
1665.00	93.87\\
1771.00	93.73\\
1884.00	93.74\\
2005.00	93.87\\
2132.00	93.76\\
2268.00	93.87\\
2413.00	93.83\\
2730.00	94.03\\
2905.00	94.15\\
3090.00	94.11\\
3287.00	93.95\\
3496.00	94.02\\
3719.00	94.00\\
3957.00	93.78\\
5066.00	93.78\\
5389.00	93.93\\
6488.00	93.88\\
6901.00	93.86\\
7809.00	93.89\\
8307.00	93.93\\
8837.00	93.82\\
10000.00	93.85\\
};
\addplot [color=white!45!black, line width=1.2pt, forget plot]
  table[row sep=crcr]{%
1.00	100.00\\
6.00	100.00\\
7.00	85.71\\
8.00	87.50\\
9.00	77.78\\
10.00	80.00\\
11.00	81.82\\
12.00	83.33\\
13.00	84.62\\
14.00	85.71\\
15.00	86.67\\
17.00	88.24\\
19.00	89.47\\
21.00	85.71\\
23.00	86.96\\
26.00	88.46\\
28.00	89.29\\
30.00	90.00\\
32.00	87.50\\
36.00	88.89\\
38.00	89.47\\
43.00	90.70\\
49.00	91.84\\
52.00	92.31\\
56.00	91.07\\
59.00	91.53\\
63.00	90.48\\
71.00	91.55\\
76.00	90.79\\
86.00	91.86\\
91.00	90.11\\
103.00	91.26\\
110.00	90.91\\
124.00	91.94\\
132.00	92.42\\
140.00	90.71\\
149.00	90.60\\
159.00	91.19\\
169.00	90.53\\
180.00	90.56\\
191.00	91.10\\
204.00	91.18\\
217.00	91.71\\
230.00	91.74\\
245.00	92.24\\
261.00	92.34\\
277.00	92.78\\
295.00	92.88\\
314.00	93.31\\
334.00	93.41\\
378.00	93.12\\
402.00	93.53\\
427.00	93.68\\
455.00	93.41\\
484.00	93.80\\
515.00	93.79\\
547.00	93.42\\
582.00	93.47\\
619.00	93.38\\
659.00	93.17\\
701.00	93.44\\
746.00	93.43\\
793.00	93.57\\
844.00	93.60\\
897.00	93.31\\
1016.00	92.91\\
1080.00	93.06\\
1149.00	93.30\\
1222.00	93.37\\
1300.00	93.23\\
1565.00	93.61\\
1771.00	93.51\\
1884.00	93.58\\
2005.00	93.57\\
2413.00	92.83\\
2567.00	92.99\\
2730.00	92.78\\
2905.00	93.01\\
3287.00	93.09\\
3496.00	92.96\\
4209.00	93.32\\
4477.00	93.28\\
5389.00	93.51\\
5733.00	93.55\\
6099.00	93.43\\
6488.00	93.50\\
6901.00	93.41\\
7341.00	93.46\\
8837.00	93.33\\
9401.00	93.37\\
10000.00	93.33\\
};
\addplot [color=gray!85!lightgray, dashed, line width=1.2pt, forget plot]
  table[row sep=crcr]{%
1.00	100.00\\
23.00	100.00\\
25.00	96.00\\
26.00	96.15\\
28.00	92.86\\
32.00	93.75\\
36.00	94.44\\
41.00	95.12\\
46.00	95.65\\
49.00	95.92\\
52.00	94.23\\
59.00	94.92\\
67.00	95.52\\
71.00	94.37\\
76.00	94.74\\
81.00	93.83\\
91.00	94.51\\
97.00	94.85\\
103.00	93.20\\
110.00	92.73\\
117.00	93.16\\
124.00	92.74\\
140.00	93.57\\
159.00	94.34\\
169.00	94.08\\
191.00	94.76\\
204.00	94.12\\
217.00	94.01\\
230.00	94.35\\
245.00	93.88\\
261.00	93.87\\
277.00	94.22\\
295.00	94.24\\
314.00	93.95\\
334.00	93.41\\
355.00	93.52\\
402.00	93.28\\
427.00	93.44\\
455.00	92.97\\
515.00	93.79\\
547.00	93.78\\
619.00	94.51\\
746.00	94.10\\
897.00	94.43\\
955.00	94.35\\
1016.00	94.09\\
1149.00	94.17\\
1222.00	94.52\\
1472.00	94.29\\
1565.00	94.38\\
1665.00	94.29\\
2005.00	94.61\\
2132.00	94.70\\
2268.00	94.49\\
2730.00	94.54\\
2905.00	94.29\\
3287.00	94.25\\
3719.00	94.43\\
3957.00	94.39\\
4477.00	94.13\\
4763.00	93.95\\
5389.00	94.04\\
6099.00	94.03\\
6488.00	93.91\\
7341.00	93.98\\
8837.00	94.05\\
9401.00	93.97\\
10000.00	94.01\\
};
\addplot [color=white!50!darkgray, dotted, line width=1.2pt, forget plot]
  table[row sep=crcr]{%
1.00	0.00\\
2.00	50.00\\
3.00	66.67\\
4.00	75.00\\
5.00	80.00\\
6.00	83.33\\
7.00	85.71\\
8.00	87.50\\
9.00	88.89\\
10.00	90.00\\
11.00	90.91\\
12.00	91.67\\
14.00	92.86\\
16.00	93.75\\
18.00	94.44\\
21.00	95.24\\
23.00	95.65\\
26.00	96.15\\
28.00	96.43\\
30.00	93.33\\
34.00	94.12\\
38.00	94.74\\
41.00	92.68\\
43.00	90.70\\
49.00	91.84\\
52.00	90.38\\
56.00	91.07\\
59.00	89.83\\
67.00	91.04\\
76.00	92.11\\
81.00	91.36\\
91.00	92.31\\
97.00	91.75\\
103.00	92.23\\
117.00	91.45\\
124.00	91.13\\
132.00	91.67\\
140.00	91.43\\
149.00	91.28\\
169.00	92.31\\
191.00	93.19\\
204.00	93.14\\
230.00	92.17\\
245.00	92.24\\
261.00	92.72\\
277.00	92.78\\
314.00	93.63\\
334.00	93.71\\
378.00	94.44\\
402.00	94.78\\
455.00	94.51\\
484.00	94.42\\
515.00	94.56\\
582.00	94.50\\
619.00	94.51\\
659.00	94.69\\
701.00	94.58\\
793.00	94.96\\
955.00	94.87\\
1016.00	94.49\\
1080.00	94.63\\
1149.00	94.87\\
1222.00	94.60\\
1300.00	94.77\\
1383.00	94.58\\
1472.00	94.63\\
1565.00	94.89\\
1665.00	94.83\\
1771.00	94.64\\
1884.00	94.75\\
2005.00	94.66\\
2268.00	94.71\\
3090.00	94.17\\
3287.00	94.31\\
3496.00	94.36\\
3719.00	94.27\\
4477.00	94.35\\
4763.00	94.27\\
6099.00	94.24\\
6901.00	94.16\\
8837.00	94.07\\
9401.00	94.22\\
10000.00	94.19\\
};
\addplot [color=black!5!lightgray, dashdotted, line width=1.2pt, forget plot]
  table[row sep=crcr]{%
1.00	100.00\\
2.00	50.00\\
3.00	33.33\\
4.00	50.00\\
5.00	60.00\\
6.00	66.67\\
7.00	71.43\\
8.00	75.00\\
9.00	77.78\\
10.00	80.00\\
11.00	81.82\\
12.00	83.33\\
13.00	84.62\\
14.00	85.71\\
15.00	86.67\\
17.00	88.24\\
19.00	89.47\\
22.00	90.91\\
23.00	91.30\\
25.00	88.00\\
26.00	88.46\\
28.00	89.29\\
30.00	90.00\\
34.00	91.18\\
36.00	91.67\\
38.00	89.47\\
43.00	90.70\\
46.00	91.30\\
49.00	87.76\\
52.00	88.46\\
59.00	89.83\\
67.00	91.04\\
76.00	92.11\\
86.00	93.02\\
91.00	92.31\\
103.00	93.20\\
110.00	92.73\\
117.00	91.45\\
124.00	91.13\\
140.00	92.14\\
159.00	93.08\\
169.00	93.49\\
191.00	93.19\\
204.00	93.63\\
217.00	93.09\\
245.00	93.88\\
261.00	93.49\\
277.00	93.86\\
355.00	94.08\\
378.00	93.65\\
402.00	94.03\\
455.00	94.29\\
484.00	94.21\\
515.00	93.79\\
547.00	94.15\\
582.00	94.33\\
619.00	94.02\\
659.00	93.93\\
701.00	93.72\\
746.00	93.70\\
793.00	93.82\\
844.00	93.36\\
897.00	93.53\\
955.00	93.51\\
1016.00	93.60\\
1080.00	93.80\\
1222.00	93.54\\
1300.00	93.46\\
1383.00	93.28\\
1472.00	93.48\\
1565.00	93.23\\
1665.00	93.39\\
1771.00	93.45\\
1884.00	93.74\\
2005.00	93.77\\
2132.00	93.48\\
2413.00	93.54\\
2567.00	93.69\\
2730.00	93.63\\
3090.00	93.79\\
3287.00	93.95\\
3496.00	93.76\\
5389.00	93.77\\
5733.00	93.83\\
6099.00	93.72\\
6488.00	93.76\\
6901.00	93.68\\
8307.00	93.84\\
9401.00	93.79\\
10000.00	93.76\\
};
\addplot [color=white!80!black, line width=1.2pt, forget plot]
  table[row sep=crcr]{%
1.00	100.00\\
46.00	100.00\\
49.00	97.96\\
59.00	98.31\\
76.00	98.68\\
97.00	98.97\\
103.00	99.03\\
110.00	98.18\\
117.00	97.44\\
124.00	97.58\\
132.00	96.97\\
159.00	97.48\\
169.00	97.63\\
180.00	97.22\\
204.00	97.55\\
230.00	96.96\\
245.00	97.14\\
277.00	96.75\\
295.00	96.61\\
334.00	97.01\\
355.00	96.62\\
378.00	96.56\\
402.00	96.27\\
427.00	96.49\\
455.00	96.26\\
484.00	95.25\\
515.00	95.53\\
547.00	95.43\\
582.00	94.85\\
619.00	95.15\\
659.00	94.69\\
701.00	94.72\\
746.00	94.10\\
793.00	93.69\\
844.00	93.96\\
897.00	93.87\\
955.00	93.61\\
1016.00	93.70\\
1149.00	93.47\\
1300.00	93.15\\
1383.00	93.49\\
1472.00	93.41\\
1565.00	93.23\\
1771.00	93.51\\
2005.00	93.07\\
2268.00	93.12\\
2567.00	93.03\\
2730.00	93.11\\
2905.00	93.12\\
3496.00	93.42\\
4209.00	93.70\\
4477.00	93.57\\
4763.00	93.53\\
5389.00	93.75\\
6901.00	93.57\\
7809.00	93.58\\
9401.00	93.62\\
10000.00	93.64\\
};
\end{axis}

\begin{axis}[%
width=0.391\figurewidth,
height=0.265\figureheight,
at={(0.54\figurewidth,0.368\figureheight)},
scale only axis,
xmin=1.00,
xmax=10000.00,
xtick={2000.00,4000.00,6000.00,8000.00,10000.00},
xticklabels={\empty},
ymin=0.00,
ymax=105.00,
ytick={0.00,20.00,40.00,60.00,80.00,100.00},
yticklabels={\empty},
axis background/.style={fill=white},
title style={font=\bfseries\color{mycolor1}},
title={\textbf{$\sym{beta}_{12}$}},
xmajorgrids,
ymajorgrids,
grid style={dotted, opacity=0.25}
]
\addplot [color=white!10!black, line width=1.2pt, forget plot]
  table[row sep=crcr]{%
1.00	100.00\\
3.00	100.00\\
4.00	75.00\\
5.00	60.00\\
6.00	66.67\\
7.00	57.14\\
8.00	50.00\\
9.00	44.44\\
10.00	50.00\\
11.00	45.45\\
12.00	41.67\\
13.00	38.46\\
14.00	42.86\\
15.00	40.00\\
16.00	37.50\\
17.00	41.18\\
18.00	38.89\\
19.00	42.11\\
21.00	47.62\\
23.00	52.17\\
25.00	56.00\\
26.00	57.69\\
28.00	60.71\\
30.00	63.33\\
32.00	62.50\\
34.00	58.82\\
38.00	57.89\\
41.00	56.10\\
43.00	55.81\\
46.00	58.70\\
52.00	59.62\\
56.00	58.93\\
59.00	57.63\\
63.00	55.56\\
67.00	56.72\\
76.00	52.63\\
81.00	51.85\\
86.00	54.65\\
91.00	57.14\\
97.00	58.76\\
103.00	58.25\\
110.00	56.36\\
117.00	58.12\\
124.00	58.87\\
140.00	59.29\\
149.00	59.06\\
159.00	58.49\\
169.00	57.40\\
180.00	57.78\\
191.00	57.59\\
204.00	57.84\\
217.00	57.60\\
230.00	58.70\\
245.00	58.78\\
261.00	58.24\\
277.00	59.21\\
295.00	59.32\\
314.00	59.87\\
334.00	59.88\\
355.00	60.00\\
378.00	59.79\\
427.00	60.19\\
455.00	59.34\\
484.00	58.88\\
515.00	58.06\\
547.00	59.05\\
582.00	58.59\\
619.00	58.32\\
659.00	57.51\\
701.00	57.49\\
746.00	56.70\\
793.00	56.49\\
844.00	56.04\\
1016.00	55.81\\
1080.00	55.83\\
1149.00	55.70\\
1222.00	55.89\\
1300.00	56.23\\
1383.00	56.47\\
1472.00	56.52\\
1565.00	56.36\\
1665.00	56.64\\
1771.00	56.30\\
1884.00	56.58\\
2005.00	56.51\\
2132.00	57.08\\
2268.00	57.19\\
2567.00	56.72\\
2905.00	56.59\\
3090.00	56.83\\
3287.00	56.89\\
3496.00	56.75\\
4763.00	57.06\\
5066.00	56.61\\
5389.00	56.24\\
5733.00	56.36\\
6488.00	55.98\\
6901.00	55.99\\
7809.00	56.18\\
8307.00	56.24\\
8837.00	56.11\\
10000.00	56.21\\
};
\addplot [color=black!25!darkgray, dashed, line width=1.2pt, forget plot]
  table[row sep=crcr]{%
1.00	100.00\\
2.00	50.00\\
3.00	66.67\\
4.00	75.00\\
5.00	60.00\\
6.00	66.67\\
7.00	71.43\\
8.00	75.00\\
9.00	66.67\\
10.00	70.00\\
11.00	72.73\\
12.00	75.00\\
13.00	69.23\\
14.00	71.43\\
15.00	73.33\\
16.00	75.00\\
17.00	76.47\\
18.00	77.78\\
19.00	73.68\\
21.00	76.19\\
22.00	72.73\\
23.00	73.91\\
25.00	72.00\\
26.00	69.23\\
28.00	71.43\\
30.00	70.00\\
32.00	65.62\\
34.00	64.71\\
36.00	63.89\\
38.00	65.79\\
41.00	65.85\\
43.00	65.12\\
46.00	67.39\\
49.00	63.27\\
52.00	63.46\\
56.00	62.50\\
59.00	59.32\\
63.00	57.14\\
67.00	53.73\\
71.00	53.52\\
76.00	51.32\\
81.00	53.09\\
86.00	52.33\\
91.00	52.75\\
97.00	52.58\\
103.00	50.49\\
110.00	50.91\\
117.00	50.43\\
124.00	51.61\\
132.00	50.00\\
140.00	49.29\\
149.00	47.65\\
159.00	46.54\\
169.00	48.52\\
180.00	50.00\\
191.00	50.26\\
204.00	51.96\\
217.00	51.61\\
230.00	52.61\\
245.00	52.65\\
261.00	52.11\\
277.00	53.43\\
295.00	52.54\\
314.00	52.87\\
334.00	54.19\\
355.00	54.08\\
378.00	53.70\\
402.00	54.23\\
455.00	54.51\\
484.00	54.75\\
515.00	54.76\\
547.00	54.30\\
582.00	55.15\\
619.00	55.57\\
659.00	55.69\\
701.00	55.35\\
793.00	55.61\\
844.00	54.74\\
897.00	55.18\\
955.00	55.29\\
1016.00	55.12\\
1080.00	55.56\\
1149.00	55.87\\
1222.00	55.89\\
1300.00	55.31\\
1383.00	55.39\\
1472.00	55.64\\
1565.00	55.59\\
1665.00	55.74\\
1771.00	55.62\\
1884.00	56.16\\
2005.00	56.16\\
2268.00	56.48\\
2413.00	56.44\\
2567.00	56.49\\
2730.00	56.67\\
3090.00	56.44\\
3496.00	56.29\\
4209.00	56.50\\
4763.00	56.18\\
5066.00	56.20\\
6099.00	56.71\\
6488.00	56.58\\
7341.00	56.55\\
7809.00	56.68\\
8837.00	56.35\\
9401.00	56.29\\
10000.00	56.50\\
};
\addplot [color=black!45!gray, dotted, line width=1.2pt, forget plot]
  table[row sep=crcr]{%
1.00	100.00\\
6.00	100.00\\
7.00	85.71\\
8.00	87.50\\
9.00	88.89\\
10.00	90.00\\
11.00	81.82\\
12.00	83.33\\
13.00	84.62\\
14.00	78.57\\
15.00	80.00\\
16.00	75.00\\
17.00	76.47\\
18.00	72.22\\
19.00	73.68\\
21.00	71.43\\
22.00	72.73\\
23.00	69.57\\
25.00	72.00\\
26.00	73.08\\
28.00	67.86\\
30.00	66.67\\
32.00	65.62\\
34.00	64.71\\
36.00	63.89\\
38.00	65.79\\
41.00	68.29\\
43.00	69.77\\
46.00	65.22\\
49.00	67.35\\
52.00	69.23\\
56.00	69.64\\
59.00	69.49\\
63.00	68.25\\
67.00	68.66\\
71.00	70.42\\
81.00	71.60\\
86.00	72.09\\
91.00	71.43\\
97.00	70.10\\
103.00	67.96\\
110.00	66.36\\
117.00	64.10\\
124.00	63.71\\
140.00	63.57\\
149.00	62.42\\
159.00	62.26\\
169.00	63.91\\
180.00	62.22\\
191.00	62.30\\
217.00	63.13\\
230.00	62.61\\
245.00	64.08\\
261.00	63.98\\
277.00	63.54\\
295.00	61.02\\
314.00	60.51\\
334.00	61.08\\
355.00	61.13\\
378.00	60.85\\
402.00	60.95\\
427.00	60.42\\
455.00	60.00\\
484.00	59.30\\
515.00	57.86\\
547.00	58.50\\
582.00	58.42\\
619.00	59.29\\
659.00	58.57\\
701.00	58.49\\
746.00	58.85\\
793.00	58.01\\
844.00	58.29\\
897.00	57.97\\
955.00	58.53\\
1016.00	58.07\\
1080.00	58.06\\
1149.00	57.53\\
1222.00	57.69\\
1300.00	57.54\\
1383.00	57.77\\
1472.00	57.81\\
1565.00	57.32\\
1665.00	58.02\\
1771.00	58.50\\
1884.00	58.01\\
2005.00	58.15\\
2132.00	57.83\\
2413.00	57.73\\
2567.00	57.50\\
2730.00	57.47\\
2905.00	57.21\\
3090.00	57.15\\
3287.00	56.86\\
3496.00	57.21\\
3719.00	57.19\\
3957.00	57.29\\
4209.00	57.16\\
4477.00	57.27\\
5066.00	57.07\\
5389.00	57.51\\
5733.00	57.54\\
6099.00	57.71\\
6488.00	57.40\\
6901.00	57.21\\
7341.00	57.16\\
7809.00	57.36\\
8307.00	57.13\\
8837.00	57.08\\
10000.00	57.14\\
};
\addplot [color=white!15!darkgray, dashdotted, line width=1.2pt, forget plot]
  table[row sep=crcr]{%
1.00	0.00\\
3.00	0.00\\
4.00	25.00\\
5.00	20.00\\
6.00	16.67\\
7.00	28.57\\
8.00	25.00\\
9.00	22.22\\
10.00	20.00\\
11.00	18.18\\
12.00	25.00\\
13.00	23.08\\
14.00	28.57\\
15.00	33.33\\
16.00	31.25\\
17.00	35.29\\
18.00	38.89\\
19.00	42.11\\
21.00	42.86\\
23.00	39.13\\
25.00	40.00\\
26.00	38.46\\
28.00	35.71\\
30.00	40.00\\
34.00	41.18\\
36.00	41.67\\
38.00	44.74\\
41.00	46.34\\
43.00	46.51\\
46.00	45.65\\
49.00	46.94\\
52.00	46.15\\
56.00	44.64\\
59.00	47.46\\
67.00	47.76\\
71.00	49.30\\
76.00	47.37\\
86.00	46.51\\
91.00	47.25\\
97.00	49.48\\
103.00	49.51\\
110.00	50.91\\
117.00	50.43\\
124.00	51.61\\
132.00	50.76\\
140.00	51.43\\
149.00	50.34\\
159.00	51.57\\
169.00	52.07\\
180.00	50.00\\
191.00	50.26\\
204.00	51.96\\
230.00	52.17\\
245.00	52.65\\
261.00	50.96\\
277.00	52.35\\
295.00	53.56\\
314.00	52.87\\
334.00	52.40\\
355.00	52.96\\
378.00	53.17\\
402.00	52.74\\
427.00	52.69\\
455.00	52.09\\
515.00	52.04\\
547.00	51.92\\
582.00	50.86\\
619.00	51.05\\
659.00	51.75\\
701.00	52.64\\
746.00	53.35\\
793.00	53.47\\
844.00	53.44\\
897.00	52.95\\
955.00	52.77\\
1016.00	52.17\\
1080.00	52.69\\
1149.00	52.83\\
1222.00	53.27\\
1300.00	53.54\\
1383.00	54.16\\
1472.00	53.74\\
1565.00	53.74\\
1665.00	53.63\\
1771.00	53.81\\
1884.00	53.72\\
2005.00	53.97\\
2132.00	54.50\\
2413.00	54.33\\
2567.00	53.92\\
2730.00	54.25\\
3287.00	54.61\\
3496.00	54.81\\
3719.00	54.77\\
3957.00	54.44\\
4763.00	55.15\\
5389.00	55.43\\
5733.00	55.29\\
6099.00	55.42\\
6488.00	55.69\\
6901.00	55.76\\
7341.00	55.97\\
7809.00	56.01\\
8307.00	55.94\\
8837.00	56.09\\
9401.00	56.14\\
10000.00	55.98\\
};
\addplot [color=white!45!black, line width=1.2pt, forget plot]
  table[row sep=crcr]{%
1.00	0.00\\
2.00	50.00\\
3.00	66.67\\
4.00	50.00\\
5.00	60.00\\
6.00	66.67\\
7.00	57.14\\
8.00	62.50\\
9.00	66.67\\
10.00	60.00\\
11.00	63.64\\
12.00	66.67\\
13.00	69.23\\
14.00	64.29\\
15.00	60.00\\
16.00	62.50\\
17.00	64.71\\
18.00	61.11\\
19.00	57.89\\
21.00	57.14\\
22.00	54.55\\
23.00	56.52\\
25.00	60.00\\
26.00	61.54\\
28.00	64.29\\
30.00	66.67\\
32.00	62.50\\
34.00	61.76\\
36.00	63.89\\
38.00	63.16\\
41.00	63.41\\
43.00	62.79\\
46.00	60.87\\
49.00	57.14\\
52.00	57.69\\
56.00	58.93\\
59.00	59.32\\
63.00	60.32\\
67.00	58.21\\
71.00	57.75\\
76.00	56.58\\
81.00	55.56\\
86.00	58.14\\
91.00	56.04\\
110.00	58.18\\
117.00	55.56\\
124.00	55.65\\
132.00	56.06\\
140.00	55.71\\
149.00	54.36\\
159.00	54.09\\
169.00	54.44\\
180.00	54.44\\
191.00	54.97\\
204.00	55.39\\
217.00	56.22\\
230.00	55.65\\
245.00	55.51\\
261.00	55.94\\
277.00	55.96\\
295.00	55.59\\
314.00	54.78\\
334.00	54.79\\
355.00	55.21\\
378.00	55.56\\
427.00	57.38\\
455.00	56.70\\
484.00	57.02\\
515.00	56.89\\
547.00	57.40\\
582.00	57.73\\
619.00	58.32\\
659.00	58.27\\
701.00	57.77\\
746.00	56.97\\
793.00	56.87\\
844.00	56.87\\
897.00	56.74\\
955.00	56.75\\
1016.00	56.40\\
1080.00	55.83\\
1149.00	55.79\\
1222.00	56.22\\
1300.00	56.23\\
1383.00	56.54\\
1472.00	55.98\\
1665.00	56.04\\
1771.00	56.41\\
1884.00	56.63\\
2005.00	56.66\\
2268.00	55.78\\
2413.00	55.66\\
2730.00	56.01\\
2905.00	55.49\\
3090.00	55.53\\
3287.00	55.83\\
3496.00	55.75\\
3957.00	55.24\\
4209.00	55.33\\
4477.00	55.15\\
4763.00	55.45\\
5066.00	55.35\\
5389.00	55.48\\
6488.00	56.18\\
6901.00	56.30\\
7341.00	56.19\\
8307.00	56.10\\
8837.00	55.97\\
9401.00	56.15\\
10000.00	56.15\\
};
\addplot [color=gray!85!lightgray, dashed, line width=1.2pt, forget plot]
  table[row sep=crcr]{%
1.00	0.00\\
2.00	50.00\\
3.00	66.67\\
4.00	75.00\\
5.00	60.00\\
6.00	66.67\\
7.00	57.14\\
8.00	62.50\\
9.00	55.56\\
10.00	60.00\\
11.00	63.64\\
12.00	66.67\\
13.00	61.54\\
14.00	64.29\\
15.00	66.67\\
16.00	68.75\\
17.00	70.59\\
18.00	72.22\\
19.00	73.68\\
21.00	66.67\\
22.00	68.18\\
23.00	65.22\\
25.00	64.00\\
26.00	65.38\\
28.00	67.86\\
30.00	66.67\\
32.00	68.75\\
34.00	70.59\\
36.00	66.67\\
38.00	65.79\\
41.00	63.41\\
43.00	62.79\\
46.00	63.04\\
49.00	59.18\\
52.00	61.54\\
56.00	62.50\\
59.00	59.32\\
63.00	61.90\\
67.00	61.19\\
71.00	61.97\\
86.00	61.63\\
91.00	60.44\\
97.00	59.79\\
103.00	60.19\\
110.00	60.00\\
117.00	60.68\\
124.00	58.87\\
132.00	58.33\\
140.00	60.00\\
149.00	59.06\\
159.00	61.01\\
169.00	59.17\\
180.00	60.00\\
191.00	59.69\\
204.00	58.82\\
217.00	58.53\\
230.00	60.00\\
245.00	60.00\\
261.00	58.62\\
277.00	59.57\\
295.00	60.00\\
314.00	59.87\\
334.00	58.98\\
355.00	57.46\\
378.00	57.67\\
402.00	57.71\\
427.00	58.55\\
455.00	58.46\\
484.00	59.30\\
515.00	59.03\\
582.00	58.76\\
619.00	59.61\\
659.00	59.18\\
701.00	58.63\\
746.00	58.18\\
793.00	58.01\\
844.00	57.58\\
897.00	57.64\\
955.00	57.59\\
1016.00	57.28\\
1149.00	57.18\\
1222.00	57.36\\
1300.00	57.85\\
1383.00	57.56\\
1472.00	57.54\\
1565.00	57.76\\
1771.00	57.26\\
1884.00	57.43\\
2005.00	58.30\\
2268.00	58.47\\
2413.00	58.97\\
2730.00	59.34\\
2905.00	59.83\\
3090.00	59.61\\
3287.00	59.93\\
3719.00	60.15\\
3957.00	60.20\\
4209.00	60.16\\
4477.00	59.91\\
4763.00	59.90\\
5389.00	60.18\\
5733.00	60.42\\
6099.00	60.22\\
6488.00	60.16\\
6901.00	60.01\\
7341.00	60.05\\
7809.00	60.19\\
8307.00	60.21\\
8837.00	60.34\\
10000.00	60.46\\
};
\addplot [color=white!50!darkgray, dotted, line width=1.2pt, forget plot]
  table[row sep=crcr]{%
1.00	100.00\\
4.00	100.00\\
5.00	80.00\\
6.00	83.33\\
7.00	85.71\\
8.00	75.00\\
9.00	66.67\\
10.00	60.00\\
11.00	54.55\\
12.00	58.33\\
13.00	53.85\\
14.00	57.14\\
15.00	53.33\\
16.00	50.00\\
17.00	52.94\\
18.00	55.56\\
19.00	52.63\\
21.00	52.38\\
22.00	54.55\\
23.00	52.17\\
25.00	52.00\\
26.00	50.00\\
28.00	50.00\\
30.00	46.67\\
32.00	43.75\\
34.00	47.06\\
36.00	47.22\\
38.00	50.00\\
41.00	51.22\\
43.00	51.16\\
46.00	50.00\\
49.00	48.98\\
52.00	50.00\\
56.00	48.21\\
59.00	49.15\\
63.00	50.79\\
71.00	50.70\\
76.00	51.32\\
81.00	50.62\\
86.00	51.16\\
91.00	50.55\\
97.00	50.52\\
103.00	53.40\\
117.00	53.85\\
124.00	55.65\\
132.00	54.55\\
140.00	52.86\\
149.00	51.68\\
159.00	52.83\\
169.00	53.25\\
180.00	53.89\\
191.00	52.36\\
204.00	51.47\\
217.00	51.61\\
230.00	52.17\\
245.00	51.02\\
261.00	50.19\\
295.00	50.17\\
314.00	50.32\\
334.00	51.80\\
355.00	52.39\\
378.00	52.91\\
402.00	53.98\\
427.00	54.33\\
455.00	54.51\\
484.00	54.96\\
515.00	54.37\\
547.00	54.48\\
582.00	54.47\\
619.00	54.77\\
659.00	54.63\\
701.00	54.78\\
746.00	56.03\\
793.00	55.86\\
844.00	56.16\\
897.00	56.08\\
955.00	56.86\\
1016.00	57.28\\
1080.00	56.67\\
1222.00	56.30\\
1300.00	55.92\\
1383.00	55.68\\
1472.00	56.25\\
1565.00	55.85\\
1665.00	55.68\\
1771.00	56.13\\
1884.00	56.21\\
2005.00	55.71\\
2132.00	56.10\\
2268.00	56.26\\
2413.00	56.24\\
2567.00	56.45\\
2730.00	57.11\\
2905.00	56.73\\
3090.00	56.96\\
3287.00	56.95\\
3496.00	56.81\\
3719.00	56.84\\
3957.00	57.01\\
4209.00	56.69\\
4477.00	56.58\\
4763.00	56.81\\
5066.00	56.65\\
5389.00	56.63\\
5733.00	56.51\\
6099.00	56.24\\
6901.00	56.38\\
8837.00	56.49\\
9401.00	56.37\\
10000.00	56.38\\
};
\addplot [color=black!5!lightgray, dashdotted, line width=1.2pt, forget plot]
  table[row sep=crcr]{%
1.00	0.00\\
2.00	0.00\\
3.00	33.33\\
4.00	50.00\\
5.00	40.00\\
6.00	33.33\\
7.00	42.86\\
8.00	50.00\\
9.00	44.44\\
10.00	50.00\\
11.00	54.55\\
12.00	50.00\\
13.00	46.15\\
14.00	50.00\\
15.00	53.33\\
16.00	50.00\\
17.00	52.94\\
18.00	55.56\\
19.00	57.89\\
21.00	61.90\\
22.00	59.09\\
23.00	56.52\\
25.00	56.00\\
26.00	53.85\\
28.00	53.57\\
30.00	56.67\\
32.00	56.25\\
34.00	58.82\\
36.00	61.11\\
38.00	60.53\\
41.00	58.54\\
43.00	55.81\\
46.00	54.35\\
49.00	55.10\\
52.00	55.77\\
56.00	53.57\\
59.00	52.54\\
63.00	53.97\\
67.00	56.72\\
71.00	56.34\\
76.00	55.26\\
86.00	55.81\\
91.00	53.85\\
97.00	52.58\\
103.00	52.43\\
110.00	53.64\\
117.00	54.70\\
124.00	56.45\\
132.00	56.06\\
149.00	58.39\\
159.00	57.86\\
169.00	56.21\\
180.00	56.11\\
191.00	53.93\\
204.00	53.43\\
217.00	54.84\\
230.00	53.91\\
245.00	54.29\\
261.00	53.26\\
277.00	53.43\\
295.00	54.24\\
314.00	53.50\\
334.00	54.49\\
355.00	55.21\\
378.00	55.03\\
402.00	54.73\\
427.00	54.80\\
455.00	54.73\\
484.00	54.96\\
515.00	54.76\\
547.00	54.66\\
619.00	54.60\\
659.00	55.08\\
701.00	54.49\\
746.00	54.29\\
793.00	54.60\\
844.00	54.50\\
897.00	55.41\\
955.00	54.97\\
1016.00	54.72\\
1080.00	55.09\\
1149.00	55.00\\
1222.00	55.32\\
1300.00	55.92\\
1472.00	55.64\\
1565.00	56.10\\
1665.00	56.46\\
1771.00	56.35\\
1884.00	56.53\\
2132.00	56.14\\
2268.00	56.53\\
2413.00	56.36\\
2567.00	56.88\\
2730.00	56.37\\
2905.00	56.76\\
3090.00	56.96\\
3287.00	56.37\\
3496.00	56.64\\
3719.00	56.68\\
3957.00	56.99\\
4477.00	57.14\\
4763.00	57.09\\
5389.00	56.76\\
6099.00	56.67\\
6488.00	56.43\\
6901.00	56.56\\
7809.00	56.59\\
8307.00	56.46\\
9401.00	56.45\\
10000.00	56.37\\
};
\addplot [color=white!80!black, line width=1.2pt, forget plot]
  table[row sep=crcr]{%
1.00	100.00\\
2.00	100.00\\
3.00	66.67\\
4.00	75.00\\
5.00	80.00\\
6.00	66.67\\
7.00	57.14\\
8.00	50.00\\
9.00	44.44\\
10.00	40.00\\
11.00	45.45\\
12.00	41.67\\
13.00	46.15\\
14.00	42.86\\
15.00	46.67\\
16.00	43.75\\
17.00	47.06\\
18.00	50.00\\
19.00	47.37\\
21.00	47.62\\
22.00	50.00\\
23.00	47.83\\
25.00	48.00\\
26.00	50.00\\
28.00	50.00\\
30.00	53.33\\
32.00	56.25\\
34.00	55.88\\
36.00	52.78\\
38.00	55.26\\
41.00	53.66\\
43.00	53.49\\
46.00	54.35\\
49.00	55.10\\
52.00	55.77\\
56.00	58.93\\
59.00	57.63\\
63.00	55.56\\
67.00	53.73\\
71.00	54.93\\
81.00	55.56\\
86.00	55.81\\
91.00	57.14\\
97.00	55.67\\
103.00	57.28\\
110.00	57.27\\
117.00	58.12\\
124.00	58.06\\
132.00	58.33\\
140.00	57.86\\
149.00	57.05\\
159.00	57.86\\
169.00	58.58\\
191.00	59.16\\
204.00	59.80\\
217.00	60.83\\
230.00	60.87\\
261.00	59.00\\
277.00	59.21\\
295.00	58.31\\
314.00	59.87\\
334.00	58.68\\
355.00	58.87\\
378.00	57.94\\
427.00	58.55\\
455.00	58.24\\
484.00	58.26\\
515.00	59.22\\
547.00	59.23\\
582.00	59.11\\
619.00	59.13\\
659.00	58.42\\
701.00	58.77\\
746.00	58.45\\
793.00	58.76\\
844.00	58.77\\
897.00	58.97\\
955.00	59.69\\
1016.00	59.74\\
1080.00	59.44\\
1149.00	59.62\\
1222.00	59.25\\
1300.00	59.46\\
1383.00	58.86\\
1472.00	59.10\\
1665.00	58.92\\
1771.00	59.06\\
1884.00	58.86\\
2005.00	58.40\\
2132.00	58.49\\
2268.00	58.33\\
2413.00	58.02\\
2567.00	58.01\\
2730.00	57.88\\
2905.00	58.14\\
3287.00	57.71\\
3496.00	58.07\\
3719.00	57.92\\
3957.00	57.34\\
4209.00	57.23\\
4477.00	56.96\\
4763.00	56.92\\
5066.00	57.13\\
5733.00	56.99\\
6099.00	56.65\\
7341.00	56.68\\
7809.00	56.58\\
8307.00	56.78\\
9401.00	56.49\\
10000.00	56.55\\
};
\end{axis}

\begin{axis}[%
width=0.391\figurewidth,
height=0.265\figureheight,
at={(0\figurewidth,0\figureheight)},
scale only axis,
xmin=1.00,
xmax=10000.00,
xlabel style={font=\color{mycolor1}},
xlabel={Iterationen $\sym{n}_{\sym{idx_MC}}$},
ymin=0.00,
ymax=105.00,
ylabel style={font=\color{mycolor1}},
ylabel={\sym{power} [\%]},
axis background/.style={fill=white},
title style={font=\bfseries\color{mycolor1}},
title={\textbf{$\sym{beta}_{11}$}},
xmajorgrids,
ymajorgrids,
grid style={dotted, opacity=0.25}
]
\addplot [color=white!10!black, line width=1.2pt, forget plot]
  table[row sep=crcr]{%
1.00	100.00\\
2.00	50.00\\
3.00	33.33\\
4.00	50.00\\
5.00	60.00\\
6.00	66.67\\
7.00	57.14\\
8.00	62.50\\
9.00	66.67\\
10.00	70.00\\
11.00	63.64\\
12.00	66.67\\
13.00	61.54\\
14.00	64.29\\
15.00	66.67\\
16.00	68.75\\
17.00	70.59\\
18.00	66.67\\
19.00	63.16\\
21.00	61.90\\
22.00	59.09\\
23.00	60.87\\
25.00	60.00\\
26.00	61.54\\
28.00	64.29\\
30.00	63.33\\
32.00	65.62\\
34.00	67.65\\
36.00	69.44\\
38.00	71.05\\
41.00	70.73\\
43.00	67.44\\
46.00	69.57\\
49.00	67.35\\
52.00	65.38\\
56.00	62.50\\
59.00	64.41\\
63.00	65.08\\
67.00	67.16\\
71.00	67.61\\
76.00	69.74\\
86.00	70.93\\
91.00	71.43\\
97.00	69.07\\
103.00	69.90\\
117.00	70.09\\
124.00	70.97\\
132.00	69.70\\
140.00	70.00\\
149.00	69.13\\
159.00	69.18\\
169.00	68.64\\
180.00	68.89\\
191.00	69.63\\
204.00	69.61\\
217.00	68.20\\
230.00	68.70\\
245.00	68.57\\
261.00	68.58\\
277.00	69.31\\
295.00	68.81\\
314.00	67.83\\
334.00	68.86\\
355.00	69.01\\
378.00	68.52\\
402.00	68.91\\
427.00	69.56\\
455.00	69.01\\
515.00	68.54\\
547.00	68.92\\
582.00	68.38\\
659.00	67.68\\
701.00	67.48\\
746.00	67.69\\
793.00	67.72\\
844.00	67.54\\
897.00	67.00\\
955.00	67.02\\
1016.00	66.54\\
1080.00	66.76\\
1149.00	66.67\\
1222.00	66.94\\
1300.00	66.77\\
1383.00	66.81\\
1472.00	67.46\\
1565.00	67.28\\
1665.00	67.27\\
1771.00	67.02\\
1884.00	67.30\\
2005.00	67.28\\
2132.00	67.45\\
2413.00	67.43\\
2567.00	67.24\\
2730.00	67.51\\
2905.00	67.30\\
3090.00	67.38\\
3287.00	68.12\\
3496.00	67.96\\
4209.00	67.81\\
4477.00	68.10\\
4763.00	68.30\\
5066.00	68.12\\
5733.00	68.53\\
6099.00	68.59\\
6901.00	68.48\\
7341.00	68.53\\
8837.00	68.42\\
10000.00	68.59\\
};
\addplot [color=black!25!darkgray, dashed, line width=1.2pt, forget plot]
  table[row sep=crcr]{%
1.00	100.00\\
2.00	100.00\\
3.00	66.67\\
4.00	75.00\\
5.00	60.00\\
6.00	66.67\\
7.00	71.43\\
8.00	75.00\\
9.00	77.78\\
10.00	70.00\\
11.00	63.64\\
12.00	66.67\\
13.00	61.54\\
14.00	64.29\\
15.00	66.67\\
16.00	68.75\\
17.00	70.59\\
18.00	72.22\\
19.00	73.68\\
21.00	71.43\\
23.00	73.91\\
25.00	72.00\\
26.00	69.23\\
28.00	71.43\\
30.00	70.00\\
32.00	71.88\\
34.00	73.53\\
36.00	72.22\\
38.00	73.68\\
41.00	73.17\\
43.00	74.42\\
46.00	71.74\\
49.00	69.39\\
52.00	71.15\\
56.00	69.64\\
59.00	71.19\\
63.00	69.84\\
67.00	70.15\\
71.00	71.83\\
76.00	71.05\\
86.00	72.09\\
91.00	72.53\\
97.00	71.13\\
103.00	69.90\\
110.00	69.09\\
124.00	67.74\\
132.00	69.70\\
140.00	67.86\\
149.00	67.11\\
169.00	67.46\\
180.00	66.67\\
191.00	66.49\\
217.00	66.82\\
230.00	67.39\\
245.00	67.76\\
261.00	67.05\\
277.00	66.79\\
295.00	66.10\\
314.00	65.61\\
334.00	65.87\\
355.00	65.92\\
378.00	66.14\\
402.00	66.67\\
427.00	66.74\\
455.00	67.69\\
484.00	66.94\\
515.00	67.77\\
547.00	67.82\\
582.00	68.21\\
619.00	67.69\\
659.00	67.37\\
701.00	67.33\\
746.00	67.02\\
897.00	67.56\\
955.00	67.23\\
1080.00	67.04\\
1149.00	66.58\\
1222.00	66.69\\
1300.00	66.08\\
1383.00	65.58\\
1472.00	65.29\\
1565.00	65.11\\
1665.00	65.41\\
1771.00	65.10\\
1884.00	65.50\\
2005.00	65.19\\
2132.00	65.15\\
2268.00	64.95\\
2413.00	64.90\\
2567.00	64.59\\
2730.00	63.96\\
2905.00	63.92\\
3090.00	64.08\\
3287.00	63.86\\
3496.00	64.04\\
3957.00	64.22\\
4209.00	64.31\\
6099.00	64.40\\
6488.00	64.60\\
7341.00	64.39\\
7809.00	64.52\\
10000.00	64.51\\
};
\addplot [color=black!45!gray, dotted, line width=1.2pt, forget plot]
  table[row sep=crcr]{%
1.00	100.00\\
2.00	50.00\\
3.00	66.67\\
4.00	75.00\\
5.00	80.00\\
6.00	83.33\\
7.00	85.71\\
8.00	87.50\\
9.00	88.89\\
10.00	90.00\\
11.00	90.91\\
12.00	91.67\\
13.00	84.62\\
14.00	85.71\\
15.00	86.67\\
16.00	81.25\\
17.00	76.47\\
18.00	77.78\\
19.00	73.68\\
21.00	76.19\\
22.00	72.73\\
23.00	73.91\\
25.00	76.00\\
26.00	76.92\\
28.00	75.00\\
30.00	76.67\\
32.00	75.00\\
34.00	76.47\\
36.00	75.00\\
38.00	71.05\\
41.00	70.73\\
43.00	69.77\\
46.00	71.74\\
49.00	73.47\\
52.00	73.08\\
56.00	73.21\\
59.00	72.88\\
71.00	73.24\\
76.00	72.37\\
81.00	70.37\\
86.00	72.09\\
91.00	72.53\\
97.00	73.20\\
103.00	71.84\\
110.00	72.73\\
124.00	72.58\\
140.00	72.86\\
149.00	71.81\\
159.00	71.70\\
169.00	71.01\\
191.00	70.16\\
204.00	72.06\\
217.00	72.35\\
230.00	70.87\\
245.00	71.02\\
261.00	72.41\\
277.00	73.29\\
295.00	73.56\\
314.00	73.57\\
334.00	72.16\\
355.00	70.99\\
378.00	70.11\\
402.00	69.65\\
427.00	70.26\\
455.00	70.77\\
484.00	70.45\\
515.00	70.49\\
547.00	70.38\\
582.00	70.45\\
619.00	70.27\\
659.00	70.41\\
701.00	71.18\\
746.00	70.91\\
793.00	70.49\\
844.00	71.21\\
897.00	71.57\\
955.00	72.25\\
1016.00	72.15\\
1080.00	72.69\\
1149.00	72.67\\
1222.00	72.75\\
1300.00	72.69\\
1383.00	73.10\\
1472.00	73.44\\
1565.00	73.16\\
1771.00	72.78\\
1884.00	73.04\\
2005.00	72.52\\
2132.00	72.70\\
2268.00	72.75\\
2413.00	73.23\\
2567.00	73.35\\
2730.00	73.59\\
2905.00	74.11\\
3090.00	73.95\\
3287.00	73.65\\
3719.00	73.60\\
3957.00	73.29\\
4477.00	73.67\\
4763.00	73.67\\
6099.00	74.39\\
8837.00	74.29\\
9401.00	74.16\\
10000.00	74.19\\
};
\addplot [color=white!15!darkgray, dashdotted, line width=1.2pt, forget plot]
  table[row sep=crcr]{%
1.00	100.00\\
2.00	50.00\\
3.00	66.67\\
4.00	75.00\\
5.00	80.00\\
6.00	83.33\\
7.00	85.71\\
8.00	87.50\\
9.00	88.89\\
10.00	80.00\\
11.00	81.82\\
12.00	75.00\\
13.00	76.92\\
14.00	78.57\\
15.00	73.33\\
16.00	68.75\\
17.00	70.59\\
18.00	66.67\\
19.00	68.42\\
21.00	61.90\\
22.00	63.64\\
23.00	60.87\\
25.00	64.00\\
26.00	65.38\\
28.00	67.86\\
30.00	66.67\\
32.00	65.62\\
34.00	67.65\\
36.00	63.89\\
38.00	63.16\\
41.00	63.41\\
43.00	65.12\\
52.00	65.38\\
56.00	66.07\\
59.00	64.41\\
63.00	65.08\\
67.00	64.18\\
71.00	63.38\\
76.00	63.16\\
81.00	61.73\\
86.00	62.79\\
91.00	63.74\\
97.00	63.92\\
103.00	62.14\\
117.00	61.54\\
124.00	62.10\\
132.00	63.64\\
140.00	65.00\\
149.00	64.43\\
159.00	64.78\\
169.00	63.31\\
180.00	63.33\\
191.00	62.30\\
204.00	62.25\\
217.00	62.67\\
230.00	63.48\\
245.00	64.08\\
261.00	63.98\\
277.00	64.62\\
295.00	64.75\\
314.00	64.65\\
334.00	65.27\\
355.00	65.35\\
378.00	65.08\\
402.00	64.18\\
427.00	64.17\\
455.00	64.40\\
484.00	64.46\\
515.00	63.69\\
547.00	63.25\\
701.00	63.20\\
746.00	61.80\\
793.00	62.30\\
844.00	61.97\\
897.00	61.76\\
955.00	61.99\\
1016.00	61.81\\
1080.00	61.85\\
1149.00	61.53\\
1222.00	61.87\\
1300.00	61.85\\
1383.00	62.11\\
1472.00	61.55\\
1565.00	61.85\\
1665.00	62.04\\
1884.00	61.73\\
2132.00	62.29\\
2413.00	62.58\\
2567.00	62.37\\
2730.00	62.38\\
2905.00	62.31\\
3287.00	62.49\\
3496.00	62.24\\
3719.00	62.41\\
3957.00	62.14\\
4209.00	62.01\\
4477.00	62.41\\
4763.00	62.08\\
5066.00	61.98\\
5389.00	62.18\\
6099.00	61.65\\
6901.00	61.72\\
7341.00	61.98\\
7809.00	61.92\\
8837.00	62.13\\
10000.00	61.84\\
};
\addplot [color=white!45!black, line width=1.2pt, forget plot]
  table[row sep=crcr]{%
1.00	0.00\\
2.00	50.00\\
3.00	33.33\\
4.00	50.00\\
5.00	40.00\\
6.00	50.00\\
7.00	57.14\\
8.00	62.50\\
9.00	66.67\\
10.00	70.00\\
11.00	63.64\\
12.00	66.67\\
13.00	69.23\\
14.00	64.29\\
15.00	66.67\\
16.00	62.50\\
17.00	58.82\\
18.00	61.11\\
19.00	57.89\\
21.00	57.14\\
22.00	54.55\\
23.00	52.17\\
25.00	52.00\\
26.00	53.85\\
28.00	57.14\\
30.00	56.67\\
32.00	59.38\\
34.00	61.76\\
38.00	60.53\\
41.00	58.54\\
43.00	58.14\\
46.00	58.70\\
49.00	57.14\\
52.00	57.69\\
56.00	58.93\\
59.00	59.32\\
63.00	61.90\\
67.00	61.19\\
71.00	63.38\\
76.00	63.16\\
81.00	60.49\\
86.00	60.47\\
91.00	59.34\\
97.00	59.79\\
103.00	59.22\\
110.00	57.27\\
117.00	56.41\\
124.00	56.45\\
132.00	59.09\\
140.00	60.00\\
149.00	60.40\\
159.00	61.64\\
169.00	60.36\\
180.00	62.22\\
204.00	64.71\\
217.00	64.98\\
230.00	64.78\\
245.00	65.31\\
261.00	66.28\\
277.00	66.06\\
295.00	67.80\\
314.00	68.15\\
334.00	68.26\\
355.00	68.73\\
378.00	68.52\\
402.00	69.65\\
427.00	70.26\\
455.00	70.11\\
484.00	70.04\\
547.00	70.93\\
582.00	70.79\\
619.00	70.92\\
659.00	70.56\\
701.00	70.47\\
746.00	70.91\\
793.00	70.49\\
844.00	70.85\\
955.00	70.47\\
1080.00	70.83\\
1149.00	71.37\\
1222.00	71.28\\
1300.00	71.08\\
1383.00	70.72\\
1472.00	71.06\\
1565.00	70.73\\
1665.00	70.75\\
1771.00	70.86\\
1884.00	70.70\\
2005.00	70.42\\
2268.00	70.72\\
2413.00	70.70\\
2567.00	70.55\\
2905.00	70.43\\
3090.00	70.26\\
3287.00	69.88\\
3496.00	69.59\\
4209.00	69.45\\
4477.00	69.15\\
4763.00	69.22\\
5389.00	69.01\\
6099.00	69.32\\
6488.00	69.39\\
6901.00	69.37\\
7341.00	69.19\\
9401.00	69.14\\
10000.00	69.26\\
};
\addplot [color=gray!85!lightgray, dashed, line width=1.2pt, forget plot]
  table[row sep=crcr]{%
1.00	0.00\\
2.00	50.00\\
3.00	66.67\\
4.00	50.00\\
5.00	40.00\\
6.00	33.33\\
7.00	42.86\\
8.00	50.00\\
9.00	55.56\\
10.00	60.00\\
11.00	54.55\\
12.00	58.33\\
13.00	61.54\\
14.00	64.29\\
15.00	66.67\\
16.00	62.50\\
17.00	58.82\\
18.00	61.11\\
19.00	63.16\\
21.00	61.90\\
22.00	63.64\\
23.00	60.87\\
25.00	60.00\\
26.00	61.54\\
28.00	64.29\\
30.00	66.67\\
32.00	62.50\\
34.00	61.76\\
36.00	63.89\\
38.00	65.79\\
41.00	63.41\\
43.00	62.79\\
46.00	63.04\\
49.00	59.18\\
52.00	57.69\\
56.00	58.93\\
59.00	59.32\\
63.00	57.14\\
67.00	53.73\\
71.00	53.52\\
76.00	55.26\\
81.00	58.02\\
86.00	59.30\\
91.00	60.44\\
97.00	60.82\\
103.00	60.19\\
110.00	60.91\\
117.00	60.68\\
124.00	61.29\\
132.00	61.36\\
140.00	62.86\\
159.00	64.78\\
169.00	65.09\\
191.00	64.92\\
204.00	64.22\\
217.00	64.06\\
230.00	63.48\\
245.00	64.08\\
261.00	64.37\\
295.00	64.07\\
314.00	64.65\\
334.00	64.97\\
355.00	65.07\\
378.00	65.61\\
402.00	65.67\\
427.00	66.28\\
455.00	66.81\\
484.00	66.12\\
515.00	65.83\\
582.00	65.81\\
619.00	65.43\\
659.00	64.49\\
701.00	64.76\\
746.00	64.88\\
793.00	64.69\\
897.00	63.55\\
955.00	62.72\\
1016.00	63.78\\
1080.00	64.35\\
1149.00	64.75\\
1300.00	64.38\\
1383.00	64.28\\
1472.00	64.88\\
1565.00	65.24\\
1665.00	65.47\\
1771.00	65.56\\
1884.00	65.50\\
2005.00	65.84\\
2132.00	65.81\\
2413.00	65.52\\
2567.00	65.56\\
2730.00	65.42\\
2905.00	65.51\\
3090.00	65.40\\
3287.00	65.38\\
3496.00	65.59\\
3719.00	65.93\\
4477.00	66.41\\
5066.00	66.56\\
5389.00	66.38\\
5733.00	66.28\\
6099.00	65.86\\
6488.00	66.15\\
6901.00	66.08\\
8307.00	66.08\\
8837.00	66.05\\
10000.00	66.21\\
};
\addplot [color=white!50!darkgray, dotted, line width=1.2pt, forget plot]
  table[row sep=crcr]{%
1.00	100.00\\
2.00	50.00\\
3.00	66.67\\
4.00	75.00\\
5.00	60.00\\
6.00	66.67\\
7.00	71.43\\
8.00	62.50\\
9.00	66.67\\
10.00	70.00\\
11.00	72.73\\
12.00	66.67\\
13.00	61.54\\
14.00	64.29\\
15.00	60.00\\
16.00	62.50\\
17.00	64.71\\
18.00	66.67\\
19.00	63.16\\
21.00	61.90\\
23.00	65.22\\
25.00	64.00\\
26.00	65.38\\
28.00	64.29\\
30.00	63.33\\
32.00	62.50\\
36.00	61.11\\
38.00	60.53\\
41.00	63.41\\
43.00	65.12\\
46.00	65.22\\
49.00	67.35\\
52.00	69.23\\
56.00	69.64\\
59.00	67.80\\
63.00	65.08\\
67.00	65.67\\
71.00	64.79\\
81.00	64.20\\
86.00	62.79\\
91.00	62.64\\
97.00	61.86\\
103.00	62.14\\
110.00	60.00\\
117.00	61.54\\
124.00	62.90\\
132.00	62.88\\
140.00	62.14\\
149.00	62.42\\
159.00	63.52\\
169.00	62.72\\
180.00	63.33\\
191.00	63.35\\
204.00	64.22\\
217.00	63.59\\
230.00	63.91\\
245.00	65.31\\
261.00	65.90\\
277.00	65.34\\
314.00	66.88\\
334.00	67.66\\
355.00	66.76\\
378.00	66.40\\
402.00	66.42\\
427.00	66.04\\
455.00	65.49\\
484.00	65.50\\
515.00	64.08\\
547.00	64.72\\
582.00	65.29\\
619.00	64.94\\
659.00	65.10\\
701.00	64.91\\
746.00	65.28\\
793.00	64.69\\
844.00	64.81\\
897.00	64.77\\
955.00	65.34\\
1016.00	65.26\\
1080.00	64.91\\
1383.00	65.80\\
1565.00	66.01\\
1665.00	65.65\\
1771.00	66.01\\
1884.00	66.24\\
2005.00	66.63\\
2132.00	66.74\\
2413.00	67.22\\
2567.00	66.97\\
2730.00	67.14\\
2905.00	67.16\\
3090.00	67.48\\
3287.00	67.54\\
3496.00	67.79\\
3719.00	67.87\\
3957.00	67.68\\
4209.00	67.71\\
4763.00	68.05\\
5066.00	67.88\\
5389.00	67.88\\
5733.00	68.06\\
7341.00	67.59\\
10000.00	67.65\\
};
\addplot [color=black!5!lightgray, dashdotted, line width=1.2pt, forget plot]
  table[row sep=crcr]{%
1.00	100.00\\
2.00	100.00\\
3.00	66.67\\
4.00	75.00\\
5.00	60.00\\
6.00	66.67\\
7.00	71.43\\
8.00	75.00\\
9.00	77.78\\
10.00	70.00\\
11.00	72.73\\
12.00	75.00\\
13.00	76.92\\
14.00	78.57\\
15.00	80.00\\
16.00	81.25\\
17.00	82.35\\
19.00	84.21\\
21.00	85.71\\
23.00	86.96\\
25.00	88.00\\
26.00	84.62\\
28.00	82.14\\
30.00	80.00\\
32.00	81.25\\
34.00	76.47\\
36.00	77.78\\
38.00	76.32\\
41.00	78.05\\
43.00	76.74\\
46.00	73.91\\
49.00	69.39\\
52.00	71.15\\
56.00	69.64\\
59.00	67.80\\
63.00	66.67\\
67.00	67.16\\
71.00	66.20\\
76.00	65.79\\
81.00	67.90\\
86.00	68.60\\
91.00	67.03\\
97.00	65.98\\
103.00	65.05\\
110.00	62.73\\
117.00	62.39\\
124.00	62.90\\
132.00	61.36\\
140.00	61.43\\
149.00	61.74\\
159.00	61.64\\
169.00	62.72\\
180.00	62.22\\
191.00	60.21\\
204.00	60.29\\
230.00	59.57\\
245.00	60.41\\
261.00	61.69\\
277.00	60.29\\
295.00	60.34\\
314.00	59.55\\
334.00	59.88\\
378.00	60.05\\
402.00	59.95\\
427.00	59.95\\
455.00	60.66\\
484.00	61.16\\
515.00	60.19\\
547.00	60.88\\
582.00	61.51\\
619.00	61.55\\
659.00	60.70\\
701.00	60.91\\
746.00	60.99\\
793.00	61.16\\
897.00	61.09\\
955.00	61.36\\
1016.00	60.73\\
1080.00	60.19\\
1300.00	60.08\\
1472.00	60.19\\
1565.00	59.68\\
1771.00	59.85\\
1884.00	60.19\\
2132.00	59.52\\
2268.00	59.35\\
2413.00	59.35\\
2567.00	59.17\\
2905.00	59.00\\
3090.00	59.81\\
3287.00	59.45\\
3496.00	59.15\\
3719.00	59.37\\
4209.00	59.25\\
4477.00	59.26\\
4763.00	59.08\\
5066.00	59.16\\
5389.00	58.84\\
5733.00	58.80\\
6099.00	58.47\\
6488.00	58.42\\
7341.00	58.66\\
7809.00	58.56\\
8307.00	58.63\\
8837.00	58.55\\
9401.00	58.62\\
10000.00	58.46\\
};
\addplot [color=white!80!black, line width=1.2pt, forget plot]
  table[row sep=crcr]{%
1.00	100.00\\
6.00	100.00\\
7.00	85.71\\
8.00	75.00\\
9.00	77.78\\
10.00	80.00\\
11.00	81.82\\
12.00	83.33\\
13.00	84.62\\
14.00	85.71\\
15.00	80.00\\
16.00	75.00\\
17.00	70.59\\
18.00	66.67\\
19.00	68.42\\
21.00	66.67\\
22.00	68.18\\
23.00	65.22\\
25.00	64.00\\
26.00	61.54\\
28.00	60.71\\
30.00	63.33\\
32.00	65.62\\
34.00	61.76\\
36.00	58.33\\
38.00	57.89\\
41.00	58.54\\
43.00	60.47\\
46.00	58.70\\
49.00	61.22\\
52.00	63.46\\
56.00	64.29\\
59.00	62.71\\
63.00	63.49\\
67.00	62.69\\
71.00	63.38\\
86.00	62.79\\
91.00	63.74\\
103.00	64.08\\
110.00	66.36\\
117.00	67.52\\
124.00	67.74\\
132.00	68.18\\
149.00	70.47\\
159.00	71.07\\
169.00	70.41\\
180.00	68.89\\
191.00	69.63\\
204.00	69.12\\
230.00	70.00\\
245.00	71.43\\
277.00	71.12\\
295.00	69.83\\
314.00	70.70\\
334.00	70.36\\
355.00	69.30\\
378.00	69.31\\
402.00	69.15\\
427.00	70.02\\
455.00	69.67\\
484.00	68.39\\
515.00	68.35\\
547.00	68.01\\
582.00	67.01\\
619.00	66.56\\
659.00	66.62\\
701.00	66.48\\
746.00	66.89\\
793.00	66.33\\
844.00	66.71\\
897.00	67.00\\
1016.00	66.83\\
1149.00	67.28\\
1222.00	67.02\\
1300.00	66.92\\
1383.00	66.67\\
1565.00	67.03\\
1665.00	66.85\\
1771.00	66.80\\
1884.00	66.61\\
2268.00	67.50\\
2567.00	67.47\\
2905.00	67.81\\
3090.00	67.73\\
3287.00	67.48\\
3496.00	67.51\\
3719.00	67.41\\
4209.00	67.38\\
4763.00	67.33\\
5066.00	67.45\\
5389.00	67.67\\
5733.00	67.82\\
6099.00	67.63\\
6488.00	67.80\\
6901.00	67.69\\
7341.00	67.88\\
7809.00	67.60\\
8307.00	67.71\\
8837.00	67.49\\
9401.00	67.46\\
10000.00	67.60\\
};
\end{axis}

\begin{axis}[%
width=0.391\figurewidth,
height=0.265\figureheight,
at={(0.54\figurewidth,0\figureheight)},
scale only axis,
xmin=1.00,
xmax=10000.00,
xlabel style={font=\color{mycolor1}},
xlabel={Iterationen $\sym{n}_{\sym{idx_MC}}$},
ymin=0.00,
ymax=105.00,
ytick={0.00,20.00,40.00,60.00,80.00,100.00},
yticklabels={\empty},
axis background/.style={fill=white},
title style={font=\bfseries\color{mycolor1}},
title={\textbf{$\sym{beta}_{22}$}},
xmajorgrids,
ymajorgrids,
grid style={dotted, opacity=0.25}
]
\addplot [color=white!10!black, line width=1.2pt, forget plot]
  table[row sep=crcr]{%
1.00	100.00\\
2.00	100.00\\
3.00	66.67\\
4.00	50.00\\
5.00	60.00\\
6.00	66.67\\
7.00	71.43\\
8.00	75.00\\
9.00	66.67\\
10.00	60.00\\
11.00	63.64\\
12.00	66.67\\
13.00	61.54\\
14.00	64.29\\
15.00	66.67\\
16.00	68.75\\
17.00	70.59\\
18.00	72.22\\
19.00	73.68\\
21.00	71.43\\
23.00	73.91\\
25.00	76.00\\
26.00	76.92\\
28.00	78.57\\
30.00	73.33\\
32.00	75.00\\
34.00	76.47\\
36.00	75.00\\
38.00	73.68\\
41.00	73.17\\
43.00	72.09\\
46.00	73.91\\
49.00	73.47\\
52.00	71.15\\
56.00	69.64\\
59.00	69.49\\
63.00	66.67\\
67.00	68.66\\
71.00	67.61\\
81.00	66.67\\
86.00	66.28\\
91.00	67.03\\
97.00	64.95\\
103.00	62.14\\
110.00	61.82\\
117.00	62.39\\
124.00	62.10\\
132.00	62.88\\
140.00	62.86\\
159.00	64.78\\
169.00	65.09\\
180.00	64.44\\
191.00	64.40\\
217.00	63.13\\
230.00	63.04\\
245.00	64.08\\
261.00	64.75\\
277.00	66.06\\
295.00	66.44\\
314.00	67.52\\
334.00	68.56\\
355.00	68.73\\
378.00	68.52\\
402.00	68.91\\
427.00	68.85\\
455.00	68.57\\
484.00	69.21\\
515.00	68.93\\
547.00	69.10\\
619.00	69.79\\
659.00	70.26\\
701.00	69.76\\
746.00	69.71\\
793.00	69.74\\
955.00	70.99\\
1016.00	71.06\\
1080.00	71.02\\
1149.00	71.11\\
1300.00	71.85\\
1383.00	71.66\\
1472.00	70.92\\
1565.00	70.29\\
1665.00	69.91\\
1771.00	69.96\\
1884.00	69.80\\
2132.00	70.08\\
2268.00	70.06\\
2413.00	69.83\\
2567.00	69.42\\
2730.00	69.67\\
3090.00	69.84\\
3287.00	69.46\\
3496.00	69.36\\
3719.00	69.16\\
4477.00	69.09\\
4763.00	69.16\\
5066.00	68.75\\
5733.00	69.13\\
6901.00	69.42\\
7341.00	69.36\\
7809.00	69.51\\
8307.00	69.57\\
8837.00	69.53\\
10000.00	69.71\\
};
\addplot [color=black!25!darkgray, dashed, line width=1.2pt, forget plot]
  table[row sep=crcr]{%
1.00	100.00\\
3.00	100.00\\
4.00	75.00\\
5.00	80.00\\
6.00	66.67\\
7.00	71.43\\
8.00	75.00\\
9.00	77.78\\
10.00	80.00\\
11.00	81.82\\
12.00	83.33\\
13.00	76.92\\
14.00	78.57\\
15.00	80.00\\
16.00	81.25\\
17.00	76.47\\
18.00	72.22\\
19.00	73.68\\
21.00	76.19\\
22.00	72.73\\
23.00	73.91\\
25.00	72.00\\
26.00	73.08\\
28.00	67.86\\
30.00	66.67\\
32.00	65.62\\
34.00	64.71\\
36.00	63.89\\
38.00	65.79\\
41.00	63.41\\
43.00	65.12\\
46.00	65.22\\
49.00	67.35\\
52.00	69.23\\
56.00	71.43\\
59.00	71.19\\
63.00	69.84\\
67.00	70.15\\
71.00	69.01\\
76.00	68.42\\
81.00	69.14\\
86.00	67.44\\
91.00	68.13\\
97.00	67.01\\
103.00	64.08\\
117.00	64.96\\
124.00	62.90\\
132.00	63.64\\
140.00	64.29\\
149.00	63.76\\
169.00	63.31\\
180.00	65.00\\
191.00	64.40\\
204.00	65.69\\
217.00	64.98\\
230.00	65.22\\
245.00	65.31\\
261.00	64.75\\
277.00	65.70\\
295.00	67.46\\
314.00	66.88\\
334.00	67.66\\
355.00	67.89\\
378.00	68.25\\
402.00	68.16\\
427.00	69.09\\
455.00	68.79\\
484.00	68.80\\
515.00	69.90\\
547.00	70.20\\
582.00	69.59\\
619.00	69.95\\
659.00	70.11\\
701.00	70.76\\
793.00	71.12\\
844.00	71.45\\
897.00	71.24\\
955.00	71.31\\
1016.00	71.06\\
1080.00	71.30\\
1149.00	70.58\\
1222.00	70.70\\
1300.00	70.54\\
1383.00	70.50\\
1472.00	70.72\\
1565.00	70.35\\
1665.00	70.27\\
1771.00	70.30\\
1884.00	70.59\\
2132.00	70.45\\
2268.00	70.72\\
2413.00	70.62\\
2567.00	70.71\\
2730.00	70.15\\
2905.00	70.12\\
3090.00	69.81\\
3287.00	69.64\\
3496.00	69.79\\
3719.00	69.56\\
3957.00	69.47\\
4209.00	69.54\\
4477.00	69.53\\
4763.00	69.62\\
5733.00	69.56\\
6488.00	69.47\\
6901.00	69.61\\
7341.00	69.49\\
7809.00	69.48\\
8307.00	69.21\\
9401.00	69.08\\
10000.00	68.96\\
};
\addplot [color=black!45!gray, dotted, line width=1.2pt, forget plot]
  table[row sep=crcr]{%
1.00	100.00\\
2.00	50.00\\
3.00	33.33\\
4.00	50.00\\
5.00	60.00\\
6.00	66.67\\
7.00	71.43\\
8.00	75.00\\
9.00	66.67\\
10.00	70.00\\
11.00	72.73\\
12.00	75.00\\
13.00	76.92\\
14.00	78.57\\
15.00	80.00\\
16.00	81.25\\
17.00	82.35\\
19.00	84.21\\
21.00	85.71\\
23.00	86.96\\
26.00	88.46\\
28.00	89.29\\
30.00	90.00\\
32.00	87.50\\
34.00	85.29\\
36.00	80.56\\
38.00	81.58\\
41.00	82.93\\
43.00	81.40\\
46.00	80.43\\
49.00	81.63\\
52.00	78.85\\
56.00	78.57\\
59.00	77.97\\
63.00	76.19\\
67.00	77.61\\
71.00	76.06\\
76.00	76.32\\
81.00	74.07\\
86.00	74.42\\
91.00	75.82\\
103.00	76.70\\
117.00	76.07\\
140.00	77.86\\
149.00	75.84\\
159.00	74.21\\
169.00	73.96\\
180.00	74.44\\
191.00	73.30\\
204.00	74.02\\
217.00	73.73\\
230.00	72.61\\
245.00	72.24\\
261.00	71.26\\
277.00	70.76\\
295.00	72.54\\
314.00	72.29\\
334.00	71.26\\
355.00	71.55\\
378.00	71.69\\
427.00	72.60\\
455.00	72.75\\
484.00	72.52\\
515.00	72.62\\
547.00	72.21\\
582.00	72.16\\
619.00	71.08\\
659.00	70.71\\
701.00	71.04\\
746.00	70.91\\
793.00	70.24\\
844.00	70.14\\
897.00	70.23\\
955.00	69.53\\
1016.00	69.69\\
1149.00	70.41\\
1222.00	70.05\\
1300.00	69.85\\
1383.00	69.49\\
1472.00	69.50\\
1665.00	68.83\\
1771.00	68.94\\
1884.00	68.90\\
2132.00	67.82\\
2413.00	67.55\\
2730.00	67.84\\
2905.00	67.68\\
3090.00	67.38\\
3287.00	67.14\\
3496.00	67.25\\
4209.00	67.09\\
4477.00	67.23\\
4763.00	67.56\\
5066.00	67.61\\
5389.00	67.80\\
6099.00	67.91\\
6901.00	68.35\\
7809.00	68.51\\
8307.00	68.46\\
8837.00	68.54\\
10000.00	68.37\\
};
\addplot [color=white!15!darkgray, dashdotted, line width=1.2pt, forget plot]
  table[row sep=crcr]{%
1.00	100.00\\
2.00	100.00\\
3.00	66.67\\
4.00	50.00\\
5.00	40.00\\
6.00	50.00\\
7.00	42.86\\
8.00	50.00\\
9.00	55.56\\
10.00	60.00\\
11.00	63.64\\
12.00	66.67\\
13.00	69.23\\
14.00	71.43\\
15.00	73.33\\
16.00	75.00\\
17.00	76.47\\
18.00	72.22\\
19.00	73.68\\
21.00	71.43\\
23.00	73.91\\
25.00	76.00\\
26.00	76.92\\
28.00	75.00\\
30.00	70.00\\
32.00	68.75\\
34.00	67.65\\
36.00	69.44\\
38.00	71.05\\
41.00	68.29\\
43.00	69.77\\
46.00	67.39\\
49.00	69.39\\
52.00	69.23\\
56.00	69.64\\
59.00	67.80\\
63.00	68.25\\
67.00	70.15\\
71.00	69.01\\
76.00	71.05\\
81.00	70.37\\
86.00	72.09\\
91.00	69.23\\
97.00	68.04\\
103.00	67.96\\
110.00	66.36\\
117.00	64.10\\
124.00	66.13\\
132.00	67.42\\
140.00	69.29\\
149.00	69.80\\
159.00	67.92\\
169.00	68.64\\
180.00	68.33\\
191.00	69.11\\
204.00	68.63\\
217.00	68.66\\
230.00	68.26\\
245.00	68.57\\
261.00	67.82\\
277.00	67.87\\
295.00	68.14\\
334.00	69.46\\
355.00	69.58\\
378.00	70.63\\
402.00	71.14\\
427.00	70.49\\
455.00	70.99\\
484.00	70.66\\
515.00	70.87\\
547.00	69.84\\
582.00	70.27\\
619.00	70.44\\
659.00	70.71\\
701.00	70.61\\
746.00	70.38\\
793.00	70.49\\
844.00	70.38\\
897.00	70.01\\
955.00	70.47\\
1016.00	70.37\\
1149.00	70.06\\
1222.00	69.97\\
1300.00	69.46\\
1472.00	68.82\\
1665.00	69.55\\
1771.00	69.28\\
1884.00	68.74\\
2005.00	68.78\\
2268.00	69.27\\
2413.00	69.33\\
2567.00	69.22\\
4209.00	69.21\\
4477.00	69.15\\
4763.00	69.24\\
5066.00	69.11\\
5389.00	69.20\\
5733.00	68.90\\
6901.00	68.79\\
7809.00	69.02\\
8307.00	68.75\\
9401.00	69.04\\
10000.00	69.03\\
};
\addplot [color=white!45!black, line width=1.2pt, forget plot]
  table[row sep=crcr]{%
1.00	100.00\\
2.00	50.00\\
3.00	33.33\\
4.00	25.00\\
5.00	20.00\\
6.00	33.33\\
7.00	42.86\\
8.00	50.00\\
9.00	44.44\\
10.00	50.00\\
11.00	45.45\\
12.00	50.00\\
13.00	53.85\\
14.00	50.00\\
15.00	53.33\\
16.00	56.25\\
17.00	52.94\\
18.00	50.00\\
19.00	52.63\\
21.00	57.14\\
22.00	54.55\\
23.00	52.17\\
25.00	56.00\\
26.00	57.69\\
28.00	57.14\\
30.00	60.00\\
32.00	62.50\\
34.00	61.76\\
36.00	63.89\\
38.00	65.79\\
41.00	65.85\\
43.00	62.79\\
46.00	65.22\\
52.00	65.38\\
56.00	64.29\\
59.00	62.71\\
63.00	63.49\\
67.00	65.67\\
71.00	67.61\\
81.00	66.67\\
91.00	65.93\\
97.00	68.04\\
103.00	67.96\\
110.00	69.09\\
117.00	69.23\\
124.00	68.55\\
149.00	69.80\\
159.00	68.55\\
169.00	66.86\\
180.00	67.78\\
191.00	67.54\\
204.00	66.67\\
217.00	66.82\\
230.00	66.52\\
245.00	66.94\\
261.00	67.05\\
277.00	66.06\\
295.00	66.44\\
314.00	67.20\\
334.00	66.17\\
355.00	66.20\\
402.00	66.67\\
427.00	66.51\\
484.00	66.74\\
515.00	66.41\\
547.00	66.36\\
582.00	66.15\\
619.00	66.07\\
659.00	66.16\\
701.00	66.76\\
746.00	66.62\\
793.00	66.96\\
844.00	66.59\\
897.00	66.33\\
1016.00	67.32\\
1080.00	67.04\\
1149.00	67.36\\
1222.00	67.51\\
1300.00	67.77\\
1383.00	67.53\\
1472.00	67.73\\
1565.00	67.67\\
1665.00	67.81\\
1771.00	68.27\\
1884.00	67.89\\
2005.00	67.93\\
2132.00	68.06\\
2268.00	68.47\\
2413.00	68.38\\
2567.00	68.52\\
2730.00	68.13\\
2905.00	68.26\\
3287.00	68.42\\
3719.00	68.78\\
3957.00	68.61\\
4209.00	68.71\\
4477.00	68.95\\
4763.00	68.86\\
5066.00	68.93\\
5733.00	68.67\\
6099.00	68.80\\
6901.00	68.63\\
7809.00	68.81\\
8307.00	68.77\\
8837.00	68.61\\
9401.00	68.60\\
10000.00	68.69\\
};
\addplot [color=gray!85!lightgray, dashed, line width=1.2pt, forget plot]
  table[row sep=crcr]{%
1.00	100.00\\
5.00	100.00\\
6.00	83.33\\
7.00	71.43\\
8.00	75.00\\
9.00	77.78\\
10.00	70.00\\
11.00	63.64\\
12.00	58.33\\
13.00	61.54\\
14.00	64.29\\
15.00	60.00\\
16.00	62.50\\
17.00	64.71\\
18.00	61.11\\
19.00	63.16\\
21.00	61.90\\
23.00	65.22\\
25.00	68.00\\
26.00	65.38\\
28.00	64.29\\
30.00	63.33\\
32.00	65.62\\
34.00	64.71\\
36.00	66.67\\
38.00	65.79\\
41.00	68.29\\
43.00	67.44\\
46.00	67.39\\
49.00	65.31\\
52.00	63.46\\
56.00	62.50\\
59.00	62.71\\
63.00	65.08\\
71.00	66.20\\
76.00	65.79\\
81.00	66.67\\
86.00	63.95\\
91.00	64.84\\
97.00	64.95\\
110.00	67.27\\
117.00	65.81\\
124.00	65.32\\
132.00	64.39\\
140.00	63.57\\
149.00	63.76\\
159.00	65.41\\
169.00	64.50\\
180.00	66.11\\
191.00	67.02\\
204.00	66.67\\
217.00	67.28\\
230.00	66.52\\
245.00	66.94\\
261.00	66.67\\
314.00	64.65\\
334.00	65.57\\
355.00	65.92\\
378.00	65.34\\
402.00	66.67\\
427.00	66.04\\
455.00	65.71\\
484.00	67.15\\
515.00	67.77\\
547.00	67.28\\
582.00	66.84\\
619.00	67.21\\
659.00	67.07\\
701.00	67.33\\
746.00	66.89\\
793.00	67.59\\
844.00	67.42\\
897.00	67.56\\
955.00	67.54\\
1016.00	67.13\\
1149.00	67.19\\
1222.00	67.59\\
1300.00	67.62\\
1383.00	67.39\\
1472.00	67.73\\
1565.00	68.63\\
1665.00	68.11\\
1771.00	67.98\\
1884.00	67.68\\
2005.00	67.88\\
2268.00	68.21\\
2413.00	68.13\\
2567.00	67.90\\
2730.00	68.28\\
2905.00	68.12\\
3719.00	67.81\\
3957.00	67.93\\
4209.00	67.85\\
4477.00	67.43\\
5066.00	67.67\\
5733.00	67.22\\
6099.00	67.55\\
6488.00	67.52\\
7341.00	67.69\\
8307.00	67.69\\
10000.00	67.88\\
};
\addplot [color=white!50!darkgray, dotted, line width=1.2pt, forget plot]
  table[row sep=crcr]{%
1.00	100.00\\
2.00	100.00\\
3.00	66.67\\
4.00	75.00\\
5.00	80.00\\
6.00	83.33\\
7.00	85.71\\
8.00	87.50\\
9.00	88.89\\
10.00	90.00\\
11.00	90.91\\
12.00	91.67\\
14.00	92.86\\
16.00	93.75\\
17.00	94.12\\
18.00	88.89\\
19.00	89.47\\
21.00	85.71\\
23.00	86.96\\
26.00	88.46\\
28.00	89.29\\
30.00	90.00\\
32.00	87.50\\
34.00	85.29\\
36.00	86.11\\
38.00	81.58\\
41.00	75.61\\
43.00	74.42\\
46.00	76.09\\
49.00	73.47\\
52.00	73.08\\
56.00	73.21\\
59.00	71.19\\
63.00	73.02\\
67.00	74.63\\
71.00	73.24\\
76.00	71.05\\
86.00	72.09\\
91.00	72.53\\
97.00	72.16\\
103.00	69.90\\
110.00	68.18\\
117.00	68.38\\
132.00	70.45\\
140.00	70.71\\
149.00	70.47\\
159.00	70.44\\
169.00	71.01\\
180.00	72.22\\
191.00	72.25\\
204.00	71.57\\
217.00	71.43\\
230.00	72.61\\
245.00	72.65\\
277.00	71.48\\
295.00	71.86\\
314.00	71.66\\
334.00	71.86\\
355.00	71.27\\
378.00	70.90\\
427.00	70.96\\
455.00	70.77\\
484.00	70.87\\
515.00	71.26\\
547.00	71.12\\
582.00	70.10\\
619.00	70.76\\
659.00	70.71\\
701.00	70.90\\
746.00	71.72\\
793.00	72.38\\
897.00	72.24\\
1016.00	72.24\\
1080.00	72.41\\
1149.00	72.06\\
1222.00	72.42\\
1300.00	72.62\\
1383.00	72.60\\
1472.00	72.15\\
1565.00	72.91\\
1884.00	73.09\\
2005.00	73.32\\
2132.00	73.78\\
2905.00	72.67\\
3090.00	72.78\\
3287.00	72.68\\
3496.00	72.71\\
3719.00	72.30\\
3957.00	72.28\\
4209.00	72.39\\
4477.00	72.35\\
4763.00	72.45\\
5389.00	72.50\\
5733.00	72.41\\
6099.00	72.42\\
6488.00	72.53\\
6901.00	72.32\\
7341.00	72.51\\
7809.00	72.56\\
8837.00	72.47\\
9401.00	72.51\\
10000.00	72.46\\
};
\addplot [color=black!5!lightgray, dashdotted, line width=1.2pt, forget plot]
  table[row sep=crcr]{%
1.00	100.00\\
3.00	100.00\\
4.00	75.00\\
5.00	60.00\\
6.00	66.67\\
7.00	71.43\\
8.00	75.00\\
9.00	77.78\\
10.00	80.00\\
11.00	81.82\\
12.00	83.33\\
13.00	84.62\\
14.00	85.71\\
15.00	80.00\\
16.00	75.00\\
17.00	76.47\\
18.00	77.78\\
19.00	73.68\\
21.00	76.19\\
22.00	72.73\\
23.00	73.91\\
25.00	72.00\\
26.00	69.23\\
28.00	64.29\\
30.00	66.67\\
32.00	68.75\\
34.00	70.59\\
36.00	72.22\\
38.00	73.68\\
41.00	73.17\\
43.00	72.09\\
49.00	71.43\\
52.00	71.15\\
56.00	73.21\\
59.00	72.88\\
63.00	74.60\\
67.00	71.64\\
71.00	71.83\\
76.00	72.37\\
81.00	70.37\\
86.00	70.93\\
91.00	70.33\\
97.00	71.13\\
103.00	68.93\\
110.00	69.09\\
117.00	67.52\\
124.00	66.94\\
132.00	67.42\\
140.00	65.71\\
149.00	66.44\\
159.00	66.04\\
169.00	66.27\\
180.00	65.56\\
191.00	65.97\\
204.00	65.69\\
217.00	66.82\\
230.00	66.52\\
261.00	67.43\\
277.00	67.15\\
295.00	68.14\\
314.00	67.20\\
334.00	67.07\\
355.00	66.76\\
378.00	65.87\\
402.00	65.17\\
427.00	66.28\\
455.00	67.03\\
515.00	67.77\\
547.00	67.82\\
582.00	67.53\\
619.00	68.01\\
659.00	67.83\\
701.00	68.19\\
746.00	68.77\\
793.00	69.23\\
844.00	68.25\\
897.00	67.78\\
955.00	69.01\\
1016.00	69.09\\
1080.00	69.35\\
1149.00	69.71\\
1222.00	69.97\\
1300.00	69.69\\
1383.00	69.56\\
1472.00	69.90\\
1565.00	69.39\\
1665.00	69.07\\
1771.00	68.61\\
2005.00	67.93\\
2132.00	68.34\\
2413.00	67.84\\
2730.00	68.21\\
3090.00	68.16\\
3287.00	68.39\\
3719.00	68.16\\
3957.00	68.01\\
4477.00	67.88\\
4763.00	68.09\\
5066.00	67.88\\
5389.00	67.95\\
5733.00	67.89\\
6099.00	67.73\\
6488.00	67.66\\
6901.00	67.82\\
8837.00	67.81\\
10000.00	67.80\\
};
\addplot [color=white!80!black, line width=1.2pt, forget plot]
  table[row sep=crcr]{%
1.00	100.00\\
2.00	50.00\\
3.00	66.67\\
4.00	75.00\\
5.00	60.00\\
6.00	66.67\\
7.00	71.43\\
8.00	75.00\\
9.00	77.78\\
10.00	70.00\\
11.00	72.73\\
12.00	75.00\\
13.00	69.23\\
14.00	64.29\\
15.00	66.67\\
16.00	68.75\\
17.00	64.71\\
18.00	66.67\\
19.00	68.42\\
21.00	66.67\\
23.00	69.57\\
25.00	72.00\\
26.00	69.23\\
28.00	71.43\\
30.00	66.67\\
32.00	68.75\\
34.00	70.59\\
36.00	69.44\\
38.00	68.42\\
41.00	63.41\\
43.00	65.12\\
46.00	65.22\\
49.00	61.22\\
52.00	61.54\\
56.00	62.50\\
59.00	64.41\\
67.00	65.67\\
71.00	63.38\\
76.00	64.47\\
81.00	66.67\\
86.00	67.44\\
91.00	65.93\\
97.00	65.98\\
103.00	66.99\\
110.00	66.36\\
124.00	66.94\\
132.00	66.67\\
140.00	67.14\\
169.00	69.23\\
180.00	69.44\\
191.00	68.59\\
204.00	70.59\\
217.00	71.43\\
230.00	71.74\\
245.00	71.43\\
261.00	70.50\\
277.00	70.04\\
295.00	70.51\\
314.00	70.38\\
334.00	70.36\\
355.00	70.42\\
378.00	70.90\\
402.00	70.40\\
427.00	70.02\\
455.00	70.11\\
484.00	70.87\\
515.00	70.68\\
547.00	70.75\\
619.00	69.14\\
659.00	68.59\\
701.00	68.47\\
746.00	68.10\\
793.00	67.84\\
844.00	67.89\\
897.00	68.45\\
955.00	68.48\\
1016.00	69.00\\
1149.00	69.89\\
1222.00	69.48\\
1300.00	68.85\\
1383.00	68.62\\
1472.00	68.82\\
1665.00	68.17\\
1771.00	68.94\\
1884.00	68.74\\
2005.00	68.28\\
2132.00	68.34\\
2268.00	68.61\\
2413.00	68.34\\
2567.00	68.64\\
2730.00	69.05\\
2905.00	68.85\\
3090.00	68.90\\
3719.00	69.45\\
3957.00	69.42\\
4209.00	69.64\\
4477.00	69.44\\
4763.00	69.54\\
5066.00	69.54\\
5389.00	69.46\\
5733.00	69.23\\
6099.00	69.09\\
6488.00	69.11\\
6901.00	69.05\\
7341.00	69.13\\
8307.00	69.04\\
8837.00	68.65\\
9401.00	68.66\\
10000.00	68.78\\
};
\end{axis}
\end{tikzpicture}%
    \caption{Laufende Konvergenz der Trennschärfe für alle Modellparameter. Die Schätzungen stabilisieren sich bei der $\sym{n}_{\sym{idx_MC}}=10.000$ reproduzierbar.}
    \label{fig:app.c_conv_power}
\end{figure}

\begin{figure}[htbp]
    \centering
    \setlength\figurewidth{0.5\textwidth}
    \setlength\figureheight{16cm}
    % This file was created by matlab2tikz.
%
\definecolor{mycolor1}{rgb}{0.12941,0.12941,0.12941}%
%
\begin{tikzpicture}

\begin{axis}[%
width=0.411\figurewidth,
height=0.265\figureheight,
at={(0\figurewidth,0.735\figureheight)},
scale only axis,
xmin=0.00,
xmax=10000.00,
xtick={0.00,5000.00,10000.00},
xticklabels={\empty},
ymin=0.00,
ymax=40.00,
ylabel style={font=\color{mycolor1}},
ylabel={$\hat{\sym{thet}}_{\sym{idx_i}}$},
axis background/.style={fill=white},
title style={font=\bfseries\color{mycolor1}},
title={\textbf{$\sym{beta}_0$}},
xmajorgrids,
ymajorgrids,
grid style={dotted, opacity=0.25},
legend style={legend cell align=left, align=left}
]
\addplot [color=white!10!black, line width=1.2pt]
  table[row sep=crcr]{%
1.00	18.88\\
2.00	20.02\\
3.00	21.78\\
4.00	22.65\\
5.00	23.34\\
6.00	23.75\\
7.00	23.91\\
8.00	23.46\\
9.00	23.36\\
10.00	22.78\\
11.00	22.65\\
12.00	22.88\\
13.00	22.66\\
14.00	22.48\\
15.00	22.19\\
16.00	22.07\\
17.00	21.44\\
18.00	21.48\\
19.00	21.95\\
21.00	22.68\\
22.00	22.99\\
23.00	23.42\\
25.00	23.76\\
26.00	24.08\\
28.00	23.90\\
30.00	23.90\\
32.00	23.59\\
34.00	23.81\\
36.00	23.64\\
38.00	23.62\\
41.00	23.55\\
43.00	23.25\\
46.00	23.31\\
49.00	23.28\\
52.00	23.60\\
56.00	23.70\\
59.00	23.66\\
63.00	23.76\\
67.00	23.61\\
71.00	23.67\\
76.00	23.95\\
81.00	24.01\\
86.00	24.03\\
91.00	24.13\\
97.00	24.10\\
103.00	23.90\\
117.00	23.90\\
124.00	23.82\\
132.00	23.84\\
149.00	23.64\\
159.00	23.76\\
169.00	23.77\\
180.00	23.68\\
191.00	23.67\\
204.00	23.58\\
217.00	23.69\\
230.00	23.67\\
245.00	23.71\\
261.00	23.52\\
277.00	23.44\\
295.00	23.59\\
314.00	23.72\\
427.00	23.79\\
455.00	23.74\\
515.00	23.96\\
547.00	24.02\\
619.00	24.04\\
659.00	24.03\\
701.00	24.07\\
746.00	24.06\\
793.00	24.12\\
844.00	24.12\\
897.00	24.02\\
955.00	24.07\\
1080.00	24.08\\
1149.00	24.09\\
1383.00	23.98\\
1665.00	23.91\\
1884.00	23.94\\
2268.00	23.91\\
2413.00	23.96\\
2730.00	23.95\\
3719.00	23.94\\
3957.00	23.97\\
4477.00	23.95\\
6488.00	23.99\\
7341.00	24.03\\
10000.00	24.05\\
};
\addlegendentry{Basis \ac{CCD}}

\addplot [color=white!20!black, dashed, line width=1.2pt]
  table[row sep=crcr]{%
1.00	31.50\\
2.00	30.02\\
4.00	26.62\\
5.00	26.13\\
6.00	25.74\\
7.00	25.08\\
8.00	25.53\\
9.00	25.51\\
10.00	24.80\\
11.00	24.26\\
12.00	23.84\\
13.00	24.18\\
14.00	23.91\\
15.00	24.18\\
16.00	24.22\\
17.00	24.08\\
18.00	23.80\\
19.00	23.39\\
21.00	23.71\\
22.00	23.51\\
23.00	23.25\\
25.00	23.65\\
26.00	23.93\\
28.00	23.76\\
30.00	24.01\\
32.00	24.06\\
34.00	24.24\\
36.00	24.33\\
38.00	24.78\\
41.00	24.79\\
43.00	24.85\\
46.00	24.88\\
49.00	24.73\\
52.00	24.86\\
56.00	24.58\\
59.00	24.70\\
63.00	24.64\\
67.00	24.62\\
71.00	24.63\\
76.00	24.58\\
81.00	24.82\\
86.00	24.91\\
97.00	25.25\\
103.00	25.12\\
110.00	24.79\\
117.00	24.79\\
124.00	24.61\\
132.00	24.56\\
140.00	24.55\\
149.00	24.38\\
159.00	24.42\\
191.00	24.26\\
204.00	24.12\\
217.00	24.09\\
230.00	24.15\\
245.00	24.13\\
261.00	24.30\\
277.00	24.31\\
295.00	24.53\\
314.00	24.56\\
334.00	24.62\\
355.00	24.59\\
378.00	24.39\\
402.00	24.38\\
427.00	24.31\\
455.00	24.11\\
515.00	24.06\\
547.00	24.10\\
582.00	24.06\\
619.00	23.95\\
746.00	23.91\\
793.00	23.97\\
897.00	23.93\\
955.00	24.04\\
1149.00	24.08\\
1222.00	24.06\\
1383.00	24.09\\
1472.00	24.15\\
1565.00	24.13\\
1665.00	24.18\\
1771.00	24.19\\
2005.00	24.13\\
2268.00	24.15\\
2413.00	24.08\\
2730.00	24.08\\
2905.00	24.05\\
3090.00	24.06\\
3287.00	24.00\\
3719.00	24.01\\
4209.00	23.99\\
6488.00	23.97\\
6901.00	24.00\\
10000.00	23.97\\
};
\addlegendentry{Fall I ($-20\,\%$)}

\addplot [color=white!30!black, dotted, line width=1.2pt]
  table[row sep=crcr]{%
1.00	26.52\\
2.00	26.28\\
3.00	25.21\\
4.00	25.53\\
5.00	23.99\\
6.00	23.55\\
7.00	24.41\\
8.00	23.39\\
9.00	23.21\\
10.00	23.93\\
11.00	24.45\\
12.00	24.18\\
14.00	24.01\\
15.00	24.76\\
16.00	24.58\\
17.00	24.82\\
18.00	24.97\\
19.00	24.49\\
21.00	24.89\\
22.00	25.24\\
23.00	25.05\\
25.00	25.05\\
26.00	24.71\\
28.00	24.86\\
30.00	24.77\\
32.00	24.85\\
34.00	24.66\\
36.00	24.51\\
38.00	24.53\\
41.00	24.52\\
43.00	24.61\\
49.00	24.63\\
52.00	24.36\\
59.00	24.09\\
63.00	24.25\\
67.00	24.30\\
76.00	24.29\\
81.00	24.36\\
86.00	24.34\\
91.00	24.42\\
97.00	24.44\\
110.00	24.40\\
124.00	24.22\\
132.00	24.31\\
140.00	24.14\\
159.00	23.96\\
169.00	23.93\\
180.00	23.81\\
191.00	23.80\\
204.00	23.86\\
217.00	23.70\\
277.00	23.93\\
295.00	23.94\\
314.00	24.08\\
334.00	24.10\\
355.00	24.17\\
402.00	24.05\\
427.00	24.10\\
455.00	24.07\\
515.00	24.15\\
547.00	24.17\\
582.00	24.06\\
619.00	24.21\\
659.00	24.17\\
701.00	24.20\\
793.00	24.18\\
897.00	24.26\\
1016.00	24.31\\
1149.00	24.24\\
1222.00	24.24\\
1300.00	24.20\\
1383.00	24.26\\
1665.00	24.22\\
1771.00	24.27\\
1884.00	24.27\\
2005.00	24.32\\
2132.00	24.32\\
2413.00	24.23\\
2730.00	24.23\\
3287.00	24.21\\
3496.00	24.23\\
3957.00	24.20\\
4209.00	24.20\\
4763.00	24.14\\
9401.00	24.00\\
10000.00	24.01\\
};
\addlegendentry{Fall I ($+20\,\%$)}

\addplot [color=white!40!black, dashdotted, line width=1.2pt]
  table[row sep=crcr]{%
1.00	26.55\\
2.00	26.91\\
3.00	27.20\\
4.00	27.92\\
5.00	25.37\\
6.00	25.15\\
7.00	24.42\\
8.00	24.20\\
9.00	25.00\\
10.00	24.25\\
11.00	24.14\\
12.00	23.89\\
13.00	24.53\\
14.00	24.99\\
15.00	25.05\\
16.00	25.06\\
18.00	25.23\\
19.00	25.62\\
21.00	25.99\\
22.00	26.14\\
23.00	26.00\\
25.00	25.37\\
26.00	25.13\\
28.00	25.32\\
30.00	25.19\\
32.00	25.39\\
34.00	25.69\\
36.00	25.40\\
38.00	25.16\\
41.00	25.06\\
43.00	24.70\\
46.00	24.47\\
49.00	24.54\\
52.00	24.51\\
56.00	24.36\\
59.00	24.55\\
63.00	24.64\\
67.00	24.50\\
71.00	24.58\\
76.00	24.85\\
81.00	24.67\\
86.00	24.72\\
91.00	24.99\\
97.00	25.09\\
103.00	24.91\\
110.00	24.97\\
117.00	24.93\\
124.00	24.65\\
132.00	24.73\\
140.00	24.92\\
159.00	24.78\\
169.00	24.80\\
180.00	24.74\\
191.00	24.59\\
204.00	24.69\\
217.00	24.62\\
230.00	24.67\\
245.00	24.81\\
261.00	24.82\\
277.00	24.77\\
295.00	24.89\\
314.00	24.72\\
334.00	24.61\\
355.00	24.57\\
378.00	24.49\\
402.00	24.47\\
427.00	24.38\\
484.00	24.36\\
515.00	24.32\\
547.00	24.35\\
619.00	24.35\\
701.00	24.38\\
746.00	24.37\\
793.00	24.43\\
1080.00	24.23\\
1222.00	24.23\\
1383.00	24.20\\
1472.00	24.24\\
1771.00	24.17\\
2005.00	24.12\\
3496.00	24.07\\
3719.00	24.09\\
5066.00	24.01\\
5389.00	24.03\\
5733.00	24.02\\
6099.00	24.05\\
7809.00	23.99\\
10000.00	23.99\\
};
\addlegendentry{Fall I (\ac{FC})}

\addplot [color=gray, line width=1.2pt]
  table[row sep=crcr]{%
1.00	21.02\\
2.00	23.33\\
3.00	23.76\\
4.00	23.91\\
5.00	24.14\\
6.00	23.84\\
7.00	25.59\\
8.00	26.10\\
9.00	25.49\\
10.00	25.83\\
11.00	25.77\\
12.00	25.59\\
13.00	25.77\\
15.00	25.87\\
16.00	25.28\\
18.00	24.93\\
19.00	24.57\\
22.00	24.20\\
25.00	24.03\\
26.00	24.14\\
28.00	23.79\\
30.00	23.92\\
32.00	24.11\\
34.00	24.06\\
38.00	23.32\\
41.00	23.22\\
43.00	23.46\\
49.00	23.63\\
52.00	23.74\\
56.00	23.61\\
59.00	23.63\\
63.00	23.61\\
67.00	23.65\\
71.00	23.61\\
76.00	23.68\\
81.00	23.64\\
86.00	24.01\\
97.00	23.96\\
103.00	23.82\\
110.00	23.90\\
117.00	23.82\\
132.00	23.91\\
140.00	24.02\\
149.00	24.20\\
159.00	24.05\\
169.00	24.18\\
180.00	24.17\\
191.00	24.31\\
204.00	24.26\\
217.00	24.31\\
230.00	24.27\\
245.00	24.13\\
295.00	24.40\\
314.00	24.43\\
355.00	24.39\\
378.00	24.34\\
402.00	24.39\\
427.00	24.38\\
455.00	24.44\\
484.00	24.41\\
515.00	24.30\\
582.00	24.20\\
746.00	24.16\\
844.00	24.31\\
897.00	24.38\\
955.00	24.29\\
1080.00	24.28\\
1149.00	24.27\\
1300.00	24.15\\
1565.00	24.10\\
1665.00	24.05\\
2005.00	24.05\\
2413.00	24.01\\
2730.00	24.00\\
4763.00	23.99\\
6488.00	24.05\\
8307.00	24.02\\
9401.00	24.03\\
10000.00	24.03\\
};
\addlegendentry{Fall II (leicht)}

\addplot [color=white!60!black, dashed, line width=1.2pt]
  table[row sep=crcr]{%
1.00	22.23\\
2.00	27.87\\
3.00	25.21\\
4.00	26.02\\
5.00	25.53\\
6.00	25.12\\
7.00	24.21\\
8.00	23.95\\
9.00	24.07\\
10.00	24.78\\
11.00	26.10\\
12.00	25.86\\
13.00	25.76\\
14.00	25.19\\
15.00	25.48\\
16.00	25.96\\
17.00	25.74\\
18.00	25.74\\
19.00	25.28\\
21.00	25.59\\
22.00	26.14\\
23.00	26.28\\
25.00	25.92\\
26.00	25.92\\
28.00	25.26\\
30.00	25.22\\
32.00	25.34\\
34.00	25.25\\
36.00	25.25\\
38.00	25.33\\
41.00	24.85\\
43.00	24.86\\
46.00	24.76\\
52.00	24.32\\
56.00	24.52\\
59.00	24.62\\
63.00	24.59\\
67.00	24.70\\
71.00	24.34\\
76.00	24.06\\
81.00	24.30\\
91.00	24.36\\
97.00	24.30\\
103.00	24.42\\
110.00	24.52\\
117.00	24.71\\
124.00	24.65\\
132.00	24.42\\
140.00	24.55\\
159.00	24.57\\
169.00	24.51\\
180.00	24.49\\
191.00	24.37\\
204.00	24.48\\
217.00	24.43\\
230.00	24.26\\
277.00	24.53\\
334.00	24.33\\
355.00	24.18\\
378.00	24.16\\
402.00	24.20\\
427.00	24.11\\
455.00	24.07\\
484.00	24.08\\
515.00	23.95\\
619.00	24.00\\
659.00	23.93\\
701.00	23.82\\
746.00	23.80\\
844.00	23.86\\
897.00	23.81\\
1016.00	23.84\\
1383.00	23.89\\
1472.00	23.96\\
1565.00	23.97\\
1665.00	23.93\\
1771.00	24.00\\
2005.00	23.96\\
2132.00	23.94\\
2413.00	23.95\\
2567.00	23.95\\
2730.00	24.00\\
2905.00	24.02\\
3957.00	23.96\\
4209.00	23.96\\
4763.00	23.89\\
5733.00	23.83\\
6099.00	23.87\\
7341.00	23.87\\
7809.00	23.89\\
8837.00	23.87\\
10000.00	23.91\\
};
\addlegendentry{Fall II (stark)}

\addplot [color=white!70!black, dotted, line width=1.2pt]
  table[row sep=crcr]{%
1.00	26.02\\
2.00	30.32\\
3.00	30.53\\
4.00	29.40\\
5.00	29.08\\
6.00	28.64\\
7.00	29.96\\
8.00	28.95\\
9.00	28.85\\
10.00	27.86\\
11.00	27.35\\
12.00	27.31\\
13.00	26.99\\
14.00	26.53\\
15.00	26.57\\
16.00	26.55\\
17.00	26.31\\
18.00	26.01\\
19.00	25.97\\
21.00	25.38\\
22.00	25.77\\
23.00	25.39\\
25.00	25.73\\
26.00	25.93\\
28.00	25.57\\
30.00	25.29\\
32.00	25.44\\
34.00	25.08\\
36.00	25.21\\
38.00	25.09\\
41.00	25.15\\
43.00	25.15\\
46.00	24.73\\
52.00	24.20\\
56.00	24.36\\
59.00	24.60\\
63.00	24.80\\
67.00	24.81\\
71.00	24.76\\
76.00	24.79\\
86.00	24.75\\
91.00	24.38\\
97.00	24.37\\
103.00	24.57\\
110.00	24.74\\
117.00	24.56\\
124.00	24.66\\
132.00	24.88\\
149.00	24.90\\
159.00	24.68\\
169.00	24.51\\
180.00	24.49\\
191.00	24.26\\
204.00	24.51\\
217.00	24.53\\
230.00	24.44\\
245.00	24.50\\
261.00	24.68\\
277.00	24.56\\
295.00	24.50\\
314.00	24.54\\
334.00	24.51\\
355.00	24.41\\
378.00	24.25\\
427.00	24.31\\
455.00	24.20\\
484.00	24.20\\
547.00	24.06\\
746.00	23.94\\
793.00	23.98\\
1016.00	23.95\\
1080.00	24.02\\
1149.00	24.04\\
1222.00	23.97\\
1565.00	23.96\\
2268.00	23.98\\
2905.00	24.01\\
3496.00	24.01\\
5733.00	24.06\\
7341.00	24.05\\
10000.00	24.07\\
};
\addlegendentry{Fall III}

\addplot [color=white!80!black, dashdotted, line width=1.2pt]
  table[row sep=crcr]{%
1.00	22.56\\
2.00	27.87\\
3.00	26.04\\
4.00	24.80\\
5.00	24.08\\
6.00	24.37\\
7.00	24.09\\
8.00	22.96\\
9.00	22.84\\
10.00	23.42\\
11.00	22.63\\
12.00	22.57\\
13.00	22.38\\
14.00	22.49\\
15.00	22.80\\
16.00	22.49\\
17.00	22.59\\
18.00	22.62\\
19.00	22.97\\
21.00	23.08\\
22.00	22.92\\
23.00	23.00\\
26.00	23.42\\
28.00	23.53\\
30.00	23.76\\
32.00	23.85\\
34.00	23.80\\
36.00	23.61\\
38.00	23.14\\
41.00	23.31\\
43.00	23.39\\
46.00	23.43\\
49.00	23.82\\
52.00	23.97\\
59.00	24.29\\
67.00	23.65\\
71.00	23.79\\
76.00	23.73\\
86.00	24.05\\
91.00	24.02\\
97.00	24.16\\
103.00	24.00\\
110.00	24.14\\
117.00	24.18\\
124.00	24.13\\
132.00	23.97\\
140.00	24.02\\
149.00	24.21\\
159.00	24.06\\
169.00	24.05\\
180.00	23.98\\
191.00	23.87\\
204.00	23.86\\
217.00	23.81\\
230.00	23.66\\
245.00	23.71\\
261.00	23.82\\
295.00	23.79\\
314.00	23.82\\
334.00	23.70\\
402.00	23.59\\
427.00	23.54\\
455.00	23.55\\
484.00	23.64\\
659.00	23.76\\
793.00	23.68\\
844.00	23.76\\
897.00	23.76\\
955.00	23.84\\
1080.00	23.80\\
1149.00	23.84\\
1472.00	23.78\\
1665.00	23.83\\
1771.00	23.79\\
1884.00	23.79\\
2005.00	23.82\\
2132.00	23.81\\
2268.00	23.85\\
2413.00	23.85\\
2567.00	23.91\\
2905.00	23.93\\
3090.00	23.91\\
3287.00	23.93\\
3719.00	23.89\\
4477.00	23.91\\
5066.00	23.88\\
5389.00	23.91\\
9401.00	23.99\\
10000.00	23.99\\
};
\addlegendentry{Fall IV: \ac{SP}}

\addplot [color=white!90!black, line width=1.2pt]
  table[row sep=crcr]{%
1.00	26.42\\
2.00	24.68\\
3.00	24.64\\
4.00	24.02\\
5.00	20.83\\
6.00	20.61\\
7.00	21.47\\
8.00	20.56\\
9.00	20.04\\
10.00	19.89\\
11.00	19.26\\
12.00	19.01\\
13.00	18.96\\
14.00	19.10\\
15.00	19.51\\
16.00	20.43\\
17.00	20.77\\
18.00	20.76\\
19.00	20.98\\
21.00	21.14\\
22.00	21.33\\
23.00	21.86\\
25.00	22.61\\
26.00	22.68\\
28.00	23.04\\
30.00	23.06\\
32.00	23.35\\
34.00	23.47\\
36.00	23.90\\
38.00	23.85\\
41.00	24.14\\
43.00	24.00\\
46.00	24.00\\
49.00	23.85\\
52.00	23.88\\
56.00	24.13\\
59.00	24.08\\
71.00	24.34\\
76.00	24.31\\
81.00	24.32\\
86.00	24.24\\
91.00	24.33\\
97.00	24.39\\
103.00	24.29\\
110.00	24.25\\
117.00	24.08\\
124.00	24.01\\
132.00	24.22\\
140.00	24.19\\
149.00	24.28\\
159.00	24.26\\
169.00	24.09\\
191.00	24.15\\
204.00	24.11\\
217.00	24.20\\
230.00	24.19\\
245.00	24.11\\
277.00	24.16\\
314.00	24.23\\
334.00	24.16\\
355.00	24.23\\
378.00	24.27\\
427.00	24.14\\
455.00	24.13\\
515.00	24.18\\
547.00	24.30\\
619.00	24.32\\
659.00	24.36\\
746.00	24.26\\
793.00	24.20\\
844.00	24.24\\
955.00	24.17\\
1149.00	24.05\\
1222.00	24.03\\
1472.00	24.09\\
1665.00	24.06\\
1771.00	24.05\\
1884.00	24.08\\
2905.00	24.03\\
3719.00	24.02\\
3957.00	24.03\\
4209.00	24.00\\
5066.00	23.99\\
5733.00	23.98\\
6901.00	24.02\\
7809.00	24.01\\
8837.00	24.04\\
10000.00	24.04\\
};
\addlegendentry{Fall IV: \ac{ZP}}

\addplot [color=black, line width=1.2pt]
  table[row sep=crcr]{%
0.00	24.00\\
10000.00	24.00\\
};
\addlegendentry{Wahrer Wert}

\end{axis}

\begin{axis}[%
width=0.411\figurewidth,
height=0.265\figureheight,
at={(0.54\figurewidth,0.735\figureheight)},
scale only axis,
xmin=0.00,
xmax=10000.00,
xtick={0.00,5000.00,10000.00},
xticklabels={\empty},
ymin=0.00,
ymax=40.00,
ytick={0.00,10.00,20.00,30.00,40.00},
yticklabels={\empty},
axis background/.style={fill=white},
title style={font=\bfseries\color{mycolor1}},
title={\textbf{$\sym{beta}_1$}},
xmajorgrids,
ymajorgrids,
grid style={dotted, opacity=0.25}
]
\addplot [color=white!10!black, line width=1.2pt, forget plot]
  table[row sep=crcr]{%
1.00	19.31\\
2.00	15.22\\
3.00	12.65\\
4.00	12.13\\
5.00	12.74\\
7.00	12.72\\
8.00	13.47\\
9.00	12.94\\
10.00	13.61\\
11.00	13.14\\
12.00	13.58\\
13.00	13.29\\
14.00	13.06\\
15.00	13.20\\
16.00	12.77\\
17.00	12.81\\
18.00	12.37\\
19.00	12.30\\
22.00	11.89\\
23.00	11.83\\
25.00	12.06\\
26.00	12.06\\
28.00	12.12\\
30.00	12.09\\
32.00	12.52\\
34.00	12.06\\
36.00	11.52\\
38.00	12.08\\
41.00	12.11\\
43.00	12.42\\
46.00	12.43\\
49.00	12.31\\
52.00	12.45\\
56.00	12.38\\
59.00	12.45\\
63.00	12.37\\
71.00	12.08\\
76.00	12.02\\
81.00	12.00\\
86.00	12.02\\
91.00	11.97\\
103.00	11.66\\
110.00	11.79\\
117.00	11.81\\
124.00	11.99\\
132.00	11.95\\
140.00	12.16\\
159.00	12.32\\
169.00	12.23\\
180.00	12.24\\
191.00	12.20\\
204.00	12.08\\
217.00	12.08\\
230.00	11.99\\
245.00	11.99\\
277.00	11.90\\
314.00	11.95\\
378.00	11.71\\
402.00	11.62\\
427.00	11.63\\
455.00	11.54\\
484.00	11.52\\
582.00	11.72\\
619.00	11.73\\
659.00	11.81\\
701.00	11.84\\
793.00	11.79\\
1016.00	11.92\\
1080.00	11.88\\
1149.00	11.90\\
1222.00	11.85\\
1383.00	11.91\\
1472.00	11.87\\
1665.00	11.88\\
2005.00	11.88\\
2132.00	11.91\\
2268.00	11.85\\
3719.00	11.97\\
3957.00	11.97\\
4209.00	11.93\\
5066.00	12.02\\
5733.00	12.00\\
6099.00	12.03\\
6901.00	12.00\\
10000.00	12.01\\
};
\addplot [color=white!20!black, dashed, line width=1.2pt, forget plot]
  table[row sep=crcr]{%
1.00	12.04\\
2.00	15.44\\
3.00	15.32\\
4.00	13.92\\
5.00	13.96\\
6.00	13.27\\
7.00	13.52\\
8.00	12.60\\
9.00	12.53\\
10.00	12.26\\
11.00	12.48\\
12.00	11.69\\
13.00	12.24\\
14.00	12.52\\
15.00	12.28\\
16.00	12.24\\
17.00	12.27\\
18.00	12.21\\
19.00	12.36\\
21.00	11.62\\
22.00	12.02\\
23.00	12.74\\
25.00	12.73\\
26.00	12.44\\
30.00	12.19\\
32.00	11.97\\
34.00	12.00\\
38.00	11.59\\
41.00	11.25\\
43.00	11.12\\
46.00	11.30\\
49.00	11.13\\
52.00	11.29\\
56.00	11.06\\
59.00	10.96\\
63.00	10.87\\
67.00	11.11\\
71.00	11.30\\
76.00	11.39\\
81.00	11.34\\
86.00	11.41\\
91.00	11.52\\
97.00	11.55\\
103.00	11.68\\
110.00	11.73\\
117.00	11.95\\
124.00	12.11\\
140.00	11.99\\
149.00	12.22\\
159.00	12.16\\
169.00	12.21\\
180.00	12.38\\
204.00	12.50\\
217.00	12.44\\
230.00	12.46\\
245.00	12.53\\
261.00	12.54\\
277.00	12.51\\
314.00	12.35\\
378.00	12.34\\
402.00	12.40\\
427.00	12.37\\
455.00	12.41\\
484.00	12.42\\
515.00	12.36\\
582.00	12.40\\
619.00	12.33\\
659.00	12.32\\
701.00	12.27\\
746.00	12.31\\
793.00	12.29\\
844.00	12.32\\
897.00	12.30\\
955.00	12.34\\
1016.00	12.26\\
1080.00	12.31\\
1149.00	12.22\\
1222.00	12.18\\
1300.00	12.19\\
1472.00	12.09\\
1665.00	12.12\\
1771.00	12.09\\
2268.00	12.12\\
2567.00	12.05\\
2730.00	12.07\\
3090.00	12.06\\
3287.00	12.10\\
3957.00	12.02\\
6099.00	12.03\\
7341.00	12.06\\
9401.00	12.01\\
10000.00	12.03\\
};
\addplot [color=white!30!black, dotted, line width=1.2pt, forget plot]
  table[row sep=crcr]{%
1.00	12.61\\
2.00	15.53\\
3.00	16.54\\
4.00	16.43\\
5.00	15.09\\
6.00	14.43\\
7.00	14.30\\
8.00	13.89\\
9.00	13.80\\
10.00	12.95\\
11.00	13.13\\
12.00	13.15\\
13.00	13.38\\
14.00	13.10\\
15.00	13.89\\
16.00	13.08\\
17.00	12.99\\
19.00	13.56\\
21.00	13.93\\
22.00	13.94\\
23.00	13.74\\
25.00	13.56\\
26.00	13.62\\
28.00	13.71\\
30.00	13.42\\
32.00	13.44\\
34.00	13.37\\
36.00	13.35\\
41.00	13.62\\
43.00	13.58\\
46.00	13.46\\
49.00	13.20\\
52.00	13.49\\
56.00	13.64\\
63.00	13.74\\
67.00	13.72\\
71.00	13.45\\
76.00	13.34\\
81.00	13.32\\
91.00	13.11\\
110.00	12.47\\
117.00	12.43\\
124.00	12.30\\
132.00	12.08\\
140.00	12.12\\
149.00	12.40\\
159.00	12.31\\
169.00	12.16\\
180.00	12.08\\
204.00	12.12\\
217.00	12.22\\
230.00	12.16\\
245.00	12.25\\
261.00	12.20\\
277.00	12.28\\
295.00	12.41\\
314.00	12.47\\
355.00	12.40\\
378.00	12.46\\
547.00	12.16\\
582.00	12.15\\
659.00	12.22\\
746.00	12.16\\
793.00	12.21\\
844.00	12.14\\
1080.00	12.02\\
1222.00	12.09\\
1565.00	11.97\\
1665.00	11.99\\
1884.00	11.97\\
2005.00	11.99\\
2132.00	11.94\\
2413.00	12.02\\
2730.00	12.04\\
2905.00	12.00\\
3957.00	11.95\\
4209.00	11.92\\
5066.00	11.93\\
5733.00	11.96\\
10000.00	11.99\\
};
\addplot [color=white!40!black, dashdotted, line width=1.2pt, forget plot]
  table[row sep=crcr]{%
1.00	17.12\\
2.00	12.60\\
3.00	12.73\\
4.00	11.74\\
5.00	12.20\\
6.00	10.63\\
7.00	10.35\\
8.00	11.55\\
9.00	11.86\\
10.00	11.83\\
11.00	11.52\\
12.00	12.40\\
13.00	12.99\\
14.00	13.08\\
15.00	13.02\\
16.00	13.22\\
17.00	12.81\\
18.00	12.75\\
19.00	12.76\\
21.00	13.14\\
22.00	13.19\\
23.00	13.55\\
25.00	13.44\\
26.00	13.14\\
28.00	13.13\\
30.00	12.78\\
32.00	12.61\\
34.00	12.21\\
36.00	12.41\\
38.00	12.09\\
43.00	11.80\\
46.00	11.98\\
49.00	12.10\\
52.00	11.98\\
56.00	12.02\\
59.00	11.86\\
63.00	12.04\\
76.00	12.35\\
81.00	12.29\\
86.00	12.35\\
91.00	12.16\\
97.00	12.57\\
110.00	12.47\\
117.00	12.64\\
132.00	12.86\\
140.00	12.80\\
149.00	12.89\\
159.00	12.81\\
169.00	12.61\\
180.00	12.50\\
191.00	12.68\\
204.00	12.69\\
217.00	12.60\\
230.00	12.54\\
314.00	12.68\\
334.00	12.62\\
378.00	12.44\\
402.00	12.48\\
455.00	12.39\\
515.00	12.44\\
582.00	12.25\\
619.00	12.15\\
659.00	12.26\\
701.00	12.28\\
746.00	12.36\\
793.00	12.37\\
844.00	12.34\\
955.00	12.37\\
1080.00	12.22\\
1222.00	12.20\\
1300.00	12.25\\
1472.00	12.22\\
1665.00	12.28\\
1771.00	12.27\\
1884.00	12.21\\
2268.00	12.22\\
2413.00	12.18\\
2730.00	12.21\\
3090.00	12.11\\
4209.00	12.11\\
4477.00	12.13\\
6488.00	12.01\\
10000.00	12.01\\
};
\addplot [color=gray, line width=1.2pt, forget plot]
  table[row sep=crcr]{%
1.00	2.08\\
2.00	5.66\\
3.00	9.38\\
4.00	10.44\\
5.00	9.62\\
6.00	9.83\\
7.00	9.02\\
8.00	9.26\\
9.00	9.58\\
10.00	9.43\\
11.00	9.68\\
12.00	11.34\\
13.00	11.71\\
14.00	12.86\\
15.00	12.80\\
16.00	12.04\\
17.00	12.39\\
18.00	12.59\\
19.00	12.27\\
21.00	12.76\\
23.00	12.77\\
25.00	12.28\\
26.00	12.29\\
28.00	12.21\\
30.00	11.95\\
32.00	11.85\\
34.00	11.78\\
36.00	11.97\\
38.00	11.79\\
41.00	11.84\\
43.00	11.68\\
49.00	11.05\\
52.00	11.11\\
63.00	11.60\\
67.00	11.55\\
76.00	11.17\\
91.00	11.30\\
97.00	11.20\\
110.00	11.19\\
117.00	11.23\\
124.00	11.43\\
132.00	11.44\\
140.00	11.65\\
149.00	11.61\\
159.00	11.64\\
169.00	11.74\\
180.00	11.82\\
191.00	11.80\\
204.00	11.62\\
230.00	11.67\\
261.00	11.50\\
277.00	11.50\\
295.00	11.58\\
314.00	11.55\\
378.00	11.56\\
402.00	11.63\\
455.00	11.62\\
547.00	11.69\\
582.00	11.67\\
619.00	11.61\\
659.00	11.71\\
701.00	11.77\\
746.00	11.63\\
897.00	11.68\\
1016.00	11.69\\
1222.00	11.77\\
1472.00	11.72\\
1665.00	11.77\\
1771.00	11.79\\
1884.00	11.85\\
2005.00	11.85\\
2268.00	11.90\\
2567.00	11.89\\
2730.00	11.95\\
3090.00	11.93\\
3287.00	11.96\\
3496.00	11.93\\
10000.00	12.02\\
};
\addplot [color=white!60!black, dashed, line width=1.2pt, forget plot]
  table[row sep=crcr]{%
1.00	10.63\\
2.00	14.00\\
3.00	14.75\\
4.00	13.91\\
5.00	13.58\\
6.00	12.32\\
7.00	12.78\\
8.00	14.49\\
9.00	15.15\\
11.00	13.90\\
12.00	13.52\\
13.00	13.27\\
14.00	12.54\\
15.00	12.98\\
16.00	13.14\\
17.00	12.82\\
18.00	12.68\\
19.00	12.67\\
21.00	12.44\\
23.00	12.72\\
25.00	12.50\\
26.00	12.23\\
28.00	12.26\\
30.00	11.94\\
32.00	11.94\\
34.00	11.67\\
36.00	11.72\\
38.00	11.88\\
41.00	11.77\\
43.00	11.57\\
46.00	11.83\\
49.00	11.76\\
52.00	11.78\\
56.00	11.52\\
67.00	11.28\\
71.00	11.53\\
76.00	12.07\\
81.00	11.97\\
86.00	11.98\\
91.00	12.04\\
97.00	12.34\\
103.00	12.57\\
110.00	12.29\\
124.00	12.26\\
132.00	12.07\\
149.00	11.98\\
159.00	11.81\\
169.00	12.03\\
180.00	12.18\\
191.00	12.19\\
204.00	12.24\\
217.00	12.34\\
277.00	12.44\\
295.00	12.38\\
314.00	12.38\\
355.00	12.24\\
378.00	12.24\\
547.00	11.91\\
582.00	11.88\\
619.00	11.90\\
746.00	11.85\\
793.00	11.90\\
844.00	11.90\\
897.00	11.94\\
955.00	11.96\\
1016.00	11.93\\
1080.00	11.87\\
1383.00	11.98\\
2005.00	11.93\\
2132.00	11.95\\
2268.00	11.93\\
2413.00	11.98\\
2567.00	11.97\\
3090.00	12.03\\
3287.00	12.04\\
3496.00	12.00\\
3719.00	12.01\\
4209.00	11.99\\
5066.00	11.94\\
5733.00	11.99\\
10000.00	11.94\\
};
\addplot [color=white!70!black, dotted, line width=1.2pt, forget plot]
  table[row sep=crcr]{%
1.00	17.26\\
2.00	15.25\\
3.00	12.78\\
4.00	10.74\\
5.00	9.76\\
6.00	10.62\\
7.00	10.28\\
8.00	11.36\\
9.00	11.64\\
10.00	12.24\\
11.00	12.32\\
12.00	12.58\\
13.00	12.65\\
14.00	12.04\\
15.00	12.02\\
16.00	12.07\\
17.00	12.42\\
18.00	11.97\\
19.00	11.44\\
21.00	11.48\\
23.00	11.99\\
25.00	11.92\\
26.00	11.60\\
28.00	11.18\\
30.00	10.85\\
32.00	10.79\\
36.00	11.17\\
38.00	11.40\\
41.00	10.93\\
43.00	10.97\\
46.00	10.98\\
49.00	10.84\\
52.00	11.01\\
56.00	11.20\\
63.00	11.16\\
71.00	10.92\\
81.00	11.43\\
86.00	11.44\\
91.00	11.65\\
97.00	11.52\\
103.00	11.66\\
110.00	11.62\\
117.00	11.64\\
132.00	11.57\\
140.00	11.62\\
149.00	11.93\\
159.00	11.82\\
169.00	11.85\\
180.00	11.73\\
217.00	11.92\\
230.00	12.06\\
245.00	12.08\\
261.00	12.05\\
277.00	12.15\\
314.00	12.08\\
334.00	12.09\\
355.00	12.07\\
378.00	12.19\\
402.00	12.25\\
427.00	12.22\\
455.00	12.24\\
484.00	12.16\\
515.00	12.15\\
547.00	12.19\\
619.00	12.17\\
701.00	12.14\\
844.00	12.12\\
897.00	12.00\\
955.00	12.05\\
1016.00	12.06\\
1080.00	12.01\\
1149.00	12.00\\
1222.00	12.04\\
1300.00	11.99\\
1472.00	12.06\\
1665.00	11.93\\
1884.00	11.94\\
2268.00	11.92\\
2413.00	11.88\\
2730.00	11.97\\
3719.00	11.96\\
4209.00	11.96\\
5733.00	11.92\\
6901.00	11.88\\
7809.00	11.87\\
10000.00	11.92\\
};
\addplot [color=white!80!black, dashdotted, line width=1.2pt, forget plot]
  table[row sep=crcr]{%
1.00	12.55\\
2.00	8.84\\
3.00	11.10\\
4.00	11.89\\
5.00	13.61\\
6.00	12.36\\
7.00	12.85\\
8.00	13.10\\
9.00	13.42\\
10.00	13.46\\
11.00	12.54\\
12.00	12.47\\
13.00	11.65\\
14.00	11.50\\
15.00	10.81\\
16.00	11.50\\
17.00	11.31\\
18.00	11.55\\
19.00	11.49\\
21.00	11.47\\
22.00	11.68\\
23.00	11.51\\
25.00	11.83\\
26.00	11.78\\
28.00	11.79\\
30.00	11.59\\
32.00	11.65\\
34.00	11.31\\
36.00	11.85\\
38.00	11.68\\
41.00	11.60\\
43.00	11.67\\
46.00	11.82\\
49.00	11.83\\
52.00	11.68\\
56.00	11.78\\
59.00	11.74\\
63.00	11.93\\
71.00	12.11\\
76.00	12.37\\
81.00	12.07\\
86.00	12.14\\
91.00	12.27\\
97.00	12.30\\
103.00	12.23\\
110.00	12.30\\
117.00	12.26\\
124.00	12.39\\
132.00	12.26\\
159.00	12.29\\
169.00	12.11\\
180.00	12.13\\
191.00	12.05\\
204.00	12.11\\
217.00	12.09\\
245.00	11.93\\
261.00	11.89\\
277.00	12.01\\
295.00	12.02\\
314.00	11.96\\
355.00	11.95\\
378.00	12.16\\
402.00	12.00\\
427.00	12.10\\
455.00	12.11\\
582.00	11.98\\
619.00	12.00\\
701.00	11.97\\
746.00	11.86\\
844.00	11.75\\
897.00	11.88\\
955.00	11.85\\
1080.00	11.83\\
1300.00	11.96\\
1472.00	11.93\\
1665.00	11.98\\
1771.00	11.96\\
1884.00	12.00\\
3957.00	11.96\\
4477.00	11.99\\
5066.00	12.02\\
5733.00	11.97\\
6488.00	11.95\\
7341.00	11.95\\
8307.00	11.94\\
10000.00	11.96\\
};
\addplot [color=white!90!black, line width=1.2pt, forget plot]
  table[row sep=crcr]{%
1.00	7.72\\
2.00	9.15\\
3.00	9.95\\
5.00	10.54\\
7.00	12.75\\
8.00	11.67\\
9.00	11.51\\
10.00	11.56\\
11.00	10.53\\
12.00	10.91\\
13.00	11.00\\
14.00	10.47\\
15.00	10.39\\
16.00	10.64\\
17.00	10.72\\
18.00	11.14\\
19.00	11.28\\
21.00	11.33\\
22.00	11.56\\
23.00	11.72\\
25.00	11.93\\
26.00	11.76\\
28.00	11.79\\
30.00	12.20\\
36.00	11.85\\
38.00	12.08\\
41.00	11.91\\
43.00	11.53\\
46.00	11.26\\
49.00	11.22\\
52.00	11.05\\
56.00	11.07\\
59.00	11.24\\
63.00	11.00\\
67.00	10.88\\
71.00	10.73\\
76.00	10.88\\
81.00	11.10\\
86.00	11.07\\
91.00	11.34\\
97.00	11.30\\
117.00	11.35\\
124.00	11.22\\
132.00	11.33\\
140.00	11.28\\
149.00	11.27\\
159.00	11.46\\
169.00	11.56\\
180.00	11.61\\
191.00	11.76\\
204.00	11.82\\
217.00	11.95\\
230.00	12.01\\
245.00	11.92\\
261.00	11.86\\
277.00	11.75\\
334.00	11.86\\
355.00	11.66\\
378.00	11.67\\
402.00	11.56\\
455.00	11.63\\
547.00	11.85\\
582.00	11.94\\
619.00	11.92\\
701.00	12.00\\
746.00	11.92\\
897.00	11.89\\
955.00	11.93\\
1016.00	11.86\\
1080.00	11.94\\
1149.00	11.94\\
1222.00	12.00\\
1300.00	12.03\\
1665.00	11.88\\
1771.00	11.89\\
1884.00	11.87\\
3090.00	11.97\\
3496.00	12.01\\
3719.00	12.02\\
3957.00	11.99\\
4477.00	11.98\\
5389.00	11.93\\
5733.00	11.97\\
6099.00	11.91\\
8307.00	11.90\\
8837.00	11.91\\
9401.00	11.89\\
10000.00	11.92\\
};
\addplot [color=black, line width=1.2pt, forget plot]
  table[row sep=crcr]{%
0.00	12.00\\
10000.00	12.00\\
};
\end{axis}

\begin{axis}[%
width=0.411\figurewidth,
height=0.265\figureheight,
at={(0\figurewidth,0.368\figureheight)},
scale only axis,
xmin=0.00,
xmax=10000.00,
xtick={0.00,5000.00,10000.00},
xticklabels={\empty},
ymin=0.00,
ymax=40.00,
ylabel style={font=\color{mycolor1}},
ylabel={$\hat{\sym{thet}}_{\sym{idx_i}}$},
axis background/.style={fill=white},
title style={font=\bfseries\color{mycolor1}},
title={\textbf{$\sym{beta}_2$}},
xmajorgrids,
ymajorgrids,
grid style={dotted, opacity=0.25}
]
\addplot [color=white!10!black, line width=1.2pt, forget plot]
  table[row sep=crcr]{%
1.00	16.91\\
2.00	17.51\\
3.00	14.96\\
4.00	16.07\\
5.00	15.16\\
6.00	15.51\\
7.00	14.89\\
8.00	14.17\\
9.00	14.25\\
10.00	14.16\\
11.00	15.16\\
12.00	16.03\\
13.00	15.89\\
14.00	15.69\\
15.00	15.76\\
16.00	16.16\\
17.00	16.41\\
18.00	16.27\\
19.00	16.43\\
21.00	15.97\\
22.00	16.15\\
23.00	16.08\\
25.00	16.08\\
26.00	16.17\\
28.00	15.75\\
30.00	15.21\\
32.00	14.74\\
34.00	14.64\\
36.00	14.59\\
38.00	14.65\\
41.00	13.85\\
43.00	13.83\\
46.00	13.52\\
49.00	13.31\\
52.00	13.22\\
56.00	13.23\\
59.00	12.81\\
63.00	12.75\\
67.00	12.92\\
71.00	12.95\\
76.00	12.73\\
81.00	12.47\\
86.00	12.27\\
91.00	12.31\\
97.00	12.62\\
103.00	12.61\\
110.00	12.52\\
117.00	12.67\\
124.00	12.59\\
132.00	12.72\\
140.00	12.69\\
149.00	12.54\\
159.00	12.55\\
169.00	12.37\\
180.00	12.26\\
204.00	12.17\\
217.00	12.03\\
230.00	12.02\\
261.00	12.17\\
277.00	12.10\\
295.00	12.13\\
314.00	12.23\\
355.00	12.24\\
378.00	12.25\\
402.00	12.21\\
427.00	12.23\\
484.00	12.12\\
515.00	12.19\\
547.00	12.21\\
619.00	12.07\\
659.00	12.12\\
701.00	12.03\\
746.00	12.02\\
793.00	11.96\\
844.00	11.93\\
1080.00	12.10\\
1149.00	12.05\\
1383.00	12.06\\
1565.00	12.07\\
1665.00	12.05\\
1771.00	12.08\\
2005.00	12.05\\
2413.00	12.06\\
2567.00	12.10\\
2905.00	12.06\\
3287.00	12.11\\
3496.00	12.12\\
3719.00	12.09\\
4209.00	12.11\\
4477.00	12.08\\
4763.00	12.10\\
5066.00	12.06\\
5733.00	12.06\\
6099.00	12.06\\
6901.00	12.08\\
10000.00	12.10\\
};
\addplot [color=white!20!black, dashed, line width=1.2pt, forget plot]
  table[row sep=crcr]{%
1.00	5.02\\
2.00	11.14\\
3.00	12.22\\
4.00	13.19\\
5.00	12.84\\
6.00	13.78\\
7.00	13.37\\
8.00	12.92\\
10.00	14.60\\
11.00	14.64\\
12.00	14.54\\
13.00	14.38\\
14.00	13.83\\
15.00	14.00\\
16.00	13.57\\
18.00	12.96\\
19.00	12.80\\
21.00	12.87\\
22.00	13.02\\
23.00	12.98\\
25.00	12.35\\
26.00	12.25\\
28.00	12.34\\
30.00	12.58\\
32.00	12.27\\
36.00	12.54\\
38.00	12.76\\
41.00	12.79\\
43.00	12.61\\
46.00	12.88\\
49.00	12.85\\
52.00	12.93\\
56.00	12.66\\
59.00	12.57\\
63.00	12.30\\
67.00	12.36\\
71.00	12.09\\
76.00	12.17\\
81.00	12.22\\
86.00	12.24\\
97.00	12.38\\
103.00	12.30\\
110.00	12.42\\
117.00	12.48\\
124.00	12.28\\
132.00	12.15\\
140.00	11.99\\
149.00	11.94\\
169.00	11.96\\
180.00	12.15\\
191.00	12.12\\
204.00	12.29\\
217.00	12.35\\
261.00	11.99\\
277.00	12.05\\
295.00	11.97\\
314.00	11.97\\
334.00	11.94\\
378.00	12.04\\
427.00	12.07\\
484.00	12.19\\
547.00	12.16\\
582.00	12.13\\
659.00	12.13\\
746.00	11.99\\
793.00	12.13\\
897.00	12.25\\
1149.00	12.06\\
1222.00	11.98\\
1383.00	12.02\\
1565.00	12.02\\
1665.00	12.05\\
1771.00	12.00\\
1884.00	12.01\\
2413.00	12.18\\
2730.00	12.15\\
3719.00	11.99\\
4209.00	11.99\\
5733.00	11.96\\
10000.00	12.00\\
};
\addplot [color=white!30!black, dotted, line width=1.2pt, forget plot]
  table[row sep=crcr]{%
1.00	13.88\\
2.00	14.14\\
3.00	13.97\\
5.00	12.91\\
6.00	12.63\\
7.00	11.49\\
8.00	12.02\\
9.00	11.65\\
10.00	11.60\\
11.00	12.30\\
12.00	11.94\\
13.00	11.80\\
14.00	11.77\\
15.00	12.14\\
17.00	11.63\\
18.00	11.87\\
19.00	11.52\\
21.00	11.67\\
22.00	11.56\\
23.00	11.77\\
25.00	10.99\\
26.00	10.99\\
28.00	10.55\\
30.00	10.64\\
32.00	11.43\\
34.00	11.73\\
36.00	11.82\\
38.00	11.70\\
41.00	11.95\\
43.00	11.92\\
46.00	11.79\\
49.00	11.84\\
52.00	11.61\\
56.00	11.59\\
59.00	11.65\\
63.00	11.51\\
71.00	12.02\\
76.00	11.94\\
86.00	12.05\\
97.00	11.76\\
103.00	11.66\\
110.00	11.73\\
117.00	11.64\\
124.00	11.62\\
140.00	12.05\\
159.00	12.10\\
169.00	12.11\\
180.00	11.85\\
191.00	11.90\\
204.00	11.81\\
217.00	11.74\\
230.00	11.74\\
245.00	11.81\\
261.00	11.84\\
277.00	11.80\\
295.00	11.66\\
314.00	11.56\\
355.00	11.57\\
378.00	11.69\\
427.00	11.77\\
484.00	11.72\\
515.00	11.74\\
547.00	11.69\\
619.00	11.72\\
659.00	11.74\\
746.00	11.69\\
793.00	11.76\\
844.00	11.72\\
1080.00	11.78\\
1149.00	11.86\\
1300.00	11.93\\
1472.00	11.89\\
1771.00	11.86\\
1884.00	11.93\\
2132.00	11.98\\
2268.00	11.95\\
2413.00	11.98\\
2567.00	11.96\\
2730.00	12.00\\
2905.00	12.00\\
3090.00	12.04\\
3287.00	12.02\\
3496.00	12.05\\
3719.00	12.04\\
3957.00	12.09\\
4477.00	12.06\\
5066.00	12.05\\
5389.00	12.06\\
5733.00	12.02\\
6488.00	12.04\\
8837.00	12.02\\
10000.00	12.02\\
};
\addplot [color=white!40!black, dashdotted, line width=1.2pt, forget plot]
  table[row sep=crcr]{%
1.00	10.38\\
3.00	12.97\\
4.00	11.52\\
5.00	13.13\\
6.00	13.03\\
7.00	12.67\\
8.00	12.17\\
9.00	11.85\\
10.00	11.26\\
11.00	11.70\\
12.00	11.97\\
13.00	11.56\\
15.00	10.96\\
16.00	10.54\\
17.00	11.11\\
18.00	11.85\\
19.00	11.57\\
21.00	11.86\\
22.00	11.76\\
23.00	11.95\\
25.00	11.90\\
26.00	11.62\\
28.00	12.17\\
30.00	12.30\\
32.00	12.03\\
34.00	11.71\\
36.00	11.69\\
38.00	11.82\\
41.00	11.71\\
43.00	11.47\\
46.00	11.40\\
52.00	11.44\\
56.00	11.35\\
59.00	11.20\\
63.00	11.55\\
71.00	11.20\\
76.00	11.28\\
81.00	11.19\\
86.00	11.36\\
91.00	11.48\\
97.00	11.52\\
103.00	11.44\\
110.00	11.56\\
117.00	11.48\\
124.00	11.61\\
132.00	11.63\\
140.00	11.53\\
149.00	11.66\\
169.00	11.77\\
180.00	11.76\\
191.00	11.70\\
204.00	11.68\\
217.00	11.74\\
230.00	11.76\\
245.00	11.83\\
261.00	11.87\\
277.00	11.82\\
295.00	11.87\\
314.00	11.82\\
334.00	11.71\\
378.00	11.78\\
402.00	11.87\\
455.00	11.96\\
484.00	11.91\\
547.00	11.90\\
619.00	11.89\\
659.00	11.97\\
701.00	11.98\\
746.00	12.04\\
793.00	11.97\\
844.00	11.95\\
955.00	12.03\\
1016.00	12.02\\
1149.00	12.05\\
1383.00	11.98\\
1565.00	11.89\\
1771.00	11.87\\
2132.00	11.96\\
2268.00	12.02\\
2413.00	12.00\\
2567.00	12.05\\
2730.00	12.02\\
3090.00	12.03\\
3287.00	12.04\\
3496.00	12.00\\
4209.00	12.04\\
5066.00	12.11\\
10000.00	12.05\\
};
\addplot [color=gray, line width=1.2pt, forget plot]
  table[row sep=crcr]{%
1.00	10.38\\
2.00	9.62\\
3.00	12.22\\
4.00	12.93\\
5.00	14.40\\
6.00	15.00\\
7.00	13.16\\
8.00	12.61\\
9.00	11.44\\
10.00	10.80\\
11.00	11.08\\
12.00	10.60\\
13.00	11.36\\
14.00	11.43\\
15.00	11.38\\
16.00	11.84\\
17.00	12.02\\
18.00	12.14\\
19.00	11.96\\
21.00	11.49\\
23.00	11.67\\
25.00	11.95\\
26.00	11.81\\
28.00	11.96\\
30.00	11.84\\
32.00	11.81\\
34.00	11.88\\
36.00	11.85\\
38.00	12.19\\
41.00	12.50\\
43.00	12.38\\
46.00	12.43\\
49.00	12.74\\
56.00	12.75\\
59.00	12.68\\
67.00	12.41\\
71.00	12.70\\
76.00	12.38\\
81.00	12.40\\
86.00	12.37\\
91.00	12.26\\
103.00	12.49\\
124.00	12.42\\
132.00	12.50\\
140.00	12.23\\
149.00	12.13\\
159.00	12.19\\
169.00	11.95\\
180.00	11.92\\
191.00	12.00\\
204.00	12.01\\
217.00	12.06\\
230.00	12.15\\
245.00	12.29\\
261.00	12.28\\
277.00	12.37\\
295.00	12.24\\
314.00	12.17\\
334.00	12.20\\
355.00	12.16\\
378.00	12.03\\
402.00	12.08\\
427.00	12.01\\
455.00	12.06\\
515.00	12.05\\
547.00	11.98\\
659.00	11.94\\
746.00	11.75\\
844.00	11.69\\
955.00	11.70\\
1016.00	11.65\\
1080.00	11.73\\
1149.00	11.78\\
1222.00	11.78\\
1300.00	11.73\\
1383.00	11.80\\
1771.00	11.89\\
2005.00	11.90\\
2132.00	11.85\\
2268.00	11.86\\
2413.00	11.82\\
3287.00	11.84\\
4209.00	11.92\\
4477.00	11.90\\
5733.00	11.95\\
8307.00	11.93\\
10000.00	11.97\\
};
\addplot [color=white!60!black, dashed, line width=1.2pt, forget plot]
  table[row sep=crcr]{%
1.00	8.85\\
2.00	13.10\\
3.00	14.40\\
4.00	13.36\\
5.00	12.92\\
6.00	12.22\\
7.00	10.98\\
8.00	10.98\\
9.00	11.06\\
10.00	11.53\\
11.00	11.61\\
12.00	11.21\\
13.00	12.06\\
14.00	12.53\\
15.00	12.81\\
16.00	12.18\\
17.00	12.16\\
18.00	11.76\\
19.00	11.66\\
21.00	11.97\\
22.00	11.81\\
23.00	12.06\\
25.00	11.87\\
26.00	11.81\\
28.00	11.29\\
30.00	11.41\\
34.00	11.20\\
36.00	10.93\\
41.00	10.96\\
43.00	10.85\\
46.00	10.91\\
52.00	10.82\\
56.00	11.06\\
59.00	11.19\\
63.00	11.00\\
67.00	11.05\\
71.00	10.94\\
76.00	10.98\\
81.00	10.81\\
86.00	11.09\\
97.00	11.48\\
103.00	11.42\\
110.00	11.42\\
124.00	11.56\\
140.00	11.83\\
149.00	11.79\\
159.00	11.97\\
169.00	12.04\\
180.00	12.21\\
204.00	12.25\\
230.00	12.19\\
245.00	12.09\\
261.00	12.11\\
314.00	11.91\\
334.00	11.98\\
355.00	12.00\\
378.00	11.86\\
427.00	11.73\\
484.00	11.87\\
515.00	12.02\\
547.00	11.98\\
619.00	12.07\\
659.00	11.99\\
793.00	11.92\\
844.00	11.94\\
1016.00	11.90\\
1080.00	11.94\\
1149.00	11.92\\
1222.00	11.98\\
1383.00	11.98\\
1472.00	11.92\\
2005.00	12.06\\
2132.00	12.01\\
2268.00	12.01\\
2413.00	12.06\\
3090.00	12.10\\
3287.00	12.05\\
3957.00	12.04\\
4209.00	11.98\\
5066.00	12.02\\
5389.00	12.04\\
6099.00	12.03\\
6488.00	12.03\\
7341.00	12.06\\
10000.00	12.00\\
};
\addplot [color=white!70!black, dotted, line width=1.2pt, forget plot]
  table[row sep=crcr]{%
1.00	1.51\\
2.00	3.62\\
3.00	5.34\\
4.00	8.05\\
5.00	8.03\\
6.00	9.73\\
7.00	10.98\\
8.00	10.58\\
9.00	9.97\\
10.00	10.63\\
11.00	11.15\\
12.00	11.76\\
13.00	11.78\\
14.00	12.50\\
15.00	12.49\\
16.00	12.66\\
17.00	12.55\\
18.00	12.36\\
19.00	12.38\\
21.00	12.16\\
22.00	11.97\\
23.00	12.05\\
25.00	12.31\\
26.00	12.41\\
28.00	12.39\\
30.00	11.82\\
32.00	11.65\\
34.00	11.58\\
38.00	12.27\\
41.00	12.28\\
43.00	12.02\\
46.00	12.05\\
49.00	12.18\\
56.00	12.02\\
59.00	11.83\\
63.00	12.22\\
67.00	12.22\\
71.00	12.08\\
76.00	12.02\\
81.00	11.84\\
86.00	11.98\\
91.00	11.96\\
97.00	11.75\\
103.00	11.66\\
110.00	11.70\\
117.00	11.65\\
124.00	11.78\\
132.00	11.68\\
140.00	11.72\\
149.00	11.66\\
180.00	11.78\\
191.00	11.66\\
204.00	11.65\\
217.00	11.57\\
230.00	11.55\\
245.00	11.57\\
261.00	11.68\\
277.00	11.66\\
295.00	11.73\\
314.00	11.71\\
355.00	11.75\\
378.00	11.72\\
427.00	11.82\\
455.00	11.80\\
484.00	11.73\\
582.00	11.86\\
619.00	11.84\\
746.00	11.95\\
844.00	11.97\\
897.00	11.89\\
955.00	11.90\\
1016.00	11.84\\
1149.00	11.92\\
1222.00	11.91\\
1300.00	11.96\\
1383.00	11.92\\
1565.00	11.95\\
2005.00	11.95\\
2567.00	11.91\\
2730.00	11.95\\
3090.00	11.92\\
3719.00	11.99\\
3957.00	11.97\\
4477.00	11.98\\
4763.00	11.96\\
5389.00	12.02\\
5733.00	12.00\\
6901.00	12.03\\
10000.00	12.03\\
};
\addplot [color=white!80!black, dashdotted, line width=1.2pt, forget plot]
  table[row sep=crcr]{%
1.00	10.84\\
2.00	7.69\\
3.00	5.96\\
4.00	5.65\\
5.00	7.83\\
6.00	9.86\\
7.00	9.89\\
8.00	10.33\\
9.00	9.71\\
10.00	9.68\\
11.00	10.46\\
12.00	11.14\\
13.00	10.60\\
14.00	10.50\\
16.00	10.50\\
17.00	10.68\\
18.00	11.11\\
19.00	11.11\\
21.00	10.91\\
22.00	10.92\\
23.00	10.89\\
26.00	10.53\\
28.00	10.62\\
30.00	10.75\\
32.00	10.71\\
34.00	10.94\\
36.00	10.76\\
38.00	10.46\\
41.00	10.50\\
43.00	10.62\\
46.00	10.40\\
49.00	10.07\\
52.00	10.31\\
56.00	10.59\\
63.00	10.96\\
67.00	11.38\\
71.00	11.58\\
76.00	11.54\\
81.00	11.80\\
86.00	11.65\\
97.00	11.65\\
103.00	11.87\\
117.00	11.87\\
124.00	12.13\\
149.00	12.22\\
159.00	12.22\\
169.00	12.27\\
191.00	12.28\\
204.00	12.25\\
217.00	12.28\\
245.00	12.25\\
295.00	12.01\\
334.00	11.99\\
355.00	12.04\\
378.00	11.99\\
427.00	12.07\\
455.00	12.05\\
515.00	11.90\\
582.00	11.92\\
619.00	11.83\\
659.00	11.93\\
793.00	11.92\\
844.00	11.82\\
897.00	11.91\\
1016.00	11.88\\
1149.00	11.91\\
1383.00	11.91\\
1472.00	11.98\\
1565.00	11.99\\
1884.00	11.94\\
2005.00	11.94\\
2132.00	11.88\\
2567.00	11.89\\
2905.00	11.87\\
3090.00	11.89\\
3287.00	11.87\\
4477.00	11.95\\
4763.00	11.93\\
5733.00	11.99\\
6488.00	11.97\\
7341.00	11.95\\
8837.00	11.99\\
10000.00	11.98\\
};
\addplot [color=white!90!black, line width=1.2pt, forget plot]
  table[row sep=crcr]{%
1.00	13.68\\
2.00	12.85\\
4.00	12.96\\
5.00	13.25\\
7.00	13.55\\
8.00	14.07\\
9.00	13.56\\
10.00	13.35\\
11.00	13.06\\
12.00	13.48\\
13.00	13.02\\
14.00	13.12\\
16.00	12.49\\
17.00	12.56\\
18.00	12.77\\
19.00	12.31\\
21.00	13.00\\
22.00	12.80\\
23.00	12.94\\
25.00	13.13\\
26.00	13.12\\
28.00	13.06\\
30.00	12.72\\
32.00	12.81\\
34.00	12.76\\
36.00	13.26\\
38.00	13.15\\
41.00	13.17\\
43.00	13.22\\
46.00	13.12\\
49.00	12.84\\
52.00	13.07\\
56.00	13.15\\
59.00	13.27\\
63.00	13.34\\
67.00	13.14\\
71.00	13.02\\
76.00	12.94\\
81.00	13.07\\
86.00	12.85\\
91.00	12.73\\
97.00	12.88\\
103.00	12.80\\
110.00	12.60\\
117.00	12.49\\
124.00	12.47\\
132.00	12.24\\
140.00	12.26\\
149.00	12.39\\
159.00	12.38\\
169.00	12.43\\
180.00	12.42\\
191.00	12.16\\
204.00	12.20\\
217.00	12.31\\
230.00	12.30\\
245.00	12.39\\
295.00	12.48\\
314.00	12.45\\
334.00	12.48\\
355.00	12.47\\
378.00	12.52\\
402.00	12.46\\
455.00	12.37\\
484.00	12.30\\
515.00	12.19\\
582.00	12.19\\
619.00	12.23\\
701.00	12.23\\
746.00	12.11\\
793.00	12.04\\
897.00	12.10\\
955.00	12.06\\
1016.00	12.07\\
1080.00	12.02\\
1149.00	12.04\\
1222.00	11.98\\
1300.00	11.95\\
1472.00	12.08\\
1565.00	12.02\\
1771.00	12.01\\
2005.00	12.01\\
2268.00	11.96\\
2567.00	11.92\\
2730.00	11.98\\
2905.00	11.96\\
3090.00	11.98\\
3496.00	11.99\\
3957.00	12.04\\
4209.00	12.08\\
4763.00	12.11\\
5066.00	12.13\\
6901.00	12.06\\
8837.00	12.09\\
10000.00	12.05\\
};
\addplot [color=black, line width=1.2pt, forget plot]
  table[row sep=crcr]{%
0.00	12.00\\
10000.00	12.00\\
};
\end{axis}

\begin{axis}[%
width=0.411\figurewidth,
height=0.265\figureheight,
at={(0.54\figurewidth,0.368\figureheight)},
scale only axis,
xmin=0.00,
xmax=6.00,
xtick={0.00,2.00,4.00,6.00},
xticklabels={\empty},
ymin=0.00,
ymax=40.00,
ytick={0.00,10.00,20.00,30.00,40.00},
yticklabels={\empty},
axis background/.style={fill=white},
title style={font=\bfseries\color{mycolor1}},
title={\textbf{$\sym{beta}_{12}$}},
xmajorgrids,
ymajorgrids,
grid style={dotted, opacity=0.25}
]
\addplot [color=white!40!black, dashdotted, line width=1.2pt, forget plot]
  table[row sep=crcr]{%
1.00	1.21\\
2.00	2.28\\
3.00	1.31\\
4.00	-0.50\\
5.00	0.34\\
6.00	-0.52\\
};
\addplot [color=gray, line width=1.2pt, forget plot]
  table[row sep=crcr]{%
1.00	0.99\\
1.63	-4.00\\
};
\addplot [color=black, line width=1.2pt, forget plot]
  table[row sep=crcr]{%
0.00	-6.00\\
6.00	-6.00\\
};
\end{axis}

\begin{axis}[%
width=0.411\figurewidth,
height=0.265\figureheight,
at={(0\figurewidth,0\figureheight)},
scale only axis,
xmin=0.00,
xmax=1.00,
xlabel style={font=\color{mycolor1}},
xlabel={Iterationen $\sym{n}_{\sym{idx_MC}}$},
ymin=0.00,
ymax=40.00,
ylabel style={font=\color{mycolor1}},
ylabel={$\hat{\sym{thet}}_{\sym{idx_i}}$},
axis background/.style={fill=white},
title style={font=\bfseries\color{mycolor1}},
title={\textbf{$\sym{beta}_{11}$}},
xmajorgrids,
ymajorgrids,
grid style={dotted, opacity=0.25}
]
\addplot [color=black, line width=1.2pt, forget plot]
  table[row sep=crcr]{%
0.00	-6.00\\
1.00	-6.00\\
};
\end{axis}

\begin{axis}[%
width=0.411\figurewidth,
height=0.265\figureheight,
at={(0.54\figurewidth,0\figureheight)},
scale only axis,
xmin=0.00,
xmax=10000.00,
xlabel style={font=\color{mycolor1}},
xlabel={Iterationen $\sym{n}_{\sym{idx_MC}}$},
ymin=0.00,
ymax=40.00,
ytick={0.00,10.00,20.00,30.00,40.00},
yticklabels={\empty},
axis background/.style={fill=white},
title style={font=\bfseries\color{mycolor1}},
title={\textbf{$\sym{beta}_{22}$}},
xmajorgrids,
ymajorgrids,
grid style={dotted, opacity=0.25}
]
\addplot [color=black, line width=1.2pt, forget plot]
  table[row sep=crcr]{%
0.00	-6.00\\
10000.00	-6.00\\
};
\addplot [color=black, line width=1.2pt, forget plot]
  table[row sep=crcr]{%
0.00	24.00\\
10000.00	24.00\\
};
\addplot [color=black, line width=1.2pt, forget plot]
  table[row sep=crcr]{%
0.00	12.00\\
10000.00	12.00\\
};
\addplot [color=black, line width=1.2pt, forget plot]
  table[row sep=crcr]{%
0.00	12.00\\
10000.00	12.00\\
};
\addplot [color=black, line width=1.2pt, forget plot]
  table[row sep=crcr]{%
0.00	-6.00\\
6.00	-6.00\\
};
\addplot [color=black, line width=1.2pt, forget plot]
  table[row sep=crcr]{%
0.00	-6.00\\
1.00	-6.00\\
};
\end{axis}
\end{tikzpicture}%
    \caption{Konvergenz der \ac{OLS}-Parameterschätzwerte ($\hat{\sym{beta}}_i$). Alle Modelle erweisen sich unabhängig von der geometrischen Verzerrung als erwartungstreu. Die Schätzungen stabilisieren sich bei der $\sym{n}_{\sym{idx_MC}}=10.000$ reproduzierbar.}
    \label{fig:app.c_conv_beta}
\end{figure}

\begin{figure}[htbp]
    \centering
    \setlength\figurewidth{0.5\textwidth}
    \setlength\figureheight{16cm}
    % This file was created by matlab2tikz.
%
\definecolor{mycolor1}{rgb}{0.12941,0.12941,0.12941}%
%
\begin{tikzpicture}

\begin{axis}[%
width=0.411\figurewidth,
height=0.264\figureheight,
at={(0\figurewidth,0.736\figureheight)},
scale only axis,
xmin=0.000,
xmax=1.000,
xtick={0.000,0.500,1.000},
xticklabels={\empty},
ymin=0.000,
ymax=1.000,
ylabel style={font=\color{mycolor1}},
ylabel={Häufigkeit},
axis background/.style={fill=white},
title style={font=\bfseries\color{mycolor1}},
title={\textbf{$\sym{beta}_0$}},
xmajorgrids,
ymajorgrids,
grid style={dotted, opacity=0.25},
legend style={legend cell align=left, align=left}
]
\addplot[const plot, color=white!10!black, line width=1.2pt] table[row sep=crcr] {%
0.000	0.997\\
0.026	0.001\\
0.051	0.000\\
0.077	0.000\\
0.103	0.000\\
0.128	0.000\\
0.154	0.000\\
0.179	0.000\\
0.205	0.000\\
0.231	0.000\\
0.256	0.000\\
0.282	0.000\\
0.308	0.000\\
0.333	0.000\\
0.359	0.000\\
0.385	0.000\\
0.410	0.000\\
0.436	0.000\\
0.462	0.000\\
0.487	0.000\\
0.513	0.000\\
0.538	0.000\\
0.564	0.000\\
0.590	0.000\\
0.615	0.000\\
0.641	0.000\\
0.667	0.000\\
0.692	0.000\\
0.718	0.000\\
0.744	0.000\\
0.769	0.000\\
0.795	0.000\\
0.821	0.000\\
0.846	0.000\\
0.872	0.000\\
0.897	0.000\\
0.923	0.000\\
0.949	0.000\\
0.974	0.000\\
1.000	0.000\\
};
\addlegendentry{Default \ac{CCD}}

\addplot[const plot, color=white!20!black, dashed, line width=1.2pt] table[row sep=crcr] {%
0.000	0.998\\
0.026	0.001\\
0.051	0.000\\
0.077	0.000\\
0.103	0.000\\
0.128	0.000\\
0.154	0.000\\
0.179	0.000\\
0.205	0.000\\
0.231	0.000\\
0.256	0.000\\
0.282	0.000\\
0.308	0.000\\
0.333	0.000\\
0.359	0.000\\
0.385	0.000\\
0.410	0.000\\
0.436	0.000\\
0.462	0.000\\
0.487	0.000\\
0.513	0.000\\
0.538	0.000\\
0.564	0.000\\
0.590	0.000\\
0.615	0.000\\
0.641	0.000\\
0.667	0.000\\
0.692	0.000\\
0.718	0.000\\
0.744	0.000\\
0.769	0.000\\
0.795	0.000\\
0.821	0.000\\
0.846	0.000\\
0.872	0.000\\
0.897	0.000\\
0.923	0.000\\
0.949	0.000\\
0.974	0.000\\
1.000	0.000\\
};
\addlegendentry{Fall I ($-20\%$)}

\addplot[const plot, color=white!30!black, dotted, line width=1.2pt] table[row sep=crcr] {%
0.000	0.998\\
0.026	0.001\\
0.051	0.000\\
0.077	0.000\\
0.103	0.000\\
0.128	0.000\\
0.154	0.000\\
0.179	0.000\\
0.205	0.000\\
0.231	0.000\\
0.256	0.000\\
0.282	0.000\\
0.308	0.000\\
0.333	0.000\\
0.359	0.000\\
0.385	0.000\\
0.410	0.000\\
0.436	0.000\\
0.462	0.000\\
0.487	0.000\\
0.513	0.000\\
0.538	0.000\\
0.564	0.000\\
0.590	0.000\\
0.615	0.000\\
0.641	0.000\\
0.667	0.000\\
0.692	0.000\\
0.718	0.000\\
0.744	0.000\\
0.769	0.000\\
0.795	0.000\\
0.821	0.000\\
0.846	0.000\\
0.872	0.000\\
0.897	0.000\\
0.923	0.000\\
0.949	0.000\\
0.974	0.000\\
1.000	0.000\\
};
\addlegendentry{Fall I ($+20\%$)}

\addplot[const plot, color=white!40!black, dashdotted, line width=1.2pt] table[row sep=crcr] {%
0.000	0.998\\
0.026	0.000\\
0.051	0.001\\
0.077	0.000\\
0.103	0.000\\
0.128	0.000\\
0.154	0.000\\
0.179	0.000\\
0.205	0.000\\
0.231	0.000\\
0.256	0.000\\
0.282	0.000\\
0.308	0.000\\
0.333	0.000\\
0.359	0.000\\
0.385	0.000\\
0.410	0.000\\
0.436	0.000\\
0.462	0.000\\
0.487	0.000\\
0.513	0.000\\
0.538	0.000\\
0.564	0.000\\
0.590	0.000\\
0.615	0.000\\
0.641	0.000\\
0.667	0.000\\
0.692	0.000\\
0.718	0.000\\
0.744	0.000\\
0.769	0.000\\
0.795	0.000\\
0.821	0.000\\
0.846	0.000\\
0.872	0.000\\
0.897	0.000\\
0.923	0.000\\
0.949	0.000\\
0.974	0.000\\
1.000	0.000\\
};
\addlegendentry{Fall I (\ac{FC})}

\addplot[const plot, color=gray, line width=1.2pt] table[row sep=crcr] {%
0.000	0.998\\
0.026	0.001\\
0.051	0.000\\
0.077	0.000\\
0.103	0.000\\
0.128	0.000\\
0.154	0.000\\
0.179	0.000\\
0.205	0.000\\
0.231	0.000\\
0.256	0.000\\
0.282	0.000\\
0.308	0.000\\
0.333	0.000\\
0.359	0.000\\
0.385	0.000\\
0.410	0.000\\
0.436	0.000\\
0.462	0.000\\
0.487	0.000\\
0.513	0.000\\
0.538	0.000\\
0.564	0.000\\
0.590	0.000\\
0.615	0.000\\
0.641	0.000\\
0.667	0.000\\
0.692	0.000\\
0.718	0.000\\
0.744	0.000\\
0.769	0.000\\
0.795	0.000\\
0.821	0.000\\
0.846	0.000\\
0.872	0.000\\
0.897	0.000\\
0.923	0.000\\
0.949	0.000\\
0.974	0.000\\
1.000	0.000\\
};
\addlegendentry{Fall II (leicht)}

\addplot[const plot, color=white!60!black, dashed, line width=1.2pt] table[row sep=crcr] {%
0.000	0.998\\
0.026	0.001\\
0.051	0.000\\
0.077	0.000\\
0.103	0.000\\
0.128	0.000\\
0.154	0.000\\
0.179	0.000\\
0.205	0.000\\
0.231	0.000\\
0.256	0.000\\
0.282	0.000\\
0.308	0.000\\
0.333	0.000\\
0.359	0.000\\
0.385	0.000\\
0.410	0.000\\
0.436	0.000\\
0.462	0.000\\
0.487	0.000\\
0.513	0.000\\
0.538	0.000\\
0.564	0.000\\
0.590	0.000\\
0.615	0.000\\
0.641	0.000\\
0.667	0.000\\
0.692	0.000\\
0.718	0.000\\
0.744	0.000\\
0.769	0.000\\
0.795	0.000\\
0.821	0.000\\
0.846	0.000\\
0.872	0.000\\
0.897	0.000\\
0.923	0.000\\
0.949	0.000\\
0.974	0.000\\
1.000	0.000\\
};
\addlegendentry{Fall II (stark)}

\addplot[const plot, color=white!70!black, dotted, line width=1.2pt] table[row sep=crcr] {%
0.000	0.998\\
0.026	0.001\\
0.051	0.000\\
0.077	0.000\\
0.103	0.000\\
0.128	0.000\\
0.154	0.000\\
0.179	0.000\\
0.205	0.000\\
0.231	0.000\\
0.256	0.000\\
0.282	0.000\\
0.308	0.000\\
0.333	0.000\\
0.359	0.000\\
0.385	0.000\\
0.410	0.000\\
0.436	0.000\\
0.462	0.000\\
0.487	0.000\\
0.513	0.000\\
0.538	0.000\\
0.564	0.000\\
0.590	0.000\\
0.615	0.000\\
0.641	0.000\\
0.667	0.000\\
0.692	0.000\\
0.718	0.000\\
0.744	0.000\\
0.769	0.000\\
0.795	0.000\\
0.821	0.000\\
0.846	0.000\\
0.872	0.000\\
0.897	0.000\\
0.923	0.000\\
0.949	0.000\\
0.974	0.000\\
1.000	0.000\\
};
\addlegendentry{Fall III}

\addplot[const plot, color=white!80!black, dashdotted, line width=1.2pt] table[row sep=crcr] {%
0.000	0.998\\
0.026	0.000\\
0.051	0.001\\
0.077	0.000\\
0.103	0.000\\
0.128	0.000\\
0.154	0.000\\
0.179	0.000\\
0.205	0.000\\
0.231	0.000\\
0.256	0.000\\
0.282	0.000\\
0.308	0.000\\
0.333	0.000\\
0.359	0.000\\
0.385	0.000\\
0.410	0.000\\
0.436	0.000\\
0.462	0.000\\
0.487	0.000\\
0.513	0.000\\
0.538	0.000\\
0.564	0.000\\
0.590	0.000\\
0.615	0.000\\
0.641	0.000\\
0.667	0.000\\
0.692	0.000\\
0.718	0.000\\
0.744	0.000\\
0.769	0.000\\
0.795	0.000\\
0.821	0.000\\
0.846	0.000\\
0.872	0.000\\
0.897	0.000\\
0.923	0.000\\
0.949	0.000\\
0.974	0.000\\
1.000	0.000\\
};
\addlegendentry{Fall IV: \ac{SP}}

\addplot[const plot, color=white!90!black, line width=1.2pt] table[row sep=crcr] {%
0.000	0.997\\
0.026	0.001\\
0.051	0.001\\
0.077	0.000\\
0.103	0.000\\
0.128	0.000\\
0.154	0.000\\
0.179	0.000\\
0.205	0.000\\
0.231	0.000\\
0.256	0.000\\
0.282	0.000\\
0.308	0.000\\
0.333	0.000\\
0.359	0.000\\
0.385	0.000\\
0.410	0.000\\
0.436	0.000\\
0.462	0.000\\
0.487	0.000\\
0.513	0.000\\
0.538	0.000\\
0.564	0.000\\
0.590	0.000\\
0.615	0.000\\
0.641	0.000\\
0.667	0.000\\
0.692	0.000\\
0.718	0.000\\
0.744	0.000\\
0.769	0.000\\
0.795	0.000\\
0.821	0.000\\
0.846	0.000\\
0.872	0.000\\
0.897	0.000\\
0.923	0.000\\
0.949	0.000\\
0.974	0.000\\
1.000	0.000\\
};
\addlegendentry{Fall IV: \ac{ZP}}

\addplot [color=black, line width=1.2pt, forget plot]
  table[row sep=crcr]{%
0.050	0.000\\
0.050	1.000\\
};
\end{axis}

\begin{axis}[%
width=0.411\figurewidth,
height=0.264\figureheight,
at={(0.541\figurewidth,0.736\figureheight)},
scale only axis,
xmin=0.000,
xmax=1.000,
xtick={0.000,0.500,1.000},
xticklabels={\empty},
ymin=0.000,
ymax=1.000,
ytick={0.000,0.200,0.400,0.600,0.800,1.000},
yticklabels={\empty},
axis background/.style={fill=white},
title style={font=\bfseries\color{mycolor1}},
title={\textbf{$\sym{beta}_1$}},
xmajorgrids,
ymajorgrids,
grid style={dotted, opacity=0.25}
]
\addplot[const plot, color=white!10!black, line width=1.2pt, forget plot] table[row sep=crcr] {%
0.000	0.677\\
0.026	0.055\\
0.051	0.035\\
0.077	0.022\\
0.103	0.016\\
0.128	0.015\\
0.154	0.013\\
0.179	0.010\\
0.205	0.009\\
0.231	0.010\\
0.256	0.009\\
0.282	0.006\\
0.308	0.007\\
0.333	0.007\\
0.359	0.006\\
0.385	0.006\\
0.410	0.005\\
0.436	0.005\\
0.462	0.005\\
0.487	0.006\\
0.513	0.005\\
0.538	0.004\\
0.564	0.004\\
0.590	0.004\\
0.615	0.003\\
0.641	0.003\\
0.667	0.005\\
0.692	0.004\\
0.718	0.004\\
0.744	0.004\\
0.769	0.003\\
0.795	0.005\\
0.821	0.004\\
0.846	0.003\\
0.872	0.004\\
0.897	0.003\\
0.923	0.004\\
0.949	0.004\\
0.974	0.003\\
1.000	0.003\\
};
\addplot[const plot, color=white!20!black, dashed, line width=1.2pt, forget plot] table[row sep=crcr] {%
0.000	0.657\\
0.026	0.059\\
0.051	0.034\\
0.077	0.022\\
0.103	0.017\\
0.128	0.018\\
0.154	0.011\\
0.179	0.013\\
0.205	0.009\\
0.231	0.009\\
0.256	0.008\\
0.282	0.007\\
0.308	0.010\\
0.333	0.007\\
0.359	0.006\\
0.385	0.006\\
0.410	0.005\\
0.436	0.005\\
0.462	0.005\\
0.487	0.006\\
0.513	0.006\\
0.538	0.005\\
0.564	0.004\\
0.590	0.003\\
0.615	0.005\\
0.641	0.005\\
0.667	0.004\\
0.692	0.004\\
0.718	0.006\\
0.744	0.005\\
0.769	0.004\\
0.795	0.004\\
0.821	0.004\\
0.846	0.005\\
0.872	0.004\\
0.897	0.004\\
0.923	0.004\\
0.949	0.004\\
0.974	0.004\\
1.000	0.004\\
};
\addplot[const plot, color=white!30!black, dotted, line width=1.2pt, forget plot] table[row sep=crcr] {%
0.000	0.686\\
0.026	0.050\\
0.051	0.032\\
0.077	0.021\\
0.103	0.019\\
0.128	0.014\\
0.154	0.011\\
0.179	0.010\\
0.205	0.009\\
0.231	0.009\\
0.256	0.007\\
0.282	0.007\\
0.308	0.009\\
0.333	0.007\\
0.359	0.006\\
0.385	0.005\\
0.410	0.006\\
0.436	0.005\\
0.462	0.005\\
0.487	0.004\\
0.513	0.006\\
0.538	0.004\\
0.564	0.004\\
0.590	0.004\\
0.615	0.003\\
0.641	0.004\\
0.667	0.004\\
0.692	0.005\\
0.718	0.005\\
0.744	0.004\\
0.769	0.004\\
0.795	0.004\\
0.821	0.004\\
0.846	0.004\\
0.872	0.004\\
0.897	0.004\\
0.923	0.003\\
0.949	0.004\\
0.974	0.004\\
1.000	0.004\\
};
\addplot[const plot, color=white!40!black, dashdotted, line width=1.2pt, forget plot] table[row sep=crcr] {%
0.000	0.638\\
0.026	0.059\\
0.051	0.035\\
0.077	0.022\\
0.103	0.019\\
0.128	0.018\\
0.154	0.015\\
0.179	0.014\\
0.205	0.010\\
0.231	0.008\\
0.256	0.009\\
0.282	0.010\\
0.308	0.007\\
0.333	0.007\\
0.359	0.008\\
0.385	0.008\\
0.410	0.005\\
0.436	0.007\\
0.462	0.006\\
0.487	0.005\\
0.513	0.006\\
0.538	0.005\\
0.564	0.005\\
0.590	0.005\\
0.615	0.006\\
0.641	0.004\\
0.667	0.004\\
0.692	0.007\\
0.718	0.004\\
0.744	0.004\\
0.769	0.005\\
0.795	0.004\\
0.821	0.005\\
0.846	0.004\\
0.872	0.005\\
0.897	0.005\\
0.923	0.004\\
0.949	0.005\\
0.974	0.004\\
1.000	0.004\\
};
\addplot[const plot, color=gray, line width=1.2pt, forget plot] table[row sep=crcr] {%
0.000	0.673\\
0.026	0.057\\
0.051	0.033\\
0.077	0.022\\
0.103	0.018\\
0.128	0.014\\
0.154	0.013\\
0.179	0.011\\
0.205	0.009\\
0.231	0.009\\
0.256	0.008\\
0.282	0.007\\
0.308	0.009\\
0.333	0.006\\
0.359	0.006\\
0.385	0.006\\
0.410	0.006\\
0.436	0.005\\
0.462	0.005\\
0.487	0.004\\
0.513	0.004\\
0.538	0.005\\
0.564	0.004\\
0.590	0.005\\
0.615	0.005\\
0.641	0.004\\
0.667	0.004\\
0.692	0.004\\
0.718	0.004\\
0.744	0.003\\
0.769	0.004\\
0.795	0.004\\
0.821	0.003\\
0.846	0.004\\
0.872	0.003\\
0.897	0.004\\
0.923	0.004\\
0.949	0.004\\
0.974	0.004\\
1.000	0.004\\
};
\addplot[const plot, color=white!60!black, dashed, line width=1.2pt, forget plot] table[row sep=crcr] {%
0.000	0.646\\
0.026	0.056\\
0.051	0.034\\
0.077	0.026\\
0.103	0.020\\
0.128	0.017\\
0.154	0.015\\
0.179	0.011\\
0.205	0.010\\
0.231	0.010\\
0.256	0.010\\
0.282	0.010\\
0.308	0.007\\
0.333	0.007\\
0.359	0.006\\
0.385	0.006\\
0.410	0.006\\
0.436	0.007\\
0.462	0.006\\
0.487	0.006\\
0.513	0.004\\
0.538	0.005\\
0.564	0.006\\
0.590	0.005\\
0.615	0.004\\
0.641	0.006\\
0.667	0.005\\
0.692	0.004\\
0.718	0.004\\
0.744	0.004\\
0.769	0.004\\
0.795	0.004\\
0.821	0.004\\
0.846	0.004\\
0.872	0.004\\
0.897	0.004\\
0.923	0.004\\
0.949	0.004\\
0.974	0.005\\
1.000	0.005\\
};
\addplot[const plot, color=white!70!black, dotted, line width=1.2pt, forget plot] table[row sep=crcr] {%
0.000	0.660\\
0.026	0.057\\
0.051	0.034\\
0.077	0.023\\
0.103	0.015\\
0.128	0.015\\
0.154	0.013\\
0.179	0.012\\
0.205	0.012\\
0.231	0.008\\
0.256	0.009\\
0.282	0.008\\
0.308	0.008\\
0.333	0.007\\
0.359	0.006\\
0.385	0.006\\
0.410	0.005\\
0.436	0.005\\
0.462	0.006\\
0.487	0.006\\
0.513	0.005\\
0.538	0.005\\
0.564	0.004\\
0.590	0.005\\
0.615	0.004\\
0.641	0.005\\
0.667	0.005\\
0.692	0.003\\
0.718	0.005\\
0.744	0.005\\
0.769	0.005\\
0.795	0.005\\
0.821	0.005\\
0.846	0.003\\
0.872	0.004\\
0.897	0.004\\
0.923	0.005\\
0.949	0.004\\
0.974	0.004\\
1.000	0.004\\
};
\addplot[const plot, color=white!80!black, dashdotted, line width=1.2pt, forget plot] table[row sep=crcr] {%
0.000	0.570\\
0.026	0.070\\
0.051	0.040\\
0.077	0.029\\
0.103	0.021\\
0.128	0.021\\
0.154	0.016\\
0.179	0.017\\
0.205	0.015\\
0.231	0.012\\
0.256	0.010\\
0.282	0.012\\
0.308	0.008\\
0.333	0.007\\
0.359	0.008\\
0.385	0.007\\
0.410	0.007\\
0.436	0.008\\
0.462	0.006\\
0.487	0.007\\
0.513	0.007\\
0.538	0.006\\
0.564	0.006\\
0.590	0.007\\
0.615	0.005\\
0.641	0.006\\
0.667	0.007\\
0.692	0.005\\
0.718	0.006\\
0.744	0.005\\
0.769	0.005\\
0.795	0.004\\
0.821	0.005\\
0.846	0.006\\
0.872	0.004\\
0.897	0.005\\
0.923	0.005\\
0.949	0.007\\
0.974	0.006\\
1.000	0.006\\
};
\addplot[const plot, color=white!90!black, line width=1.2pt, forget plot] table[row sep=crcr] {%
0.000	0.666\\
0.026	0.059\\
0.051	0.030\\
0.077	0.022\\
0.103	0.017\\
0.128	0.016\\
0.154	0.014\\
0.179	0.011\\
0.205	0.011\\
0.231	0.009\\
0.256	0.009\\
0.282	0.007\\
0.308	0.006\\
0.333	0.007\\
0.359	0.006\\
0.385	0.007\\
0.410	0.005\\
0.436	0.005\\
0.462	0.006\\
0.487	0.004\\
0.513	0.005\\
0.538	0.005\\
0.564	0.005\\
0.590	0.005\\
0.615	0.005\\
0.641	0.004\\
0.667	0.005\\
0.692	0.004\\
0.718	0.004\\
0.744	0.005\\
0.769	0.004\\
0.795	0.004\\
0.821	0.003\\
0.846	0.004\\
0.872	0.004\\
0.897	0.004\\
0.923	0.005\\
0.949	0.003\\
0.974	0.004\\
1.000	0.004\\
};
\addplot [color=black, line width=1.2pt, forget plot]
  table[row sep=crcr]{%
0.050	0.000\\
0.050	1.000\\
};
\end{axis}

\begin{axis}[%
width=0.411\figurewidth,
height=0.264\figureheight,
at={(0\figurewidth,0.368\figureheight)},
scale only axis,
xmin=0.000,
xmax=1.000,
xtick={0.000,0.500,1.000},
xticklabels={\empty},
ymin=0.000,
ymax=1.000,
ylabel style={font=\color{mycolor1}},
ylabel={Häufigkeit},
axis background/.style={fill=white},
title style={font=\bfseries\color{mycolor1}},
title={\textbf{$\sym{beta}_2$}},
xmajorgrids,
ymajorgrids,
grid style={dotted, opacity=0.25}
]
\addplot[const plot, color=white!10!black, line width=1.2pt, forget plot] table[row sep=crcr] {%
0.000	0.683\\
0.026	0.052\\
0.051	0.031\\
0.077	0.022\\
0.103	0.016\\
0.128	0.016\\
0.154	0.011\\
0.179	0.010\\
0.205	0.009\\
0.231	0.009\\
0.256	0.007\\
0.282	0.008\\
0.308	0.007\\
0.333	0.006\\
0.359	0.007\\
0.385	0.006\\
0.410	0.005\\
0.436	0.004\\
0.462	0.004\\
0.487	0.005\\
0.513	0.004\\
0.538	0.005\\
0.564	0.005\\
0.590	0.005\\
0.615	0.005\\
0.641	0.003\\
0.667	0.005\\
0.692	0.004\\
0.718	0.004\\
0.744	0.004\\
0.769	0.005\\
0.795	0.004\\
0.821	0.004\\
0.846	0.004\\
0.872	0.004\\
0.897	0.003\\
0.923	0.003\\
0.949	0.003\\
0.974	0.004\\
1.000	0.004\\
};
\addplot[const plot, color=white!20!black, dashed, line width=1.2pt, forget plot] table[row sep=crcr] {%
0.000	0.676\\
0.026	0.049\\
0.051	0.030\\
0.077	0.024\\
0.103	0.017\\
0.128	0.016\\
0.154	0.013\\
0.179	0.010\\
0.205	0.010\\
0.231	0.008\\
0.256	0.008\\
0.282	0.008\\
0.308	0.007\\
0.333	0.007\\
0.359	0.006\\
0.385	0.005\\
0.410	0.005\\
0.436	0.006\\
0.462	0.005\\
0.487	0.005\\
0.513	0.005\\
0.538	0.005\\
0.564	0.006\\
0.590	0.005\\
0.615	0.004\\
0.641	0.004\\
0.667	0.004\\
0.692	0.004\\
0.718	0.005\\
0.744	0.004\\
0.769	0.004\\
0.795	0.004\\
0.821	0.003\\
0.846	0.004\\
0.872	0.005\\
0.897	0.004\\
0.923	0.005\\
0.949	0.004\\
0.974	0.004\\
1.000	0.004\\
};
\addplot[const plot, color=white!30!black, dotted, line width=1.2pt, forget plot] table[row sep=crcr] {%
0.000	0.680\\
0.026	0.053\\
0.051	0.029\\
0.077	0.022\\
0.103	0.018\\
0.128	0.015\\
0.154	0.013\\
0.179	0.009\\
0.205	0.009\\
0.231	0.011\\
0.256	0.008\\
0.282	0.008\\
0.308	0.006\\
0.333	0.006\\
0.359	0.008\\
0.385	0.006\\
0.410	0.007\\
0.436	0.006\\
0.462	0.005\\
0.487	0.006\\
0.513	0.005\\
0.538	0.005\\
0.564	0.006\\
0.590	0.003\\
0.615	0.004\\
0.641	0.004\\
0.667	0.005\\
0.692	0.005\\
0.718	0.004\\
0.744	0.004\\
0.769	0.003\\
0.795	0.004\\
0.821	0.003\\
0.846	0.004\\
0.872	0.003\\
0.897	0.004\\
0.923	0.005\\
0.949	0.003\\
0.974	0.003\\
1.000	0.003\\
};
\addplot[const plot, color=white!40!black, dashdotted, line width=1.2pt, forget plot] table[row sep=crcr] {%
0.000	0.674\\
0.026	0.058\\
0.051	0.033\\
0.077	0.022\\
0.103	0.018\\
0.128	0.013\\
0.154	0.011\\
0.179	0.012\\
0.205	0.007\\
0.231	0.009\\
0.256	0.008\\
0.282	0.006\\
0.308	0.007\\
0.333	0.008\\
0.359	0.007\\
0.385	0.005\\
0.410	0.005\\
0.436	0.006\\
0.462	0.005\\
0.487	0.005\\
0.513	0.005\\
0.538	0.005\\
0.564	0.005\\
0.590	0.005\\
0.615	0.005\\
0.641	0.005\\
0.667	0.005\\
0.692	0.004\\
0.718	0.004\\
0.744	0.004\\
0.769	0.004\\
0.795	0.004\\
0.821	0.003\\
0.846	0.003\\
0.872	0.004\\
0.897	0.004\\
0.923	0.004\\
0.949	0.004\\
0.974	0.004\\
1.000	0.004\\
};
\addplot[const plot, color=gray, line width=1.2pt, forget plot] table[row sep=crcr] {%
0.000	0.675\\
0.026	0.053\\
0.051	0.032\\
0.077	0.021\\
0.103	0.016\\
0.128	0.017\\
0.154	0.013\\
0.179	0.010\\
0.205	0.009\\
0.231	0.009\\
0.256	0.009\\
0.282	0.008\\
0.308	0.007\\
0.333	0.008\\
0.359	0.008\\
0.385	0.004\\
0.410	0.006\\
0.436	0.006\\
0.462	0.005\\
0.487	0.005\\
0.513	0.004\\
0.538	0.004\\
0.564	0.005\\
0.590	0.004\\
0.615	0.005\\
0.641	0.004\\
0.667	0.004\\
0.692	0.005\\
0.718	0.004\\
0.744	0.003\\
0.769	0.004\\
0.795	0.004\\
0.821	0.004\\
0.846	0.004\\
0.872	0.005\\
0.897	0.004\\
0.923	0.003\\
0.949	0.006\\
0.974	0.003\\
1.000	0.003\\
};
\addplot[const plot, color=white!60!black, dashed, line width=1.2pt, forget plot] table[row sep=crcr] {%
0.000	0.679\\
0.026	0.054\\
0.051	0.030\\
0.077	0.024\\
0.103	0.016\\
0.128	0.013\\
0.154	0.013\\
0.179	0.010\\
0.205	0.010\\
0.231	0.008\\
0.256	0.009\\
0.282	0.007\\
0.308	0.005\\
0.333	0.007\\
0.359	0.005\\
0.385	0.007\\
0.410	0.007\\
0.436	0.004\\
0.462	0.004\\
0.487	0.004\\
0.513	0.004\\
0.538	0.004\\
0.564	0.005\\
0.590	0.004\\
0.615	0.004\\
0.641	0.004\\
0.667	0.004\\
0.692	0.003\\
0.718	0.004\\
0.744	0.005\\
0.769	0.005\\
0.795	0.004\\
0.821	0.005\\
0.846	0.004\\
0.872	0.002\\
0.897	0.004\\
0.923	0.005\\
0.949	0.005\\
0.974	0.004\\
1.000	0.004\\
};
\addplot[const plot, color=white!70!black, dotted, line width=1.2pt, forget plot] table[row sep=crcr] {%
0.000	0.687\\
0.026	0.054\\
0.051	0.031\\
0.077	0.020\\
0.103	0.017\\
0.128	0.016\\
0.154	0.013\\
0.179	0.010\\
0.205	0.009\\
0.231	0.008\\
0.256	0.008\\
0.282	0.007\\
0.308	0.007\\
0.333	0.007\\
0.359	0.005\\
0.385	0.007\\
0.410	0.006\\
0.436	0.005\\
0.462	0.005\\
0.487	0.005\\
0.513	0.004\\
0.538	0.004\\
0.564	0.005\\
0.590	0.004\\
0.615	0.004\\
0.641	0.005\\
0.667	0.004\\
0.692	0.004\\
0.718	0.004\\
0.744	0.004\\
0.769	0.004\\
0.795	0.003\\
0.821	0.003\\
0.846	0.003\\
0.872	0.004\\
0.897	0.004\\
0.923	0.004\\
0.949	0.003\\
0.974	0.003\\
1.000	0.003\\
};
\addplot[const plot, color=white!80!black, dashdotted, line width=1.2pt, forget plot] table[row sep=crcr] {%
0.000	0.678\\
0.026	0.054\\
0.051	0.034\\
0.077	0.021\\
0.103	0.017\\
0.128	0.015\\
0.154	0.011\\
0.179	0.011\\
0.205	0.009\\
0.231	0.009\\
0.256	0.008\\
0.282	0.007\\
0.308	0.008\\
0.333	0.006\\
0.359	0.006\\
0.385	0.006\\
0.410	0.004\\
0.436	0.005\\
0.462	0.006\\
0.487	0.005\\
0.513	0.004\\
0.538	0.004\\
0.564	0.005\\
0.590	0.004\\
0.615	0.005\\
0.641	0.004\\
0.667	0.004\\
0.692	0.005\\
0.718	0.004\\
0.744	0.004\\
0.769	0.005\\
0.795	0.004\\
0.821	0.003\\
0.846	0.004\\
0.872	0.004\\
0.897	0.004\\
0.923	0.003\\
0.949	0.005\\
0.974	0.004\\
1.000	0.004\\
};
\addplot[const plot, color=white!90!black, line width=1.2pt, forget plot] table[row sep=crcr] {%
0.000	0.673\\
0.026	0.055\\
0.051	0.032\\
0.077	0.022\\
0.103	0.015\\
0.128	0.016\\
0.154	0.012\\
0.179	0.012\\
0.205	0.010\\
0.231	0.009\\
0.256	0.009\\
0.282	0.008\\
0.308	0.008\\
0.333	0.007\\
0.359	0.006\\
0.385	0.005\\
0.410	0.005\\
0.436	0.005\\
0.462	0.005\\
0.487	0.005\\
0.513	0.006\\
0.538	0.005\\
0.564	0.006\\
0.590	0.004\\
0.615	0.004\\
0.641	0.005\\
0.667	0.004\\
0.692	0.004\\
0.718	0.004\\
0.744	0.003\\
0.769	0.004\\
0.795	0.005\\
0.821	0.004\\
0.846	0.004\\
0.872	0.004\\
0.897	0.005\\
0.923	0.004\\
0.949	0.003\\
0.974	0.005\\
1.000	0.005\\
};
\addplot [color=black, line width=1.2pt, forget plot]
  table[row sep=crcr]{%
0.050	0.000\\
0.050	1.000\\
};
\end{axis}

\begin{axis}[%
width=0.411\figurewidth,
height=0.264\figureheight,
at={(0.541\figurewidth,0.368\figureheight)},
scale only axis,
xmin=0.000,
xmax=1.000,
xtick={0.000,0.500,1.000},
xticklabels={\empty},
ymin=0.000,
ymax=1.000,
ytick={0.000,0.200,0.400,0.600,0.800,1.000},
yticklabels={\empty},
axis background/.style={fill=white},
title style={font=\bfseries\color{mycolor1}},
title={\textbf{$\sym{beta}_{12}$}},
xmajorgrids,
ymajorgrids,
grid style={dotted, opacity=0.25}
]
\addplot[const plot, color=white!10!black, line width=1.2pt, forget plot] table[row sep=crcr] {%
0.000	0.534\\
0.026	0.075\\
0.051	0.049\\
0.077	0.028\\
0.103	0.028\\
0.128	0.020\\
0.154	0.018\\
0.179	0.015\\
0.205	0.014\\
0.231	0.012\\
0.256	0.011\\
0.282	0.011\\
0.308	0.009\\
0.333	0.011\\
0.359	0.010\\
0.385	0.008\\
0.410	0.009\\
0.436	0.008\\
0.462	0.007\\
0.487	0.008\\
0.513	0.006\\
0.538	0.008\\
0.564	0.007\\
0.590	0.006\\
0.615	0.007\\
0.641	0.007\\
0.667	0.006\\
0.692	0.006\\
0.718	0.005\\
0.744	0.006\\
0.769	0.006\\
0.795	0.005\\
0.821	0.006\\
0.846	0.006\\
0.872	0.005\\
0.897	0.005\\
0.923	0.004\\
0.949	0.007\\
0.974	0.005\\
1.000	0.005\\
};
\addplot[const plot, color=white!20!black, dashed, line width=1.2pt, forget plot] table[row sep=crcr] {%
0.000	0.523\\
0.026	0.074\\
0.051	0.044\\
0.077	0.034\\
0.103	0.025\\
0.128	0.023\\
0.154	0.018\\
0.179	0.016\\
0.205	0.014\\
0.231	0.013\\
0.256	0.012\\
0.282	0.011\\
0.308	0.012\\
0.333	0.009\\
0.359	0.010\\
0.385	0.009\\
0.410	0.008\\
0.436	0.011\\
0.462	0.007\\
0.487	0.007\\
0.513	0.009\\
0.538	0.007\\
0.564	0.006\\
0.590	0.008\\
0.615	0.008\\
0.641	0.006\\
0.667	0.007\\
0.692	0.006\\
0.718	0.005\\
0.744	0.005\\
0.769	0.006\\
0.795	0.007\\
0.821	0.007\\
0.846	0.006\\
0.872	0.005\\
0.897	0.005\\
0.923	0.005\\
0.949	0.005\\
0.974	0.005\\
1.000	0.005\\
};
\addplot[const plot, color=white!30!black, dotted, line width=1.2pt, forget plot] table[row sep=crcr] {%
0.000	0.528\\
0.026	0.072\\
0.051	0.046\\
0.077	0.033\\
0.103	0.023\\
0.128	0.020\\
0.154	0.019\\
0.179	0.017\\
0.205	0.013\\
0.231	0.014\\
0.256	0.012\\
0.282	0.011\\
0.308	0.010\\
0.333	0.009\\
0.359	0.009\\
0.385	0.010\\
0.410	0.007\\
0.436	0.007\\
0.462	0.008\\
0.487	0.008\\
0.513	0.008\\
0.538	0.010\\
0.564	0.006\\
0.590	0.006\\
0.615	0.008\\
0.641	0.006\\
0.667	0.007\\
0.692	0.006\\
0.718	0.007\\
0.744	0.005\\
0.769	0.006\\
0.795	0.007\\
0.821	0.007\\
0.846	0.005\\
0.872	0.005\\
0.897	0.006\\
0.923	0.007\\
0.949	0.005\\
0.974	0.007\\
1.000	0.007\\
};
\addplot[const plot, color=white!40!black, dashdotted, line width=1.2pt, forget plot] table[row sep=crcr] {%
0.000	0.536\\
0.026	0.072\\
0.051	0.045\\
0.077	0.030\\
0.103	0.026\\
0.128	0.019\\
0.154	0.019\\
0.179	0.015\\
0.205	0.014\\
0.231	0.014\\
0.256	0.010\\
0.282	0.009\\
0.308	0.011\\
0.333	0.011\\
0.359	0.010\\
0.385	0.009\\
0.410	0.009\\
0.436	0.008\\
0.462	0.006\\
0.487	0.007\\
0.513	0.007\\
0.538	0.007\\
0.564	0.007\\
0.590	0.006\\
0.615	0.007\\
0.641	0.006\\
0.667	0.006\\
0.692	0.006\\
0.718	0.007\\
0.744	0.005\\
0.769	0.007\\
0.795	0.007\\
0.821	0.006\\
0.846	0.006\\
0.872	0.006\\
0.897	0.006\\
0.923	0.006\\
0.949	0.005\\
0.974	0.004\\
1.000	0.004\\
};
\addplot[const plot, color=gray, line width=1.2pt, forget plot] table[row sep=crcr] {%
0.000	0.537\\
0.026	0.069\\
0.051	0.044\\
0.077	0.034\\
0.103	0.027\\
0.128	0.020\\
0.154	0.018\\
0.179	0.015\\
0.205	0.012\\
0.231	0.013\\
0.256	0.012\\
0.282	0.011\\
0.308	0.010\\
0.333	0.010\\
0.359	0.009\\
0.385	0.009\\
0.410	0.009\\
0.436	0.008\\
0.462	0.009\\
0.487	0.009\\
0.513	0.008\\
0.538	0.007\\
0.564	0.007\\
0.590	0.008\\
0.615	0.008\\
0.641	0.006\\
0.667	0.005\\
0.692	0.005\\
0.718	0.005\\
0.744	0.005\\
0.769	0.005\\
0.795	0.005\\
0.821	0.005\\
0.846	0.004\\
0.872	0.005\\
0.897	0.005\\
0.923	0.005\\
0.949	0.005\\
0.974	0.006\\
1.000	0.006\\
};
\addplot[const plot, color=white!60!black, dashed, line width=1.2pt, forget plot] table[row sep=crcr] {%
0.000	0.565\\
0.026	0.069\\
0.051	0.039\\
0.077	0.029\\
0.103	0.024\\
0.128	0.021\\
0.154	0.019\\
0.179	0.016\\
0.205	0.014\\
0.231	0.012\\
0.256	0.010\\
0.282	0.013\\
0.308	0.009\\
0.333	0.009\\
0.359	0.008\\
0.385	0.008\\
0.410	0.007\\
0.436	0.007\\
0.462	0.005\\
0.487	0.007\\
0.513	0.007\\
0.538	0.006\\
0.564	0.006\\
0.590	0.004\\
0.615	0.005\\
0.641	0.005\\
0.667	0.006\\
0.692	0.007\\
0.718	0.006\\
0.744	0.005\\
0.769	0.006\\
0.795	0.005\\
0.821	0.005\\
0.846	0.006\\
0.872	0.005\\
0.897	0.006\\
0.923	0.006\\
0.949	0.005\\
0.974	0.006\\
1.000	0.006\\
};
\addplot[const plot, color=white!70!black, dotted, line width=1.2pt, forget plot] table[row sep=crcr] {%
0.000	0.532\\
0.026	0.070\\
0.051	0.044\\
0.077	0.033\\
0.103	0.028\\
0.128	0.021\\
0.154	0.017\\
0.179	0.016\\
0.205	0.015\\
0.231	0.013\\
0.256	0.012\\
0.282	0.011\\
0.308	0.010\\
0.333	0.011\\
0.359	0.008\\
0.385	0.009\\
0.410	0.009\\
0.436	0.008\\
0.462	0.009\\
0.487	0.008\\
0.513	0.007\\
0.538	0.006\\
0.564	0.005\\
0.590	0.006\\
0.615	0.006\\
0.641	0.007\\
0.667	0.006\\
0.692	0.006\\
0.718	0.007\\
0.744	0.006\\
0.769	0.006\\
0.795	0.006\\
0.821	0.006\\
0.846	0.006\\
0.872	0.006\\
0.897	0.006\\
0.923	0.005\\
0.949	0.006\\
0.974	0.005\\
1.000	0.005\\
};
\addplot[const plot, color=white!80!black, dashdotted, line width=1.2pt, forget plot] table[row sep=crcr] {%
0.000	0.518\\
0.026	0.077\\
0.051	0.049\\
0.077	0.033\\
0.103	0.025\\
0.128	0.024\\
0.154	0.018\\
0.179	0.016\\
0.205	0.015\\
0.231	0.012\\
0.256	0.013\\
0.282	0.012\\
0.308	0.011\\
0.333	0.012\\
0.359	0.008\\
0.385	0.010\\
0.410	0.008\\
0.436	0.008\\
0.462	0.008\\
0.487	0.007\\
0.513	0.009\\
0.538	0.008\\
0.564	0.007\\
0.590	0.007\\
0.615	0.007\\
0.641	0.006\\
0.667	0.007\\
0.692	0.006\\
0.718	0.005\\
0.744	0.005\\
0.769	0.005\\
0.795	0.005\\
0.821	0.006\\
0.846	0.005\\
0.872	0.005\\
0.897	0.006\\
0.923	0.005\\
0.949	0.006\\
0.974	0.005\\
1.000	0.005\\
};
\addplot[const plot, color=white!90!black, line width=1.2pt, forget plot] table[row sep=crcr] {%
0.000	0.524\\
0.026	0.076\\
0.051	0.044\\
0.077	0.033\\
0.103	0.025\\
0.128	0.022\\
0.154	0.019\\
0.179	0.016\\
0.205	0.013\\
0.231	0.014\\
0.256	0.010\\
0.282	0.012\\
0.308	0.010\\
0.333	0.011\\
0.359	0.010\\
0.385	0.009\\
0.410	0.007\\
0.436	0.009\\
0.462	0.009\\
0.487	0.008\\
0.513	0.008\\
0.538	0.007\\
0.564	0.006\\
0.590	0.008\\
0.615	0.006\\
0.641	0.006\\
0.667	0.006\\
0.692	0.007\\
0.718	0.007\\
0.744	0.006\\
0.769	0.005\\
0.795	0.006\\
0.821	0.007\\
0.846	0.006\\
0.872	0.006\\
0.897	0.006\\
0.923	0.005\\
0.949	0.005\\
0.974	0.005\\
1.000	0.005\\
};
\addplot [color=black, line width=1.2pt, forget plot]
  table[row sep=crcr]{%
0.050	0.000\\
0.050	1.000\\
};
\end{axis}

\begin{axis}[%
width=0.411\figurewidth,
height=0.264\figureheight,
at={(0\figurewidth,0\figureheight)},
scale only axis,
xmin=0.000,
xmax=1.000,
xlabel style={font=\color{mycolor1}},
xlabel={$\sym{p-Wert}$},
ymin=0.000,
ymax=1.000,
ylabel style={font=\color{mycolor1}},
ylabel={Häufigkeit},
axis background/.style={fill=white},
title style={font=\bfseries\color{mycolor1}},
title={\textbf{$\sym{beta}_{11}$}},
xmajorgrids,
ymajorgrids,
grid style={dotted, opacity=0.25}
]
\addplot[const plot, color=white!10!black, line width=1.2pt, forget plot] table[row sep=crcr] {%
0.000	0.673\\
0.026	0.054\\
0.051	0.033\\
0.077	0.023\\
0.103	0.017\\
0.128	0.014\\
0.154	0.012\\
0.179	0.011\\
0.205	0.011\\
0.231	0.009\\
0.256	0.008\\
0.282	0.007\\
0.308	0.008\\
0.333	0.006\\
0.359	0.006\\
0.385	0.006\\
0.410	0.006\\
0.436	0.005\\
0.462	0.006\\
0.487	0.005\\
0.513	0.004\\
0.538	0.005\\
0.564	0.004\\
0.590	0.004\\
0.615	0.006\\
0.641	0.004\\
0.667	0.004\\
0.692	0.004\\
0.718	0.004\\
0.744	0.004\\
0.769	0.004\\
0.795	0.004\\
0.821	0.004\\
0.846	0.004\\
0.872	0.005\\
0.897	0.003\\
0.923	0.005\\
0.949	0.004\\
0.974	0.004\\
1.000	0.004\\
};
\addplot[const plot, color=white!20!black, dashed, line width=1.2pt, forget plot] table[row sep=crcr] {%
0.000	0.620\\
0.026	0.062\\
0.051	0.038\\
0.077	0.029\\
0.103	0.018\\
0.128	0.015\\
0.154	0.017\\
0.179	0.012\\
0.205	0.011\\
0.231	0.011\\
0.256	0.010\\
0.282	0.008\\
0.308	0.007\\
0.333	0.009\\
0.359	0.008\\
0.385	0.005\\
0.410	0.006\\
0.436	0.006\\
0.462	0.007\\
0.487	0.006\\
0.513	0.006\\
0.538	0.005\\
0.564	0.005\\
0.590	0.004\\
0.615	0.006\\
0.641	0.005\\
0.667	0.004\\
0.692	0.006\\
0.718	0.005\\
0.744	0.004\\
0.769	0.005\\
0.795	0.005\\
0.821	0.004\\
0.846	0.004\\
0.872	0.004\\
0.897	0.005\\
0.923	0.004\\
0.949	0.005\\
0.974	0.005\\
1.000	0.005\\
};
\addplot[const plot, color=white!30!black, dotted, line width=1.2pt, forget plot] table[row sep=crcr] {%
0.000	0.726\\
0.026	0.046\\
0.051	0.027\\
0.077	0.019\\
0.103	0.014\\
0.128	0.012\\
0.154	0.010\\
0.179	0.009\\
0.205	0.009\\
0.231	0.008\\
0.256	0.008\\
0.282	0.006\\
0.308	0.006\\
0.333	0.006\\
0.359	0.005\\
0.385	0.006\\
0.410	0.005\\
0.436	0.005\\
0.462	0.004\\
0.487	0.004\\
0.513	0.004\\
0.538	0.004\\
0.564	0.004\\
0.590	0.004\\
0.615	0.003\\
0.641	0.004\\
0.667	0.004\\
0.692	0.003\\
0.718	0.004\\
0.744	0.004\\
0.769	0.004\\
0.795	0.004\\
0.821	0.004\\
0.846	0.004\\
0.872	0.003\\
0.897	0.003\\
0.923	0.003\\
0.949	0.004\\
0.974	0.003\\
1.000	0.003\\
};
\addplot[const plot, color=white!40!black, dashdotted, line width=1.2pt, forget plot] table[row sep=crcr] {%
0.000	0.614\\
0.026	0.065\\
0.051	0.038\\
0.077	0.025\\
0.103	0.021\\
0.128	0.016\\
0.154	0.016\\
0.179	0.011\\
0.205	0.010\\
0.231	0.011\\
0.256	0.011\\
0.282	0.009\\
0.308	0.007\\
0.333	0.008\\
0.359	0.007\\
0.385	0.007\\
0.410	0.006\\
0.436	0.008\\
0.462	0.007\\
0.487	0.007\\
0.513	0.005\\
0.538	0.006\\
0.564	0.005\\
0.590	0.005\\
0.615	0.005\\
0.641	0.006\\
0.667	0.006\\
0.692	0.005\\
0.718	0.004\\
0.744	0.006\\
0.769	0.006\\
0.795	0.004\\
0.821	0.004\\
0.846	0.004\\
0.872	0.005\\
0.897	0.005\\
0.923	0.004\\
0.949	0.003\\
0.974	0.005\\
1.000	0.005\\
};
\addplot[const plot, color=gray, line width=1.2pt, forget plot] table[row sep=crcr] {%
0.000	0.670\\
0.026	0.055\\
0.051	0.032\\
0.077	0.024\\
0.103	0.018\\
0.128	0.015\\
0.154	0.011\\
0.179	0.012\\
0.205	0.010\\
0.231	0.008\\
0.256	0.010\\
0.282	0.008\\
0.308	0.007\\
0.333	0.006\\
0.359	0.006\\
0.385	0.005\\
0.410	0.007\\
0.436	0.005\\
0.462	0.005\\
0.487	0.005\\
0.513	0.005\\
0.538	0.006\\
0.564	0.005\\
0.590	0.005\\
0.615	0.004\\
0.641	0.005\\
0.667	0.005\\
0.692	0.004\\
0.718	0.004\\
0.744	0.004\\
0.769	0.004\\
0.795	0.003\\
0.821	0.004\\
0.846	0.004\\
0.872	0.005\\
0.897	0.004\\
0.923	0.005\\
0.949	0.004\\
0.974	0.003\\
1.000	0.003\\
};
\addplot[const plot, color=white!60!black, dashed, line width=1.2pt, forget plot] table[row sep=crcr] {%
0.000	0.645\\
0.026	0.058\\
0.051	0.037\\
0.077	0.023\\
0.103	0.018\\
0.128	0.015\\
0.154	0.015\\
0.179	0.013\\
0.205	0.009\\
0.231	0.010\\
0.256	0.010\\
0.282	0.009\\
0.308	0.007\\
0.333	0.007\\
0.359	0.006\\
0.385	0.006\\
0.410	0.007\\
0.436	0.006\\
0.462	0.007\\
0.487	0.005\\
0.513	0.005\\
0.538	0.004\\
0.564	0.005\\
0.590	0.004\\
0.615	0.006\\
0.641	0.005\\
0.667	0.004\\
0.692	0.005\\
0.718	0.004\\
0.744	0.005\\
0.769	0.004\\
0.795	0.004\\
0.821	0.005\\
0.846	0.004\\
0.872	0.005\\
0.897	0.004\\
0.923	0.005\\
0.949	0.005\\
0.974	0.005\\
1.000	0.005\\
};
\addplot[const plot, color=white!70!black, dotted, line width=1.2pt, forget plot] table[row sep=crcr] {%
0.000	0.667\\
0.026	0.053\\
0.051	0.034\\
0.077	0.022\\
0.103	0.017\\
0.128	0.014\\
0.154	0.015\\
0.179	0.010\\
0.205	0.010\\
0.231	0.009\\
0.256	0.008\\
0.282	0.009\\
0.308	0.007\\
0.333	0.006\\
0.359	0.007\\
0.385	0.005\\
0.410	0.006\\
0.436	0.006\\
0.462	0.005\\
0.487	0.006\\
0.513	0.005\\
0.538	0.005\\
0.564	0.005\\
0.590	0.004\\
0.615	0.005\\
0.641	0.003\\
0.667	0.004\\
0.692	0.004\\
0.718	0.004\\
0.744	0.004\\
0.769	0.003\\
0.795	0.005\\
0.821	0.004\\
0.846	0.003\\
0.872	0.005\\
0.897	0.006\\
0.923	0.005\\
0.949	0.005\\
0.974	0.003\\
1.000	0.003\\
};
\addplot[const plot, color=white!80!black, dashdotted, line width=1.2pt, forget plot] table[row sep=crcr] {%
0.000	0.553\\
0.026	0.071\\
0.051	0.045\\
0.077	0.033\\
0.103	0.023\\
0.128	0.021\\
0.154	0.016\\
0.179	0.014\\
0.205	0.015\\
0.231	0.012\\
0.256	0.011\\
0.282	0.012\\
0.308	0.012\\
0.333	0.010\\
0.359	0.009\\
0.385	0.009\\
0.410	0.008\\
0.436	0.009\\
0.462	0.007\\
0.487	0.006\\
0.513	0.007\\
0.538	0.007\\
0.564	0.005\\
0.590	0.008\\
0.615	0.006\\
0.641	0.007\\
0.667	0.007\\
0.692	0.005\\
0.718	0.005\\
0.744	0.004\\
0.769	0.006\\
0.795	0.005\\
0.821	0.005\\
0.846	0.005\\
0.872	0.003\\
0.897	0.006\\
0.923	0.004\\
0.949	0.004\\
0.974	0.004\\
1.000	0.004\\
};
\addplot[const plot, color=white!90!black, line width=1.2pt, forget plot] table[row sep=crcr] {%
0.000	0.667\\
0.026	0.056\\
0.051	0.032\\
0.077	0.024\\
0.103	0.018\\
0.128	0.013\\
0.154	0.014\\
0.179	0.012\\
0.205	0.011\\
0.231	0.009\\
0.256	0.009\\
0.282	0.008\\
0.308	0.007\\
0.333	0.006\\
0.359	0.008\\
0.385	0.006\\
0.410	0.004\\
0.436	0.005\\
0.462	0.005\\
0.487	0.004\\
0.513	0.007\\
0.538	0.004\\
0.564	0.005\\
0.590	0.004\\
0.615	0.004\\
0.641	0.005\\
0.667	0.004\\
0.692	0.004\\
0.718	0.004\\
0.744	0.004\\
0.769	0.004\\
0.795	0.005\\
0.821	0.004\\
0.846	0.004\\
0.872	0.004\\
0.897	0.004\\
0.923	0.004\\
0.949	0.004\\
0.974	0.004\\
1.000	0.004\\
};
\addplot [color=black, line width=1.2pt, forget plot]
  table[row sep=crcr]{%
0.050	0.000\\
0.050	1.000\\
};
\end{axis}

\begin{axis}[%
width=0.411\figurewidth,
height=0.264\figureheight,
at={(0.541\figurewidth,0\figureheight)},
scale only axis,
xmin=0.000,
xmax=1.000,
xlabel style={font=\color{mycolor1}},
xlabel={$\sym{p-Wert}$},
ymin=0.000,
ymax=1.000,
ytick={0.000,0.200,0.400,0.600,0.800,1.000},
yticklabels={\empty},
axis background/.style={fill=white},
title style={font=\bfseries\color{mycolor1}},
title={\textbf{$\sym{beta}_{22}$}},
xmajorgrids,
ymajorgrids,
grid style={dotted, opacity=0.25}
]
\addplot[const plot, color=white!10!black, line width=1.2pt, forget plot] table[row sep=crcr] {%
0.000	0.668\\
0.026	0.051\\
0.051	0.031\\
0.077	0.024\\
0.103	0.018\\
0.128	0.015\\
0.154	0.013\\
0.179	0.011\\
0.205	0.009\\
0.231	0.010\\
0.256	0.010\\
0.282	0.009\\
0.308	0.008\\
0.333	0.005\\
0.359	0.006\\
0.385	0.006\\
0.410	0.007\\
0.436	0.006\\
0.462	0.004\\
0.487	0.005\\
0.513	0.005\\
0.538	0.006\\
0.564	0.005\\
0.590	0.005\\
0.615	0.005\\
0.641	0.005\\
0.667	0.004\\
0.692	0.004\\
0.718	0.005\\
0.744	0.003\\
0.769	0.003\\
0.795	0.003\\
0.821	0.005\\
0.846	0.005\\
0.872	0.003\\
0.897	0.004\\
0.923	0.004\\
0.949	0.005\\
0.974	0.004\\
1.000	0.004\\
};
\addplot[const plot, color=white!20!black, dashed, line width=1.2pt, forget plot] table[row sep=crcr] {%
0.000	0.672\\
0.026	0.055\\
0.051	0.032\\
0.077	0.022\\
0.103	0.017\\
0.128	0.017\\
0.154	0.012\\
0.179	0.012\\
0.205	0.008\\
0.231	0.010\\
0.256	0.007\\
0.282	0.009\\
0.308	0.008\\
0.333	0.007\\
0.359	0.006\\
0.385	0.006\\
0.410	0.006\\
0.436	0.007\\
0.462	0.005\\
0.487	0.006\\
0.513	0.006\\
0.538	0.004\\
0.564	0.004\\
0.590	0.004\\
0.615	0.004\\
0.641	0.003\\
0.667	0.005\\
0.692	0.004\\
0.718	0.004\\
0.744	0.005\\
0.769	0.002\\
0.795	0.003\\
0.821	0.004\\
0.846	0.005\\
0.872	0.004\\
0.897	0.003\\
0.923	0.003\\
0.949	0.004\\
0.974	0.004\\
1.000	0.004\\
};
\addplot[const plot, color=white!30!black, dotted, line width=1.2pt, forget plot] table[row sep=crcr] {%
0.000	0.677\\
0.026	0.056\\
0.051	0.031\\
0.077	0.023\\
0.103	0.018\\
0.128	0.014\\
0.154	0.010\\
0.179	0.012\\
0.205	0.008\\
0.231	0.009\\
0.256	0.007\\
0.282	0.007\\
0.308	0.007\\
0.333	0.006\\
0.359	0.007\\
0.385	0.006\\
0.410	0.006\\
0.436	0.004\\
0.462	0.005\\
0.487	0.006\\
0.513	0.005\\
0.538	0.004\\
0.564	0.004\\
0.590	0.004\\
0.615	0.004\\
0.641	0.005\\
0.667	0.005\\
0.692	0.003\\
0.718	0.003\\
0.744	0.004\\
0.769	0.005\\
0.795	0.004\\
0.821	0.004\\
0.846	0.004\\
0.872	0.004\\
0.897	0.004\\
0.923	0.004\\
0.949	0.004\\
0.974	0.004\\
1.000	0.004\\
};
\addplot[const plot, color=white!40!black, dashdotted, line width=1.2pt, forget plot] table[row sep=crcr] {%
0.000	0.675\\
0.026	0.054\\
0.051	0.031\\
0.077	0.021\\
0.103	0.017\\
0.128	0.014\\
0.154	0.012\\
0.179	0.011\\
0.205	0.010\\
0.231	0.009\\
0.256	0.007\\
0.282	0.007\\
0.308	0.006\\
0.333	0.006\\
0.359	0.006\\
0.385	0.005\\
0.410	0.005\\
0.436	0.006\\
0.462	0.006\\
0.487	0.005\\
0.513	0.006\\
0.538	0.005\\
0.564	0.005\\
0.590	0.005\\
0.615	0.004\\
0.641	0.005\\
0.667	0.005\\
0.692	0.005\\
0.718	0.004\\
0.744	0.005\\
0.769	0.004\\
0.795	0.006\\
0.821	0.004\\
0.846	0.004\\
0.872	0.004\\
0.897	0.004\\
0.923	0.004\\
0.949	0.005\\
0.974	0.003\\
1.000	0.003\\
};
\addplot[const plot, color=gray, line width=1.2pt, forget plot] table[row sep=crcr] {%
0.000	0.672\\
0.026	0.054\\
0.051	0.029\\
0.077	0.022\\
0.103	0.018\\
0.128	0.016\\
0.154	0.009\\
0.179	0.011\\
0.205	0.008\\
0.231	0.011\\
0.256	0.008\\
0.282	0.008\\
0.308	0.008\\
0.333	0.007\\
0.359	0.007\\
0.385	0.008\\
0.410	0.006\\
0.436	0.006\\
0.462	0.005\\
0.487	0.005\\
0.513	0.005\\
0.538	0.005\\
0.564	0.005\\
0.590	0.003\\
0.615	0.004\\
0.641	0.004\\
0.667	0.004\\
0.692	0.004\\
0.718	0.004\\
0.744	0.004\\
0.769	0.005\\
0.795	0.005\\
0.821	0.005\\
0.846	0.004\\
0.872	0.003\\
0.897	0.005\\
0.923	0.004\\
0.949	0.004\\
0.974	0.004\\
1.000	0.004\\
};
\addplot[const plot, color=white!60!black, dashed, line width=1.2pt, forget plot] table[row sep=crcr] {%
0.000	0.653\\
0.026	0.060\\
0.051	0.035\\
0.077	0.024\\
0.103	0.019\\
0.128	0.015\\
0.154	0.013\\
0.179	0.012\\
0.205	0.010\\
0.231	0.011\\
0.256	0.008\\
0.282	0.007\\
0.308	0.007\\
0.333	0.006\\
0.359	0.006\\
0.385	0.005\\
0.410	0.007\\
0.436	0.006\\
0.462	0.006\\
0.487	0.004\\
0.513	0.006\\
0.538	0.004\\
0.564	0.004\\
0.590	0.004\\
0.615	0.006\\
0.641	0.006\\
0.667	0.003\\
0.692	0.005\\
0.718	0.005\\
0.744	0.004\\
0.769	0.004\\
0.795	0.004\\
0.821	0.005\\
0.846	0.004\\
0.872	0.004\\
0.897	0.004\\
0.923	0.004\\
0.949	0.004\\
0.974	0.003\\
1.000	0.003\\
};
\addplot[const plot, color=white!70!black, dotted, line width=1.2pt, forget plot] table[row sep=crcr] {%
0.000	0.698\\
0.026	0.052\\
0.051	0.028\\
0.077	0.019\\
0.103	0.018\\
0.128	0.013\\
0.154	0.012\\
0.179	0.008\\
0.205	0.009\\
0.231	0.010\\
0.256	0.007\\
0.282	0.007\\
0.308	0.007\\
0.333	0.006\\
0.359	0.005\\
0.385	0.004\\
0.410	0.006\\
0.436	0.006\\
0.462	0.004\\
0.487	0.005\\
0.513	0.005\\
0.538	0.005\\
0.564	0.005\\
0.590	0.004\\
0.615	0.005\\
0.641	0.004\\
0.667	0.003\\
0.692	0.004\\
0.718	0.004\\
0.744	0.004\\
0.769	0.003\\
0.795	0.004\\
0.821	0.003\\
0.846	0.004\\
0.872	0.005\\
0.897	0.004\\
0.923	0.004\\
0.949	0.004\\
0.974	0.004\\
1.000	0.004\\
};
\addplot[const plot, color=white!80!black, dashdotted, line width=1.2pt, forget plot] table[row sep=crcr] {%
0.000	0.657\\
0.026	0.058\\
0.051	0.036\\
0.077	0.024\\
0.103	0.018\\
0.128	0.016\\
0.154	0.013\\
0.179	0.011\\
0.205	0.011\\
0.231	0.006\\
0.256	0.009\\
0.282	0.007\\
0.308	0.008\\
0.333	0.007\\
0.359	0.006\\
0.385	0.006\\
0.410	0.005\\
0.436	0.006\\
0.462	0.006\\
0.487	0.005\\
0.513	0.004\\
0.538	0.005\\
0.564	0.004\\
0.590	0.005\\
0.615	0.005\\
0.641	0.005\\
0.667	0.005\\
0.692	0.006\\
0.718	0.004\\
0.744	0.005\\
0.769	0.004\\
0.795	0.004\\
0.821	0.004\\
0.846	0.004\\
0.872	0.004\\
0.897	0.004\\
0.923	0.005\\
0.949	0.004\\
0.974	0.004\\
1.000	0.004\\
};
\addplot[const plot, color=white!90!black, line width=1.2pt, forget plot] table[row sep=crcr] {%
0.000	0.661\\
0.026	0.057\\
0.051	0.034\\
0.077	0.023\\
0.103	0.020\\
0.128	0.013\\
0.154	0.013\\
0.179	0.012\\
0.205	0.009\\
0.231	0.009\\
0.256	0.008\\
0.282	0.008\\
0.308	0.007\\
0.333	0.008\\
0.359	0.007\\
0.385	0.006\\
0.410	0.006\\
0.436	0.006\\
0.462	0.004\\
0.487	0.004\\
0.513	0.006\\
0.538	0.004\\
0.564	0.004\\
0.590	0.006\\
0.615	0.005\\
0.641	0.005\\
0.667	0.006\\
0.692	0.004\\
0.718	0.005\\
0.744	0.005\\
0.769	0.003\\
0.795	0.004\\
0.821	0.004\\
0.846	0.003\\
0.872	0.004\\
0.897	0.004\\
0.923	0.004\\
0.949	0.004\\
0.974	0.003\\
1.000	0.003\\
};
\addplot [color=black, line width=1.2pt, forget plot]
  table[row sep=crcr]{%
0.050	0.000\\
0.050	1.000\\
};
\addplot [color=black, line width=1.2pt, forget plot]
  table[row sep=crcr]{%
0.050	0.000\\
0.050	1.000\\
};
\addplot [color=black, line width=1.2pt, forget plot]
  table[row sep=crcr]{%
0.050	0.000\\
0.050	1.000\\
};
\addplot [color=black, line width=1.2pt, forget plot]
  table[row sep=crcr]{%
0.050	0.000\\
0.050	1.000\\
};
\addplot [color=black, line width=1.2pt, forget plot]
  table[row sep=crcr]{%
0.050	0.000\\
0.050	1.000\\
};
\addplot [color=black, line width=1.2pt, forget plot]
  table[row sep=crcr]{%
0.050	0.000\\
0.050	1.000\\
};
\end{axis}
\end{tikzpicture}%
    \caption{Relative Häufigkeitender in der Simulation ermittelten $\sym{p-Wert}$e der Signifikanztests.}
    \label{fig:app.c_pval_hist}
\end{figure}

\begin{figure}[htbp]
    \centering
    \setlength\figurewidth{0.5\textwidth}
    \setlength\figureheight{16cm}
    % This file was created by matlab2tikz.
%
\definecolor{mycolor1}{rgb}{0.12941,0.12941,0.12941}%
%
\begin{tikzpicture}

\begin{axis}[%
width=0.411\textwidth,
height=0.264\figureheight,
at={(0\textwidth,0.736\figureheight)},
scale only axis,
xmin=4.4830,
xmax=15.7579,
xtick={5.0000,10.0000,15.0000},
xticklabels={\empty},
ymin=-5.5700,
ymax=11.5200,
ylabel style={font=\color{mycolor1}},
ylabel={$\Delta \sym{power}$ [\%]},
axis background/.style={fill=white},
title style={font=\bfseries\color{mycolor1}},
title={\textbf{$\sym{beta}_0$} ($\sym{corr} = 0{,}75$)},
xmajorgrids,
ymajorgrids,
grid style={dotted, opacity=0.25}
]
\addplot[only marks, mark=*, mark options={}, mark size=1.3693pt, color=black, fill=black, forget plot] table[row sep=crcr]{%
x	y\\
6.9512	0.0000\\
6.8667	0.0100\\
6.9004	0.0000\\
6.8775	0.0100\\
6.8838	0.0000\\
6.9736	0.0000\\
6.8352	-0.0200\\
7.1613	0.0300\\
};
\addplot [color=black, dashed, line width=1.2pt, forget plot]
  table[row sep=crcr]{%
6.8352	-0.0061\\
6.8385	-0.0058\\
6.8418	-0.0054\\
6.8451	-0.0051\\
6.8483	-0.0047\\
6.8516	-0.0044\\
6.8549	-0.0041\\
6.8582	-0.0037\\
6.8615	-0.0034\\
6.8648	-0.0031\\
6.8681	-0.0027\\
6.8714	-0.0024\\
6.8747	-0.0020\\
6.8780	-0.0017\\
6.8813	-0.0014\\
6.8846	-0.0010\\
6.8879	-0.0007\\
6.8912	-0.0004\\
6.8945	-0.0000\\
6.8978	0.0003\\
6.9011	0.0007\\
6.9043	0.0010\\
6.9076	0.0013\\
6.9109	0.0017\\
6.9142	0.0020\\
6.9175	0.0023\\
6.9208	0.0027\\
6.9241	0.0030\\
6.9274	0.0034\\
6.9307	0.0037\\
6.9340	0.0040\\
6.9373	0.0044\\
6.9406	0.0047\\
6.9439	0.0050\\
6.9472	0.0054\\
6.9505	0.0057\\
6.9538	0.0061\\
6.9571	0.0064\\
6.9603	0.0067\\
6.9636	0.0071\\
6.9669	0.0074\\
6.9702	0.0077\\
6.9735	0.0081\\
6.9768	0.0084\\
6.9801	0.0088\\
6.9834	0.0091\\
6.9867	0.0094\\
6.9900	0.0098\\
6.9933	0.0101\\
6.9966	0.0104\\
6.9999	0.0108\\
7.0032	0.0111\\
7.0065	0.0115\\
7.0098	0.0118\\
7.0131	0.0121\\
7.0163	0.0125\\
7.0196	0.0128\\
7.0229	0.0131\\
7.0262	0.0135\\
7.0295	0.0138\\
7.0328	0.0142\\
7.0361	0.0145\\
7.0394	0.0148\\
7.0427	0.0152\\
7.0460	0.0155\\
7.0493	0.0158\\
7.0526	0.0162\\
7.0559	0.0165\\
7.0592	0.0169\\
7.0625	0.0172\\
7.0658	0.0175\\
7.0691	0.0179\\
7.0723	0.0182\\
7.0756	0.0185\\
7.0789	0.0189\\
7.0822	0.0192\\
7.0855	0.0196\\
7.0888	0.0199\\
7.0921	0.0202\\
7.0954	0.0206\\
7.0987	0.0209\\
7.1020	0.0212\\
7.1053	0.0216\\
7.1086	0.0219\\
7.1119	0.0223\\
7.1152	0.0226\\
7.1185	0.0229\\
7.1218	0.0233\\
7.1251	0.0236\\
7.1283	0.0239\\
7.1316	0.0243\\
7.1349	0.0246\\
7.1382	0.0250\\
7.1415	0.0253\\
7.1448	0.0256\\
7.1481	0.0260\\
7.1514	0.0263\\
7.1547	0.0266\\
7.1580	0.0270\\
7.1613	0.0273\\
};
\end{axis}

\begin{axis}[%
width=0.411\textwidth,
height=0.264\figureheight,
at={(0.541\textwidth,0.736\figureheight)},
scale only axis,
xmin=4.4830,
xmax=15.7579,
xtick={5.0000,10.0000,15.0000},
xticklabels={\empty},
ymin=-5.5700,
ymax=11.5200,
ytick={-5.0000,0.0000,5.0000,10.0000},
yticklabels={\empty},
axis background/.style={fill=white},
title style={font=\bfseries\color{mycolor1}},
title={\textbf{$\sym{beta}_1$} ($\sym{corr} = 1{,}00$)},
xmajorgrids,
ymajorgrids,
grid style={dotted, opacity=0.25}
]
\addplot[only marks, mark=*, mark options={}, mark size=1.3693pt, color=black, fill=black, forget plot] table[row sep=crcr]{%
x	y\\
6.9696	0.0000\\
8.0297	2.0400\\
6.6127	-0.8900\\
8.6060	2.8100\\
7.1251	0.0400\\
8.0796	2.3800\\
7.1550	0.3100\\
11.9090	8.4000\\
};
\addplot [color=black, dashed, line width=1.2pt, forget plot]
  table[row sep=crcr]{%
6.6127	-0.6227\\
6.6662	-0.5300\\
6.7197	-0.4373\\
6.7732	-0.3446\\
6.8267	-0.2519\\
6.8802	-0.1592\\
6.9337	-0.0665\\
6.9872	0.0262\\
7.0407	0.1188\\
7.0942	0.2115\\
7.1477	0.3042\\
7.2012	0.3969\\
7.2547	0.4896\\
7.3082	0.5823\\
7.3617	0.6750\\
7.4152	0.7676\\
7.4687	0.8603\\
7.5222	0.9530\\
7.5757	1.0457\\
7.6292	1.1384\\
7.6827	1.2311\\
7.7362	1.3238\\
7.7897	1.4165\\
7.8432	1.5091\\
7.8967	1.6018\\
7.9502	1.6945\\
8.0037	1.7872\\
8.0572	1.8799\\
8.1107	1.9726\\
8.1642	2.0653\\
8.2177	2.1580\\
8.2712	2.2506\\
8.3247	2.3433\\
8.3782	2.4360\\
8.4317	2.5287\\
8.4852	2.6214\\
8.5387	2.7141\\
8.5922	2.8068\\
8.6457	2.8994\\
8.6992	2.9921\\
8.7527	3.0848\\
8.8061	3.1775\\
8.8596	3.2702\\
8.9131	3.3629\\
8.9666	3.4556\\
9.0201	3.5483\\
9.0736	3.6409\\
9.1271	3.7336\\
9.1806	3.8263\\
9.2341	3.9190\\
9.2876	4.0117\\
9.3411	4.1044\\
9.3946	4.1971\\
9.4481	4.2897\\
9.5016	4.3824\\
9.5551	4.4751\\
9.6086	4.5678\\
9.6621	4.6605\\
9.7156	4.7532\\
9.7691	4.8459\\
9.8226	4.9386\\
9.8761	5.0312\\
9.9296	5.1239\\
9.9831	5.2166\\
10.0366	5.3093\\
10.0901	5.4020\\
10.1436	5.4947\\
10.1971	5.5874\\
10.2506	5.6801\\
10.3041	5.7727\\
10.3576	5.8654\\
10.4111	5.9581\\
10.4646	6.0508\\
10.5181	6.1435\\
10.5716	6.2362\\
10.6251	6.3289\\
10.6786	6.4215\\
10.7321	6.5142\\
10.7856	6.6069\\
10.8391	6.6996\\
10.8926	6.7923\\
10.9461	6.8850\\
10.9996	6.9777\\
11.0531	7.0704\\
11.1066	7.1630\\
11.1601	7.2557\\
11.2136	7.3484\\
11.2671	7.4411\\
11.3206	7.5338\\
11.3741	7.6265\\
11.4276	7.7192\\
11.4811	7.8118\\
11.5346	7.9045\\
11.5881	7.9972\\
11.6416	8.0899\\
11.6951	8.1826\\
11.7486	8.2753\\
11.8020	8.3680\\
11.8555	8.4607\\
11.9090	8.5533\\
};
\end{axis}

\begin{axis}[%
width=0.411\textwidth,
height=0.264\figureheight,
at={(0\textwidth,0.368\figureheight)},
scale only axis,
xmin=4.4830,
xmax=15.7579,
xtick={5.0000,10.0000,15.0000},
xticklabels={\empty},
ymin=-5.5700,
ymax=11.5200,
ylabel style={font=\color{mycolor1}},
ylabel={$\Delta \sym{power}$ [\%]},
axis background/.style={fill=white},
title style={font=\bfseries\color{mycolor1}},
title={\textbf{$\sym{beta}_2$} ($\sym{corr} = 0{,}71$)},
xmajorgrids,
ymajorgrids,
grid style={dotted, opacity=0.25}
]
\addplot[only marks, mark=*, mark options={}, mark size=1.3693pt, color=black, fill=black, forget plot] table[row sep=crcr]{%
x	y\\
7.1245	0.0000\\
7.0814	0.0600\\
7.0676	-0.0100\\
6.9684	0.0500\\
6.9063	-0.0100\\
6.7793	-0.7200\\
6.7143	-1.0100\\
6.9055	0.4300\\
};
\addplot [color=black, dashed, line width=1.2pt, forget plot]
  table[row sep=crcr]{%
6.7143	-0.6733\\
6.7184	-0.6638\\
6.7225	-0.6544\\
6.7267	-0.6450\\
6.7308	-0.6355\\
6.7350	-0.6261\\
6.7391	-0.6167\\
6.7433	-0.6072\\
6.7474	-0.5978\\
6.7515	-0.5883\\
6.7557	-0.5789\\
6.7598	-0.5695\\
6.7640	-0.5600\\
6.7681	-0.5506\\
6.7723	-0.5411\\
6.7764	-0.5317\\
6.7805	-0.5223\\
6.7847	-0.5128\\
6.7888	-0.5034\\
6.7930	-0.4939\\
6.7971	-0.4845\\
6.8013	-0.4751\\
6.8054	-0.4656\\
6.8096	-0.4562\\
6.8137	-0.4467\\
6.8178	-0.4373\\
6.8220	-0.4279\\
6.8261	-0.4184\\
6.8303	-0.4090\\
6.8344	-0.3996\\
6.8386	-0.3901\\
6.8427	-0.3807\\
6.8468	-0.3712\\
6.8510	-0.3618\\
6.8551	-0.3524\\
6.8593	-0.3429\\
6.8634	-0.3335\\
6.8676	-0.3240\\
6.8717	-0.3146\\
6.8758	-0.3052\\
6.8800	-0.2957\\
6.8841	-0.2863\\
6.8883	-0.2768\\
6.8924	-0.2674\\
6.8966	-0.2580\\
6.9007	-0.2485\\
6.9049	-0.2391\\
6.9090	-0.2297\\
6.9131	-0.2202\\
6.9173	-0.2108\\
6.9214	-0.2013\\
6.9256	-0.1919\\
6.9297	-0.1825\\
6.9339	-0.1730\\
6.9380	-0.1636\\
6.9421	-0.1541\\
6.9463	-0.1447\\
6.9504	-0.1353\\
6.9546	-0.1258\\
6.9587	-0.1164\\
6.9629	-0.1069\\
6.9670	-0.0975\\
6.9711	-0.0881\\
6.9753	-0.0786\\
6.9794	-0.0692\\
6.9836	-0.0598\\
6.9877	-0.0503\\
6.9919	-0.0409\\
6.9960	-0.0314\\
7.0001	-0.0220\\
7.0043	-0.0126\\
7.0084	-0.0031\\
7.0126	0.0063\\
7.0167	0.0158\\
7.0209	0.0252\\
7.0250	0.0346\\
7.0292	0.0441\\
7.0333	0.0535\\
7.0374	0.0630\\
7.0416	0.0724\\
7.0457	0.0818\\
7.0499	0.0913\\
7.0540	0.1007\\
7.0582	0.1102\\
7.0623	0.1196\\
7.0664	0.1290\\
7.0706	0.1385\\
7.0747	0.1479\\
7.0789	0.1573\\
7.0830	0.1668\\
7.0872	0.1762\\
7.0913	0.1857\\
7.0954	0.1951\\
7.0996	0.2045\\
7.1037	0.2140\\
7.1079	0.2234\\
7.1120	0.2329\\
7.1162	0.2423\\
7.1203	0.2517\\
7.1245	0.2612\\
};
\end{axis}

\begin{axis}[%
width=0.411\textwidth,
height=0.264\figureheight,
at={(0.541\textwidth,0.368\figureheight)},
scale only axis,
xmin=4.4830,
xmax=15.7579,
xtick={5.0000,10.0000,15.0000},
xticklabels={\empty},
ymin=-5.5700,
ymax=11.5200,
ytick={-5.0000,0.0000,5.0000,10.0000},
yticklabels={\empty},
axis background/.style={fill=white},
title style={font=\bfseries\color{mycolor1}},
title={\textbf{$\sym{beta}_{12}$} ($\sym{corr} = 0{,}96$)},
xmajorgrids,
ymajorgrids,
grid style={dotted, opacity=0.25}
]
\addplot[only marks, mark=*, mark options={}, mark size=1.3693pt, color=black, fill=black, forget plot] table[row sep=crcr]{%
x	y\\
13.8898	0.0000\\
13.9579	-0.0900\\
14.3253	0.0600\\
14.1819	-0.0700\\
13.9968	0.0200\\
11.7337	-3.9400\\
13.9479	-0.1600\\
13.8714	0.6200\\
};
\addplot [color=black, dashed, line width=1.2pt, forget plot]
  table[row sep=crcr]{%
11.7337	-3.7871\\
11.7598	-3.7435\\
11.7860	-3.6998\\
11.8122	-3.6562\\
11.8384	-3.6125\\
11.8645	-3.5689\\
11.8907	-3.5253\\
11.9169	-3.4816\\
11.9431	-3.4380\\
11.9693	-3.3943\\
11.9954	-3.3507\\
12.0216	-3.3070\\
12.0478	-3.2634\\
12.0740	-3.2197\\
12.1002	-3.1761\\
12.1263	-3.1324\\
12.1525	-3.0888\\
12.1787	-3.0451\\
12.2049	-3.0015\\
12.2310	-2.9578\\
12.2572	-2.9142\\
12.2834	-2.8705\\
12.3096	-2.8269\\
12.3358	-2.7832\\
12.3619	-2.7396\\
12.3881	-2.6959\\
12.4143	-2.6523\\
12.4405	-2.6086\\
12.4667	-2.5650\\
12.4928	-2.5213\\
12.5190	-2.4777\\
12.5452	-2.4340\\
12.5714	-2.3904\\
12.5976	-2.3467\\
12.6237	-2.3031\\
12.6499	-2.2594\\
12.6761	-2.2158\\
12.7023	-2.1721\\
12.7284	-2.1285\\
12.7546	-2.0848\\
12.7808	-2.0412\\
12.8070	-1.9975\\
12.8332	-1.9539\\
12.8593	-1.9102\\
12.8855	-1.8666\\
12.9117	-1.8229\\
12.9379	-1.7793\\
12.9641	-1.7356\\
12.9902	-1.6920\\
13.0164	-1.6483\\
13.0426	-1.6047\\
13.0688	-1.5610\\
13.0949	-1.5174\\
13.1211	-1.4737\\
13.1473	-1.4301\\
13.1735	-1.3864\\
13.1997	-1.3428\\
13.2258	-1.2991\\
13.2520	-1.2555\\
13.2782	-1.2118\\
13.3044	-1.1682\\
13.3306	-1.1245\\
13.3567	-1.0809\\
13.3829	-1.0372\\
13.4091	-0.9936\\
13.4353	-0.9499\\
13.4614	-0.9063\\
13.4876	-0.8626\\
13.5138	-0.8190\\
13.5400	-0.7753\\
13.5662	-0.7317\\
13.5923	-0.6880\\
13.6185	-0.6444\\
13.6447	-0.6007\\
13.6709	-0.5571\\
13.6971	-0.5134\\
13.7232	-0.4698\\
13.7494	-0.4261\\
13.7756	-0.3825\\
13.8018	-0.3388\\
13.8279	-0.2952\\
13.8541	-0.2515\\
13.8803	-0.2079\\
13.9065	-0.1642\\
13.9327	-0.1206\\
13.9588	-0.0769\\
13.9850	-0.0333\\
14.0112	0.0104\\
14.0374	0.0540\\
14.0636	0.0977\\
14.0897	0.1413\\
14.1159	0.1850\\
14.1421	0.2286\\
14.1683	0.2723\\
14.1944	0.3159\\
14.2206	0.3596\\
14.2468	0.4032\\
14.2730	0.4469\\
14.2992	0.4905\\
14.3253	0.5342\\
};
\end{axis}

\begin{axis}[%
width=0.411\textwidth,
height=0.264\figureheight,
at={(0\textwidth,0\figureheight)},
scale only axis,
xmin=4.4830,
xmax=15.7579,
xlabel style={font=\color{mycolor1}},
xlabel={$\sym{mse}$},
ymin=-5.5700,
ymax=11.5200,
ylabel style={font=\color{mycolor1}},
ylabel={$\Delta \sym{power}$ [\%]},
axis background/.style={fill=white},
title style={font=\bfseries\color{mycolor1}},
title={\textbf{$\sym{beta}_{11}$} ($\sym{corr} = 1{,}00$)},
xmajorgrids,
ymajorgrids,
grid style={dotted, opacity=0.25}
]
\addplot[only marks, mark=*, mark options={}, mark size=1.3693pt, color=black, fill=black, forget plot] table[row sep=crcr]{%
x	y\\
7.0355	0.0000\\
9.3283	4.3400\\
4.9812	-4.5700\\
10.3692	6.2500\\
7.0474	-0.0500\\
8.6204	2.7400\\
7.3558	1.1900\\
12.2925	10.5200\\
};
\addplot [color=black, dashed, line width=1.2pt, forget plot]
  table[row sep=crcr]{%
4.9812	-4.2389\\
5.0550	-4.0913\\
5.1289	-3.9437\\
5.2027	-3.7961\\
5.2766	-3.6484\\
5.3504	-3.5008\\
5.4243	-3.3532\\
5.4981	-3.2056\\
5.5720	-3.0580\\
5.6458	-2.9103\\
5.7197	-2.7627\\
5.7935	-2.6151\\
5.8674	-2.4675\\
5.9412	-2.3198\\
6.0151	-2.1722\\
6.0889	-2.0246\\
6.1628	-1.8770\\
6.2366	-1.7294\\
6.3105	-1.5817\\
6.3843	-1.4341\\
6.4582	-1.2865\\
6.5320	-1.1389\\
6.6059	-0.9913\\
6.6798	-0.8436\\
6.7536	-0.6960\\
6.8275	-0.5484\\
6.9013	-0.4008\\
6.9752	-0.2531\\
7.0490	-0.1055\\
7.1229	0.0421\\
7.1967	0.1897\\
7.2706	0.3373\\
7.3444	0.4850\\
7.4183	0.6326\\
7.4921	0.7802\\
7.5660	0.9278\\
7.6398	1.0754\\
7.7137	1.2231\\
7.7875	1.3707\\
7.8614	1.5183\\
7.9352	1.6659\\
8.0091	1.8136\\
8.0829	1.9612\\
8.1568	2.1088\\
8.2306	2.2564\\
8.3045	2.4040\\
8.3783	2.5517\\
8.4522	2.6993\\
8.5260	2.8469\\
8.5999	2.9945\\
8.6738	3.1421\\
8.7476	3.2898\\
8.8215	3.4374\\
8.8953	3.5850\\
8.9692	3.7326\\
9.0430	3.8803\\
9.1169	4.0279\\
9.1907	4.1755\\
9.2646	4.3231\\
9.3384	4.4707\\
9.4123	4.6184\\
9.4861	4.7660\\
9.5600	4.9136\\
9.6338	5.0612\\
9.7077	5.2088\\
9.7815	5.3565\\
9.8554	5.5041\\
9.9292	5.6517\\
10.0031	5.7993\\
10.0769	5.9470\\
10.1508	6.0946\\
10.2246	6.2422\\
10.2985	6.3898\\
10.3723	6.5374\\
10.4462	6.6851\\
10.5200	6.8327\\
10.5939	6.9803\\
10.6678	7.1279\\
10.7416	7.2755\\
10.8155	7.4232\\
10.8893	7.5708\\
10.9632	7.7184\\
11.0370	7.8660\\
11.1109	8.0137\\
11.1847	8.1613\\
11.2586	8.3089\\
11.3324	8.4565\\
11.4063	8.6041\\
11.4801	8.7518\\
11.5540	8.8994\\
11.6278	9.0470\\
11.7017	9.1946\\
11.7755	9.3422\\
11.8494	9.4899\\
11.9232	9.6375\\
11.9971	9.7851\\
12.0709	9.9327\\
12.1448	10.0804\\
12.2186	10.2280\\
12.2925	10.3756\\
};
\end{axis}

\begin{axis}[%
width=0.411\textwidth,
height=0.264\figureheight,
at={(0.541\textwidth,0\figureheight)},
scale only axis,
xmin=4.4830,
xmax=15.7579,
xlabel style={font=\color{mycolor1}},
xlabel={$\sym{mse}$},
ymin=-5.5700,
ymax=11.5200,
ytick={-5.0000,0.0000,5.0000,10.0000},
yticklabels={\empty},
axis background/.style={fill=white},
title style={font=\bfseries\color{mycolor1}},
title={\textbf{$\sym{beta}_{22}$} ($\sym{corr} = 0{,}98$)},
xmajorgrids,
ymajorgrids,
grid style={dotted, opacity=0.25}
]
\addplot[only marks, mark=*, mark options={}, mark size=1.3693pt, color=black, fill=black, forget plot] table[row sep=crcr]{%
x	y\\
6.9870	0.0000\\
7.1253	0.0800\\
6.8842	0.0500\\
7.2232	-0.0800\\
7.0186	-0.0700\\
7.7292	1.3500\\
5.8889	-2.9800\\
7.8135	1.7500\\
};
\addplot [color=black, dashed, line width=1.2pt, forget plot]
  table[row sep=crcr]{%
5.8889	-2.7637\\
5.9083	-2.7185\\
5.9278	-2.6734\\
5.9472	-2.6282\\
5.9666	-2.5830\\
5.9861	-2.5379\\
6.0055	-2.4927\\
6.0250	-2.4475\\
6.0444	-2.4024\\
6.0638	-2.3572\\
6.0833	-2.3120\\
6.1027	-2.2669\\
6.1222	-2.2217\\
6.1416	-2.1765\\
6.1610	-2.1313\\
6.1805	-2.0862\\
6.1999	-2.0410\\
6.2194	-1.9958\\
6.2388	-1.9507\\
6.2582	-1.9055\\
6.2777	-1.8603\\
6.2971	-1.8152\\
6.3166	-1.7700\\
6.3360	-1.7248\\
6.3554	-1.6797\\
6.3749	-1.6345\\
6.3943	-1.5893\\
6.4138	-1.5441\\
6.4332	-1.4990\\
6.4526	-1.4538\\
6.4721	-1.4086\\
6.4915	-1.3635\\
6.5110	-1.3183\\
6.5304	-1.2731\\
6.5498	-1.2280\\
6.5693	-1.1828\\
6.5887	-1.1376\\
6.6082	-1.0925\\
6.6276	-1.0473\\
6.6470	-1.0021\\
6.6665	-0.9570\\
6.6859	-0.9118\\
6.7054	-0.8666\\
6.7248	-0.8214\\
6.7442	-0.7763\\
6.7637	-0.7311\\
6.7831	-0.6859\\
6.8026	-0.6408\\
6.8220	-0.5956\\
6.8414	-0.5504\\
6.8609	-0.5053\\
6.8803	-0.4601\\
6.8998	-0.4149\\
6.9192	-0.3698\\
6.9386	-0.3246\\
6.9581	-0.2794\\
6.9775	-0.2343\\
6.9970	-0.1891\\
7.0164	-0.1439\\
7.0358	-0.0987\\
7.0553	-0.0536\\
7.0747	-0.0084\\
7.0942	0.0368\\
7.1136	0.0819\\
7.1330	0.1271\\
7.1525	0.1723\\
7.1719	0.2174\\
7.1914	0.2626\\
7.2108	0.3078\\
7.2303	0.3529\\
7.2497	0.3981\\
7.2691	0.4433\\
7.2886	0.4885\\
7.3080	0.5336\\
7.3275	0.5788\\
7.3469	0.6240\\
7.3663	0.6691\\
7.3858	0.7143\\
7.4052	0.7595\\
7.4247	0.8046\\
7.4441	0.8498\\
7.4635	0.8950\\
7.4830	0.9401\\
7.5024	0.9853\\
7.5219	1.0305\\
7.5413	1.0756\\
7.5607	1.1208\\
7.5802	1.1660\\
7.5996	1.2112\\
7.6191	1.2563\\
7.6385	1.3015\\
7.6579	1.3467\\
7.6774	1.3918\\
7.6968	1.4370\\
7.7163	1.4822\\
7.7357	1.5273\\
7.7551	1.5725\\
7.7746	1.6177\\
7.7940	1.6628\\
7.8135	1.7080\\
};
\end{axis}
\end{tikzpicture}%
    \caption{Zusammenhang zwischen simulierter Schätzvarianz (\ac{MSE}) und dem prozentualen Trennschärfe-Vverlust ($\Delta \sym{power}$) über alle Koeffizienten. Die Pearson-Korrelation $\sym{corr}$ untermauert die signifikante Kausalität.}
    \label{fig:app.c_corr_mse}
\end{figure}